% defer/rcu.tex
% mainfile: ../perfbook.tex
% SPDX-License-Identifier: CC-BY-SA-3.0

\section{Read-Copy Update（RCU）}
\label{sec:defer:Read-Copy Update (RCU)}
%
\epigraph{``免费''是一个\emph{非常}好的价格！}{Tom Peterson}

前面几节讨论的所有机制都使用了多种方法中的一种来推迟特定操作，
直到它们可以安全执行。
\cref{sec:defer:Reference Counting}中讨论的引用计数器
使用显式计数器来延迟可能干扰读者的操作，
这导致了读侧竞争，从而可扩展性差。
\cref{sec:defer:Hazard Pointers}涵盖的 hazard pointers
以每线程指针列表的形式使用隐式计数器。
这避免了读侧竞争，但要求读者执行存储和条件分支，
以及在读侧原语中使用\IXhpl{full}{memory barrier}
或在更新侧原语中使用对实时不友好的\IXacrlpl{ipi}。\footnote{
	在一些重要的特殊情况下，可以通过使用 link counting 来避免
	这些额外工作，如 Folly 开源库中实现的 UnboundedQueue 和
	ConcurrentHashMap 数据结构所示
	（\url{https://github.com/facebook/folly}）。}
\cref{sec:defer:Sequence Locks}中介绍的 sequence lock
也避免了读侧竞争，但不保护指针遍历，
并且像 hazard pointers 一样，需要在读侧原语中使用完整的 memory barriers，
或在更新侧原语中使用处理器间中断。
这些方案的不足引发了一个问题：是否可以做得更好。

本节介绍\IXBacrfst{rcu}，它提供了一个 API，
允许将读者与源代码中的区域关联，
而不是与对频繁更新的共享数据的昂贵更新关联。
本节的其余部分从多个不同角度审视 RCU。
\Cref{sec:defer:Introduction to RCU}提供了经典的 RCU 介绍，
\cref{sec:defer:RCU Fundamentals}涵盖基本的 RCU 概念，
\cref{sec:defer:RCU Linux-Kernel API}介绍 Linux 内核 API，
\cref{sec:defer:RCU Usage}介绍一些常见的 RCU 使用场景，
最后
\cref{sec:defer:RCU Related Work}涵盖与 RCU 相关的近期工作。

虽然 RCU 以微妙和困难著称，但如果正确使用，它是相当直接的。
事实上，Butler Lampson 这样的权威将其分类为简单的
并发~\cite[Chapter 3]{ButlerLampson2022ProgrammingConcurrentSystems}。
Lampson 是完全正确的：
RCU 背后的核心概念不过是保证更新者等待先前存在的
（且\emph{仅仅是}先前存在的）读者完成。
Balmau、Guerraoui、Herlihy 和 Zablotchi 则指出
``QSBR [一种特定类别的 RCU 实现] 速度快且几乎可以应用于
任何数据结构''。

另一方面，要充分利用 RCU 通常需要你以不同的方式思考
一般的并发性和你的特定使用场景。
因此，本节的大部分内容提供了将使用场景映射到 RCU 的指南。

\QuickQuiz{
	RCU 不可能这么简单！！！
	而且你到底能用一个不过是简单延迟的同步原语做什么？？？
}\QuickQuizAnswer{
	抱歉，但 RCU 的核心确实就是这么简单！

	而且你可以用 RCU 的简单延迟做许多令人惊叹的事情。
	请继续阅读本节以了解具体方法。
}\QuickQuizEnd

% defer/rcuintro.tex
% mainfile: ../perfbook.tex
% SPDX-License-Identifier: CC-BY-SA-3.0

\subsection{RCU 简介}
\label{sec:defer:Introduction to RCU}

前面几节讨论的方法为 Pre-BSD 路由表提供了良好的可扩展性，
但性能明显不够理想。
因此，本着``只有走得太远的人才知道能走多远''的精神，\footnote{
	向 T.~S.~Eliot 致歉。}
我们将一探到底，研究并发读者可能执行的汇编语言指令序列
与单线程查找完全相同的算法，尽管存在并发更新。
当然，这个崇高的目标可能会引发严重的可实现性问题，
但如果我们连尝试都不做，就不可能成功！

如果我们成功了，我们将揭开
\cpageref{sec:defer:Mysteries RCU}中提出的又一个秘密。

\subsubsection{最小插入和删除}
\label{sec:defer:Minimal Insertion and Deletion}

\begin{figure}
\centering
\resizebox{3in}{!}{\includegraphics{defer/RCUListInsertClassic}}
\caption{并发读者下的插入}
\label{fig:defer:Insertion With Concurrent Readers}
\end{figure}

为了最小化可实现性问题，我们关注一个最小数据结构，
它由一个全局指针组成，该指针要么是 \co{NULL}，
要么引用一个结构。
尽管它可能很简单，但此数据结构在
生产中被大量使用~\cite{GeoffRomer2018C++DeferredReclamationP0561R4}。
\cref{fig:defer:Insertion With Concurrent Readers}展示了
经典的插入方法，显示了四个状态，时间从上到下推进。
第一行显示初始状态，\co{gptr} 等于 \co{NULL}。
第二行中，我们分配了一个未初始化的结构，如问号所示。
第三行中，我们初始化了该结构。
最后，在第四行也是最后一行，我们更新了 \co{gptr}
以引用新分配和初始化的元素。

我们可能希望对 \co{gptr} 的赋值可以使用简单的 C 语言赋值语句。
不幸的是，\cref{sec:toolsoftrade:Shared-Variable Shenanigans}
粉碎了这些希望。
因此，更新者不能使用简单的 C 语言赋值，
而必须使用如图所示的 \co{smp_store_release()}，
或者如后面将看到的 \co{rcu_assign_pointer()}。

类似地，人们可能希望读者可以使用单个 C 语言赋值来获取 \co{gptr} 的值，
并保证要么得到旧值 \co{NULL}，要么得到新安装的指针，
但无论哪种方式都看到有效的结果。
不幸的是，\cref{sec:toolsoftrade:Shared-Variable Shenanigans}
同样粉碎了这些希望。
为了获得此保证，读者必须使用 \co{READ_ONCE()}，
或者如后面将看到的 \co{rcu_dereference()}。
然而，在大多数现代计算机系统上，这些读侧原语中的每一个
都可以用单条加载指令实现，正是在单线程代码中通常使用的指令。

从读者的角度回顾\cref{fig:defer:Insertion With Concurrent Readers}，
在前三个状态中，所有读者看到 \co{gptr} 的值为 \co{NULL}。
进入第四个状态后，一些读者可能仍然看到 \co{gptr} 的值为 \co{NULL}，
而其他读者可能看到它引用新插入的元素，
但过了一段时间后，所有读者都将看到这个新元素。
在所有时刻，所有读者都将看到 \co{gptr} 包含一个有效指针。
因此，在允许并发读者执行通常用于单线程代码的相同机器指令序列的同时，
确实可以向链式数据结构添加新数据。
这种零成本的并发读取方法提供了出色的性能和可扩展性，
也非常适合实时使用。

\begin{figure}
\centering
\resizebox{3in}{!}{\includegraphics{defer/RCUListDeleteClassic}}
\caption{并发读者下的删除}
\label{fig:defer:Deletion With Concurrent Readers}
\end{figure}

插入当然非常有用，但迟早也需要删除数据。
如\cref{fig:defer:Deletion With Concurrent Readers}所示，
第一步很简单。
再次牢记\cref{sec:toolsoftrade:Shared-Variable Shenanigans}的教训，
使用 \co{smp_store_release()} 将指针置为 \co{NULL}，
从而从图中的第一行移到第二行。
此时，先前存在的读者看到旧结构，其 \co{->addr} 为~42，
\co{->iface} 为~1，但新读者将看到 \co{NULL} 指针，
即并发读者可能对状态不一致，
如图中的``2~Versions''所示。

\QuickQuizSeries{%
\QuickQuizB{
	为什么\cref{fig:defer:Deletion With Concurrent Readers}
	在存储 \co{NULL} 指针时使用 \co{smp_store_release()}？
	鉴于没有结构初始化需要与 \co{NULL} 指针的存储排序，
	\co{WRITE_ONCE()} 在这种情况下不是同样适用吗？
}\QuickQuizAnswerB{
	是的，同样适用。

	因为赋值的是 \co{NULL} 指针，没有什么需要排序的，
	所以不需要 \co{smp_store_release()}。
	相比之下，当赋值非 \co{NULL} 指针时，
	有必要使用 \co{smp_store_release()} 以确保
	指向结构的初始化在指针赋值之前完成。

	简而言之，\co{WRITE_ONCE()} 是可行的，
	并且在某些架构上可以节省少量 CPU 时间。
	然而，正如我们将看到的，软件工程的考虑
	将促使使用与 \co{smp_store_release()} 非常相似的
	特殊 \co{rcu_assign_pointer()}。
}\QuickQuizEndB
%
\QuickQuizE{
	与彼此以及与\cref{fig:defer:Deletion With Concurrent Readers}
	中概述的过程并发运行的读者
	可能对 \co{gptr} 的值不一致。
	这难道不是有点问题吗？？？
}\QuickQuizAnswerE{
	不一定。

	正如\cref{sec:cpu:Hardware Optimizations,sec:cpu:Hardware Free Lunch?}
	中暗示的那样，光速延迟意味着计算机的数据
	与该数据旨在建模的外部现实相比总是过时的。

	因此，现实世界的算法绝对必须容忍外部现实
	与反映该现实的计算机内数据之间的不一致。
	许多此类算法还能够容忍计算机内数据的
	一定程度的不一致。
	\Cref{sec:datastruct:RCU-Protected Hash Table Discussion}
	更详细地讨论了这一点。

	请注意，容忍不一致和过时数据的需要不仅限于 RCU\@。
	它也适用于 reference counting、hazard pointers、sequence
	locks，甚至某些锁的使用场景。
	例如，如果你在持有锁时计算某个值，
	但在释放锁后使用该值，
	你很可能在使用过时的数据。
	毕竟，该值所基于的数据可能在锁释放后立即发生任意变化。

	所以是的，RCU 读者可能看到过时和不一致的数据，
	但不，这不一定是个问题。
	而且在需要时，有 RCU 使用模式可以同时避免
	过时和不一致~\cite{Arcangeli03}。
}\QuickQuizEndE
}

我们只需等待所有先前存在的读者完成，
就可以回到单一版本，如第~3 行所示。\footnote{
	这意味着\cref{sec:defer:Read-Copy Update (RCU)}中引入的
	等待读者属性毕竟确实有重要用途！}
此时，所有先前存在的读者都已完成，
没有后续读者有通往旧数据项的路径，
因此不再有任何读者引用它。
因此可以安全地释放它，如第~4 行所示。

因此，给定一种等待先前存在读者完成的方法，
就可以在链式数据结构中添加和删除数据，
尽管读者执行的机器指令序列与单线程执行相同。
所以也许走到底并没有走得太远！

但我们如何知道所有先前存在的读者实际上何时完成？
这个问题是\cref{sec:defer:Waiting for Readers}的主题。
不过首先，下一节定义 RCU 的核心 API。

\subsubsection{核心 RCU API}
\label{sec:defer:Core RCU API}

完整的 Linux 内核 API 相当广泛，有一百多个 API 成员。
然而，本节将仅限于六个核心 RCU API 成员，
这对于接下来介绍 RCU 和涵盖其基础知识的各节来说已经足够。
完整的 API 在\cref{sec:defer:RCU Linux-Kernel API}中介绍。

核心 API 的三个成员由读者使用。
\apik{rcu_read_lock()} 和 \apik{rcu_read_unlock()} 函数界定
\IXhmrpl{RCU read-side}{critical section}。
它们可以嵌套，使得一个
\co{rcu_read_lock()}--\co{rcu_read_unlock()} 对可以包含在另一个中。
在这种情况下，嵌套的 RCU 读侧临界区集合作为一个大型临界区，
覆盖嵌套集合的完整范围。
第三个读侧 API 成员 \apik{rcu_dereference()} 获取一个
\IXr{RCU-protected pointer}。
概念上，\co{rcu_dereference()} 只是从内存加载，但我们将在
\cref{sec:defer:Publish-Subscribe Mechanism}中看到，
\co{rcu_dereference()} 必须防止编译器和（在一种情况下）CPU
将其加载与后续解引用该指针的内存操作重排序。

\QuickQuiz{
	什么是 RCU-protected pointer？
}\QuickQuizAnswer{
	指向\IXr{RCU-protected data}的指针。
	RCU-protected data 则是一块动态分配的内存，
	其释放将被延迟，使得从不再有任何
	RCU 读者可访问的指向该块的指针到该块被释放之间
	会经过一个 RCU grace period。
	这确保在释放该块时没有 RCU 读者能够访问它。

	RCU-protected pointers 必须小心处理。
	例如，任何打算解引用 RCU-protected pointer 的读者
	必须使用 \co{rcu_dereference()}（或更强的）来加载该指针。
	此外，任何更新者必须使用 \co{rcu_assign_pointer()}
	（或更强的）来存储该指针。
}\QuickQuizEnd

核心 API 的另外三个成员由更新者使用。
\apik{synchronize_rcu()} 函数实现了
\cref{fig:defer:Deletion With Concurrent Readers}中的
``等待读者''操作。
\apik{call_rcu()} 函数是 \co{synchronize_rcu()} 的异步对应物，
它在所有先前存在的 RCU 读者完成后调用指定的函数。
最后，\apik{rcu_assign_pointer()} 宏用于更新 RCU-protected pointer。
概念上，这只是一条赋值语句，但我们将在
\cref{sec:defer:Publish-Subscribe Mechanism}中看到，
\co{rcu_assign_pointer()} 必须防止编译器和 CPU
将此赋值重排序到用于初始化所指向结构的任何先前赋值之前。

\begin{table*}
\renewcommand*{\arraystretch}{1.25}
\rowcolors{2}{}{lightgray}
\small
\centering
\ebresizewidth{
\begin{tabular}{llp{2.4in}}
\toprule
&
	原语 &
		用途 \\
\midrule
\cellcolor{white}\emph{读者} &
	\tco{rcu_read_lock()} &
		开始一个 RCU 读侧临界区。 \\
&
	\tco{rcu_read_unlock()} &
		结束一个 RCU 读侧临界区。 \\
\cellcolor{white}&
	\tco{rcu_dereference()} &
		安全地加载一个 RCU-protected pointer。 \\
\midrule
\emph{更新者} &
	\tco{synchronize_rcu()} &
		等待所有先前存在的 RCU 读侧临界区完成。 \\
\cellcolor{white}&
	\tco{call_rcu()} &
		在所有先前存在的 RCU 读侧临界区完成后调用指定函数。 \\
&
	\tco{rcu_assign_pointer()} &
		安全地更新一个 RCU-protected pointer。 \\
\bottomrule
\end{tabular}
}
\caption{核心 RCU API}
\label{tab:defer:Core RCU API}
\end{table*}

\QuickQuiz{
	如果 \co{synchronize_rcu()} 大约与
	\co{rcu_read_lock()} 同时开始，会怎样？
}\QuickQuizAnswer{
	如果 \co{synchronize_rcu()} 无法证明它在给定的
	\co{rcu_read_lock()} 之前开始，
	则它必须等待对应的 \co{rcu_read_unlock()}。
}\QuickQuizEnd

核心 RCU API 总结在\cref{tab:defer:Core RCU API}中以便参考。
有了这些，我们准备好继续介绍 RCU 的关键操作——等待读者。

\subsubsection{等待读者}
\label{sec:defer:Waiting for Readers}

很容易想到将 \co{synchronize_rcu()} 和 \co{call_rcu()}
的等待读者功能建立在由 \co{rcu_read_lock()} 和
\co{rcu_read_unlock()} 更新的引用计数器之上，
但\cref{chp:Counting}中的
\cref{fig:count:Atomic Increment Scalability on x86}
表明并发引用计数导致极高的开销。
这种极高开销在引用计数器的特定情况下由
\cpageref{fig:defer:Pre-BSD Routing Table Protected by Reference Counting}上的
\cref{fig:defer:Pre-BSD Routing Table Protected by Reference Counting}
所证实。
Hazard pointers 大幅降低了此开销，但正如我们在
\cpageref{fig:defer:Pre-BSD Routing Table Protected by Hazard Pointers}上的
\cref{fig:defer:Pre-BSD Routing Table Protected by Hazard Pointers}
中看到的，没有降到零。
然而，许多 RCU 实现使用精心控制缓存局部性的计数器。

第二种方法观察到内存同步是昂贵的，
因此改用寄存器，即每个 CPU 或线程的程序计数器（PC），
从而至少在没有并发更新的情况下不给读者带来开销。
更新者轮询每个相关的 PC，如果该 PC 不在读侧代码内，
则对应的 CPU 或线程处于 quiescent state，
进而表示可能访问新移除数据元素的任何读者已完成。
一旦所有 CPU 或线程的 PC 都被观察到在任何读者之外，
\IX{grace period}就完成了。
请注意，此方法带来一些严重的挑战，
包括内存排序、\emph{有时}从读者调用的函数，
以及令人兴奋的代码移动优化。
然而，据说此方法在生产中使用~\cite{MikeAsh2015Apple}。

第三种方法是简单地等待一段固定的时间，
该时间足以舒适地超过任何合理读者的
生命周期~\cite{Jacobson93,AjuJohn95}。
这在硬实时系统中可以很好地工作~\cite{YuxinRen2018RTRCU}，
但在不那么特殊的环境中，Murphy 说
即使对不合理的长寿命读者也要做好准备是至关重要的。
要理解这一点，考虑未能做到的后果：
在不合理的读者仍在引用它时，数据项将被释放，
该项很可能被立即重新分配，甚至可能作为其他类型的数据项。
不合理的读者和不知情的重新分配者将试图
将同一块内存用于两个非常不同的目的。
随之而来的混乱将极难调试。

第四种方法是永远等待，确信这样做可以容纳
即使是最不合理的读者。
此方法也称为``内存泄漏''，由于内存泄漏通常需要
不合时宜和不方便的重启，因而名声不好。
然而，当更新速率和运行时间都有严格上限时，这是一个可行的策略。
例如，此方法在高可用性集群中效果很好，
其中系统定期被崩溃以确保集群确实保持高可用性。\footnote{
	强制定期崩溃的程序有时被称为``混沌猴''：
	\url{https://netflix.github.io/chaosmonkey/}。
	然而，忽视系统运行太久导致的混沌也可能是个错误。}
在有垃圾回收器的环境中，泄漏内存也是可行的策略，
在这种情况下，垃圾回收器可以被认为是堵住了泄漏~\cite{Kung80}。
但如果你的环境缺少垃圾回收器，请继续阅读！

第五种方法通过定期``停止世界''来避免定期崩溃，
传统的 stop-the-world 垃圾回收器就是这种方法的范例。
在无处不在的连接之前的几十年中，这种方法也被大量使用，
当时在每个工作日结束时关闭系统电源是常见做法。
然而，在当今始终连接、始终在线的世界中，
停止世界会严重降低响应时间，
这也是开发并发垃圾回收器的一个动机~\cite{DavidFBacon2003RTGC}。
此外，虽然我们需要所有先前存在的读者完成，
但我们不需要它们同时完成。

这一观察引出了第六种方法，即一次停止一个 CPU 或线程。
此方法的优点是完全不降低读者的响应时间，更不用说严重降低了。
此外，许多应用程序已经有只有在所有先前存在的读者完成后
才能到达的状态（称为\emph{\IXpl{quiescent state}}）。
在事务处理系统中，一对连续事务之间的时间可能是 quiescent state。
在反应式系统中，一对连续事件之间的状态可能是 quiescent state。
在非抢占式操作系统内核中，上下文切换可以是
quiescent state~\cite{McKenney98}。
无论哪种方式，一旦所有 CPU 和/或线程都通过了 quiescent state，
系统就被认为完成了一个\emph{grace period}，
此时保证在该 grace period 开始时存在的所有读者都已完成。
因此，也保证可以安全释放在该 grace period 开始之前
移除的任何已移除数据项。\footnote{
	使用 RCU 可以做的远不止简单地推迟内存回收，
	但延迟回收是 RCU 最常见的使用场景，
	因此是一个很好的起点。
	有关延迟执行的更一般情况的示例，
	请参见\cref{sec:defer:Phased State Change}中的
	分阶段状态变更。}

在非抢占式操作系统内核中，要使上下文切换成为有效的
quiescent state，读者必须被禁止在引用通过
\cref{fig:defer:Insertion With Concurrent Readers,%
fig:defer:Deletion With Concurrent Readers}
中所示的 \co{gptr} 指针获得的数据结构的给定实例时阻塞。
此不阻塞约束与纯自旋锁的类似约束一致，
即 CPU 在持有自旋锁时被禁止阻塞。
没有此约束，所有 CPU 可能会被试图获取
被阻塞线程持有的自旋锁的自旋线程消耗。
自旋线程在获取锁之前不会释放其 CPU，
但持有锁的线程不可能释放它，
直到某个自旋线程释放一个 CPU\@。
这是一个经典的死锁（deadlock）情况，
通过禁止在持有自旋锁时阻塞来避免这种\IX{deadlock}。

同样的约束也施加在解引用 \co{gptr} 的读者线程上：
这些线程在使用完指向的数据项之前不允许阻塞。
回到\cref{fig:defer:Deletion With Concurrent Readers}的第二行，
更新者刚刚完成执行 \co{smp_store_release()}，
想象 CPU~0 执行了一次上下文切换。
因为读者在遍历链表时不允许阻塞，
我们保证所有可能在 CPU~0 上运行的先前读者都已完成。
将这条推理线扩展到其他 CPU，一旦每个 CPU
都被观察到执行了上下文切换，
我们就保证所有先前的读者都已完成，
不再有任何读者线程引用新移除的数据元素。
然后更新者可以安全释放该数据元素，
产生\cref{fig:defer:Deletion With Concurrent Readers}底部所示的状态。

\begin{figure}
\centering
\resizebox{3in}{!}{\includegraphics{defer/QSBRGracePeriod}}
\caption{QSBR\@：等待先前存在的读者}
\label{fig:defer:QSBR: Waiting for Pre-Existing Readers}
\end{figure}

此方法被称为\IXacrfst{qsbr}~\cite{ThomasEHart2006a}。
\cref{fig:defer:QSBR: Waiting for Pre-Existing Readers}中
展示了 QSBR 示意图，时间从图的顶部向底部推进。
青色方框描绘了 RCU 读侧临界区，
每个临界区从 \co{rcu_read_lock()} 开始，
到 \co{rcu_read_unlock()} 结束。
CPU~1 执行移除当前数据项的 \co{WRITE_ONCE()}
（大概之前已读取指针值并利用了适当的同步），
然后等待读者。
此等待操作导致立即上下文切换，
这是一个 quiescent state（由粉色圆圈表示），
这反过来意味着 CPU~1 上所有先前的读取都已完成。
接下来，CPU~2 执行上下文切换，
因此 CPU~1 和~2 上的所有读者现在都已知完成。
最后，CPU~3 执行上下文切换。
此时，整个系统中的所有读者都已知完成，
所以 grace period 结束，
允许 \co{synchronize_rcu()} 返回给其调用者，
进而允许 CPU~1 释放旧数据项。

\QuickQuiz{
	在\cref{fig:defer:QSBR: Waiting for Pre-Existing Readers}中，
	CPU~3 可能访问旧数据项的最后一个读者
	在 grace period 开始之前就结束了！
	那为什么还要等到 CPU~3 后来的上下文切换？？？
}\QuickQuizAnswer{
	因为那种等待正是使读者能够使用
	适合单线程情况的相同指令序列的原因。
	换句话说，这种额外的``冗余''等待
	实现了出色的读侧性能、可扩展性和实时响应。
}\QuickQuizEnd

\subsubsection{玩具实现}
\label{sec:defer:Toy Implementation}

虽然生产质量的 QSBR 实现可能相当复杂，
但一个玩具级的非抢占式 Linux 内核实现非常简单：

\begin{VerbatimN}[samepage=true]
void synchronize_rcu(void)
{
	int cpu;

	for_each_online_cpu(cpu)
		sched_setaffinity(current->pid, cpumask_of(cpu));
}
\end{VerbatimN}

\co{for_each_online_cpu()} 原语遍历所有 CPU，
\co{sched_setaffinity()} 函数使当前线程在指定的 CPU 上执行，
这强制目标 CPU 执行上下文切换。
因此，一旦 \co{for_each_online_cpu()} 完成，
每个 CPU 都执行了上下文切换，
这反过来保证所有先前存在的读者线程都已完成。

\begin{listing}
\begin{fcvlabel}[ln:defer:Insertion and Deletion With Concurrent Readers]
\begin{VerbatimL}[commandchars=\\\[\]]
struct route *gptr;

int access_route(int (*f)(struct route *rp))
{
	int ret = -1;
	struct route *rp;

	rcu_read_lock();
	rp = rcu_dereference(gptr);
	if (rp)
		ret = f(rp);		\lnlbl[access_rp]
	rcu_read_unlock();
	return ret;
}

struct route *ins_route(struct route *rp)
{
	struct route *old_rp;

	spin_lock(&route_lock);
	old_rp = gptr;
	rcu_assign_pointer(gptr, rp);
	spin_unlock(&route_lock);
	return old_rp;
}

int del_route(void)
{
	struct route *old_rp;

	spin_lock(&route_lock);
	old_rp = gptr;
	RCU_INIT_POINTER(gptr, NULL);
	spin_unlock(&route_lock);
	synchronize_rcu();
	free(old_rp);
	return !!old_rp;
}
\end{VerbatimL}
\end{fcvlabel}
\caption{并发读者下的插入和删除}
\label{lst:defer:Insertion and Deletion With Concurrent Readers}
\end{listing}

请注意此方法\emph{不是}生产质量的。
正确处理许多边界情况以及大量强大优化的需要
意味着生产质量的实现相当复杂。
此外，可抢占环境的 RCU 实现要求读者实际做些事情，
在非实时 Linux 内核环境中，这可以简单到将
\co{rcu_read_lock()} 和 \co{rcu_read_unlock()}
分别定义为 \co{preempt_disable()} 和 \co{preempt_enable()}。\footnote{
	一些处理被抢占的读侧临界区的玩具 RCU 实现
	见\cref{chp:app:``Toy'' RCU Implementations}\@。}
然而，这个简单的非抢占式方法在概念上是完整的，
并且证明了即使面对并发更新，
也确实可以零成本提供读侧同步。
事实上，\cref{lst:defer:Insertion and Deletion With Concurrent Readers}
展示了如何实现读取（\co{access_route()}）、
\cref{fig:defer:Insertion With Concurrent Readers}的
插入（\co{ins_route()}）和
\cref{fig:defer:Deletion With Concurrent Readers}的
删除（\co{del_route()}）。
（稍微更强大的路由表在
\cref{sec:defer:RCU for Pre-BSD Routing}中展示。）

\QuickQuizSeries{%
\QuickQuizB{
	\cref{lst:defer:Insertion and Deletion With Concurrent Readers}中
	\co{rcu_read_lock()} 和 \co{rcu_read_unlock()} 有什么意义？
	为什么不让 quiescent states 自己说话？
}\QuickQuizAnswerB{
	回想一下，读者不允许通过 quiescent state。
	例如，在 Linux 内核中，RCU 读者不允许执行上下文切换。
	使用 \co{rcu_read_lock()} 和 \co{rcu_read_unlock()}
	可以对不正确放置的 quiescent states 进行调试检查，
	使得容易找到否则难以发现的、间歇性的且相当具有破坏性的错误。
}\QuickQuizEndB
%
\QuickQuizE{
	\cref{lst:defer:Insertion and Deletion With Concurrent Readers}中
	\co{rcu_dereference()}、\co{rcu_assign_pointer()} 和
	\co{RCU_INIT_POINTER()} 有什么意义？
	为什么不直接分别使用 \co{READ_ONCE()}、\co{smp_store_release()}
	和 \co{WRITE_ONCE()}？
}\QuickQuizAnswerE{
	RCU 特定的 API 确实与建议的替代品具有类似的语义，
	但也启用了静态分析调试检查，
	如果对非 RCU 指针调用 RCU 特定 API 或反之则会报错。
}\QuickQuizEndE
}

回顾\cref{lst:defer:Insertion and Deletion With Concurrent Readers}，
注意 \co{route_lock} 用于同步调用 \co{ins_route()} 和
\co{del_route()} 的并发更新者。
但是，调用 \co{access_route()} 的读者不获取此锁：
读者由\cref{sec:defer:Waiting for Readers}中描述的
QSBR 技术保护。

注意 \co{ins_route()} 只是返回 \co{gptr} 的旧值，
\cref{fig:defer:Insertion With Concurrent Readers}假设它总是 \co{NULL}。
这意味着调用者有责任弄清楚如何处理非 \co{NULL} 值，
这个任务因为读者可能在不确定的时间段内仍在引用它而变得复杂。
调用者可以使用以下方法之一：

\begin{enumerate}
\item	使用 \co{synchronize_rcu()} 安全地释放指向的结构。
	虽然此方法从 RCU 的角度来看是正确的，
	但可以说存在软件工程方面的泄漏 API 问题。
\item	如果返回的指针非 \co{NULL}，触发断言。
\item	将返回的指针传递给后续的 \co{ins_route()} 调用
	以恢复先前的值。
\end{enumerate}

相比之下，\co{del_route()} 使用 \co{synchronize_rcu()} 和
\co{free()} 安全地释放新删除的数据项。

\QuickQuiz{
	但如果旧结构需要被释放，而 \co{ins_route()} 的调用者
	不能阻塞，也许是由于性能考虑，
	也许是因为调用者在 RCU 读侧临界区内执行，怎么办？
}\QuickQuizAnswer{
	\co{call_rcu()} 函数在
	\cref{sec:defer:Wait For Pre-Existing RCU Readers}
	中描述，允许异步 grace-period 等待。
}\QuickQuizEnd

此示例展示了读取和更新 RCU 保护的数据结构的一般方法，
然而存在相当多样的使用场景，
其中几个在\cref{sec:defer:RCU Usage}中介绍。

总之，确实可以创建并发链式数据结构，
使读者执行的机器指令序列与单线程读者执行的完全相同。
下一节总结 RCU 的高级属性。

\subsubsection{RCU 属性}
\label{sec:defer:RCU Properties}

RCU 的一个关键属性是读操作不需要等待更新操作。
此属性使 RCU 实现能够提供低成本甚至零成本的读者，
从而实现低开销和出色的可扩展性。
此属性还允许 RCU 读者和更新者进行有用的并发前进。
相比之下，传统同步原语必须使用昂贵的指令强制严格的互斥，
从而增加开销并降低可扩展性，
但通常也阻止读者和更新者进行有用的并发前进。

\QuickQuiz{
	\cref{sec:defer:Sequence Locks}的 seqlock
	是否也允许读者和更新者进行有用的并发前进？
}\QuickQuizAnswer{
	是也不是。
	虽然 seqlock 读者可以与 seqlock 写者并发运行，
	但每当这发生时，\co{read_seqretry()} 原语将迫使读者重试。
	这意味着与 seqlock 更新者并发运行的 seqlock 读者所做的
	任何工作都将被丢弃，然后在重试时重做。
	所以 seqlock 读者可以与更新者\emph{并发运行}，
	但在这种情况下实际上不能完成任何工作。

	相比之下，RCU 读者即使在并发 RCU 更新者存在的情况下
	也能执行有用的工作。

	然而，reference counters
	（\cref{sec:defer:Reference Counting}）
	和 hazard pointers
	（\cref{sec:defer:Hazard Pointers}）
	确实允许更新者和读者都进行有用的并发前进，
	只是成本稍高。
	请参见\cref{sec:defer:Which to Choose?}
	以比较这些不同的延迟回收问题解决方案。
}\QuickQuizEnd

如前所述，RCU 用 \co{rcu_read_lock()} 和
\co{rcu_read_unlock()} 界定读者，
通过维护对象的多个版本并使用更新侧原语（如 \co{synchronize_rcu()}）
确保对象在可能使用它们的所有读者完成之前不被释放，
从而确保每个读者对每个对象有一致的视图
（参见\cref{fig:defer:Deletion With Concurrent Readers}）。
RCU 使用 \co{rcu_assign_pointer()} 和 \co{rcu_dereference()}
分别提供发布和读取对象新版本的高效且可扩展的机制。
这些机制在读路径和更新路径之间分配工作，
使读路径极其快速，以类似于\IXpl{hazard pointer}的方式
使用复制和弱化优化，但不需要读侧重试。
在某些情况下，包括 \co{CONFIG_PREEMPT=n} 的 Linux 内核，
RCU 的读侧原语有零开销。

但这些属性在实践中真的有用吗？
下一节将探讨这个问题。

\subsubsection{实际适用性}
\label{sec:defer:Practical Applicability}

\begin{figure}
\centering
\resizebox{3in}{!}{\includegraphics{defer/linux-RCU}}
\caption{Linux 内核中的 RCU 使用}
\label{fig:defer:RCU Usage in the Linux Kernel}
\end{figure}

RCU 从 2002 年 10 月起就在 Linux 内核中
使用~\cite{Torvalds2.5.43}。
此后，RCU API 的使用显著增加，
如\cref{fig:defer:RCU Usage in the Linux Kernel}所示。
RCU 在被 Linux 内核接受之前和之后都被大量使用，
如\cref{sec:defer:RCU Related Work}中所讨论。
简而言之，RCU 具有广泛的实际适用性。

本节讨论的最小示例是 RCU 的良好入门。
然而，有效使用 RCU 通常需要你以不同的方式思考你的问题。
因此，审视 RCU 的基础知识是有用的，
下一节将承担这个任务。

% defer/rcufundamental.tex
% mainfile: ../perfbook.tex
% SPDX-License-Identifier: CC-BY-SA-3.0

\subsection{RCU基础}
\label{sec:defer:RCU Fundamentals}
\OriginallyPublished{sec:defer:RCU Fundamentals}{RCU Fundamentals}{Linux Weekly News}{PaulEMcKenney2007WhatIsRCUFundamentally}

本节重新审视前一节涵盖的内容，但不依赖于任何特定的示例或用例。
喜欢贴近实际代码的读者可以跳过本节所介绍的底层基础知识。

RCU保护链式数据结构的常见用法由三个基本机制组成，第一个用于插入，
第二个用于删除，第三个用于允许读者容忍并发的插入和删除操作。
\Cref{sec:defer:Publish-Subscribe Mechanism}
描述了用于插入的publish-subscribe机制，
\cref{sec:defer:Wait For Pre-Existing RCU Readers}
描述了等待已有RCU读者如何实现删除，\footnote{
	是的，RCU就是这样，这个最基本的RCU组件不是第一个介绍，
	而是第二个。}
而
\cref{sec:defer:Maintain Multiple Versions of Recently Updated Objects}
讨论了维护最近更新对象的多个版本如何允许并发的插入和删除操作。
最后，
\cref{sec:defer:Summary of RCU Fundamentals}
总结了RCU的基础知识。

\subsubsection{发布-订阅机制}
\label{sec:defer:Publish-Subscribe Mechanism}

由于RCU读者不会被RCU更新者排斥，RCU保护的数据结构可能在读者
访问时发生变化。
被访问的数据项可能被移动、移除或替换。
因为数据结构不会为读者"保持静止"，所以每个读者的访问可以被视为
订阅（subscribe）RCU保护数据项的当前版本。
就更新者而言，它们可以被视为发布（publish）新版本。

% @@@ Merge usage section into "Which to Choose?" and suggest choices.

\begin{figure}
\centering
\resizebox{3in}{!}{\includegraphics{defer/pubsub}}
\caption{发布/订阅约束}
\label{fig:defer:Publication/Subscription Constraints}
\end{figure}

不幸的是，正如
\cref{sec:toolsoftrade:Shared-Variable Shenanigans}
中所述以及
\cref{sec:defer:Minimal Insertion and Deletion}
中重申的那样，
对这些发布和订阅操作使用\IXplx{plain access}{es}是不明智的。
相反，必须告知编译器和CPU需要小心处理，
如\cref{fig:defer:Publication/Subscription Constraints}
所示，该图说明了
\cref{lst:defer:Insertion and Deletion With Concurrent Readers}
中\co{ins_route()}（及其调用者）和\co{access_route()}并发执行之间的交互。

\cref{fig:defer:Publication/Subscription Constraints}
中的\co{ins_route()}列显示了\co{ins_route()}的调用者分配一个新的
\co{route}结构，此时该结构包含的是初始化前的垃圾数据。
调用者随后初始化新分配的结构，然后调用\co{ins_route()}发布指向
新\co{route}结构的指针。
发布不会影响结构的内容，因此结构内容在发布后仍然有效。

同一图中的\co{access_route()}列显示了指针被订阅和解引用的过程。
这个解引用操作绝对必须看到一个有效的\co{route}结构，而不是初始化
前的垃圾数据，因为引用垃圾数据可能导致内存损坏、崩溃和挂起。
如前所述，避免此类垃圾数据意味着发布和订阅操作必须告知编译器和CPU
维持所需排序的必要性。

发布通过\co{rcu_assign_pointer()}执行，它确保\co{ins_route()}调用者的
初始化操作排在实际发布操作存储指针之前。
此外，\co{rcu_assign_pointer()}必须是原子的，即并发读者看到的要么是
指针的旧值，要么是新值，而不是两者的某种混合。
这些要求由C11 store-release操作满足，实际上在Linux内核中，
\co{rcu_assign_pointer()}是用\co{smp_store_release()}定义的，
后者类似于C11 store-release。

注意，如果需要并发更新，则需要某种同步机制来协调同一指针上的
多个并发\co{rcu_assign_pointer()}调用。
在Linux内核中，锁是首选机制，但几乎任何同步机制都可以使用。
一个特别轻量级的同步机制的例子是
\cref{chp:Data Ownership}的数据所有权（data ownership）：
如果每个指针由特定线程拥有，则该线程可以在没有额外同步开销的情况下
对该指针执行\co{rcu_assign_pointer()}。

\QuickQuiz{
	对RCU更新者使用数据所有权是否意味着更新可以使用与
	相应单线程代码完全相同的指令序列？
}\QuickQuizAnswer{
	有时候是的，例如在TSO系统（如x86或IBM大型机）上，
	store-release操作只发出一条存储指令。
	然而，弱排序系统还必须发出memory barrier或者
	store-release指令。
	此外，移除数据需要相当多的额外工作，因为在释放被移除的数据
	之前必须等待已有的读者完成。
}\QuickQuizEnd

订阅通过\co{rcu_dereference()}执行，它防止编译器和（在一种情况下）
CPU将指针加载与后续解引用该指针的内存操作重排序。
与\co{rcu_assign_pointer()}类似，\co{rcu_dereference()}必须是原子的，
即加载的值必须来自单次存储，例如，编译器不能拆分该加载。\footnote{
	也就是说，编译器不能将该加载拆分为多个更小的加载，
	如\cref{sec:toolsoftrade:Shared-Variable Shenanigans}中
	"load tearing"下所述。}
不幸的是，编译器对\co{rcu_dereference()}的支持充其量还在进行中~\cite{PaulEMcKennneyConsumeP0190R4,PaulEMcKenney2017markconsumeP0462R1,JFBastien2018P0750R1consume}。
在此期间，Linux内核依赖volatile加载、各种CPU架构的细节、
编码限制~\cite{PaulEMcKenney2014rcu-dereference}，
以及在DEC Alpha~\cite{ALPHA2002}上的一条memory-barrier指令。
然而，在其他架构上，\co{rcu_dereference()}通常只发出一条加载指令，
与等效的单线程代码相同。
编码限制在
\cref{sec:memorder:Address- and Data-Dependency Difficulties}
中有更详细的描述，不过，常见的字段选择（\qtco{->}）情况工作得很好。
不需要极致读侧性能的软件可以改用C11 \IXpl{acquire load}，
它提供了所需的排序及更多保证，但需要付出一定代价。
希望更轻量级的\co{rcu_dereference()}编译器支持将在适当时候出现。

简而言之，使用\co{rcu_assign_pointer()}发布指针和使用
\co{rcu_dereference()}订阅指针成功避免了
\cref{fig:defer:Publication/Subscription Constraints}
中描述的"Not OK"垃圾加载。
因此，这两个原语可以用于向链式结构添加新数据，而不会干扰并发读者。

\QuickQuiz{
	但假设更新者正在向链表中添加和移除多个数据项，
	同时一个读者正在遍历同一链表。
	具体来说，假设链表最初包含元素A、B和~C，
	更新者移除了元素A，然后在链表末尾添加了新元素D。
	读者很可能看到\{A, B, C, D\}，而这个元素序列实际上
	从未存在过！
	在哪个平行宇宙中这算得上"不干扰并发读者"？？？
}\QuickQuizAnswer{
	在这样一个世界中：遍历读者只需要遍历在整个遍历期间
	一直存在的元素。
	在本例中，那就是元素B和~C\@。
	因为元素A和~D各自只在遍历的一部分时间内存在，读者可以
	遍历它们，但并非必须遍历。
	注意，这支持了常见的情况，即读者只是查找单个项目，
	并不知道或关心其他项目的存在与否。

	如果需要更强的一致性，则需要更高代价的同步机制，
	例如sequence locking或读写锁。
	但如果不需要更强的一致性（实际上很多时候不需要），
	为什么要付出更高的代价呢？
}\QuickQuizEnd

在不干扰读者的情况下向链式结构添加数据是一件好事，
与单线程读者相比可以不增加读侧开销的情况更是如此。
然而，在大多数情况下还需要移除数据，这是下一节的主题。

\subsubsection{等待已有的RCU读者}
\label{sec:defer:Wait For Pre-Existing RCU Readers}

RCU最基本的组件是等待事情完成。
当然，还有很多其他方式来等待事情完成，包括reference counting、
读写锁、事件等。
RCU的巨大优势在于它可以等待（比如说）20,000个不同的事情完成，
而不必显式跟踪每一个，也不必担心显式跟踪方案固有的性能下降、
可扩展性限制、复杂的死锁场景和内存泄漏风险。

在RCU的情况下，等待的每一件事称为
\emph{\IXBhmr{RCU read-side}{critical section}}。
如\cref{tab:defer:Core RCU API}中所述，
RCU read-side critical section以\co{rcu_read_lock()}原语开始，
以相应的\co{rcu_read_unlock()}原语结束。
RCU read-side critical section可以嵌套，并且可以包含几乎任何代码，
只要该代码不包含quiescent state。
例如，在Linux内核中，在RCU read-side critical section中睡眠是
不允许的，因为上下文切换是一个quiescent state。\footnote{
	然而，一种称为SRCU~\cite{PaulEMcKenney2006c}的特殊RCU形式
	确实允许在SRCU read-side critical section中进行一般性睡眠。}
如果你遵守这些约定，你可以使用RCU来等待\emph{任何}已有的
RCU read-side critical section完成，而
\co{synchronize_rcu()}使用间接方式来执行实际的
等待~\cite{MathieuDesnoyers2012URCU,McKenney:2013:SDS:2483852.2483867}。

\begin{figure}
\centering
\resizebox{3in}{!}{\includegraphics{defer/RCUGuaranteeFwd}}
\caption{RCU读者与后续Grace Period}
\label{fig:defer:RCU Reader and Later Grace Period}
\end{figure}

RCU read-side critical section与后续RCU grace period之间的关系
是一种"如果-那么"关系，如
\cref{fig:defer:RCU Reader and Later Grace Period}所示。
如果给定critical section的任何部分先于给定grace period的开始，
那么RCU保证该critical section的全部将先于该grace period的结束。
在图中，\co{P0()}对\co{x}的访问先于\co{P1()}对同一变量的访问，
因此也先于\co{P1()}调用\co{synchronize_rcu()}产生的grace period。
因此可以保证\co{P0()}对\co{y}的访问将先于\co{P1()}的访问。
在这种情况下，如果\co{r1}的最终值为0，则\co{r2}的最终值
也保证为0。

\QuickQuiz{
	在\cref{fig:defer:RCU Reader and Later Grace Period}中，
	\co{r1}和\co{r2}还有哪些其他可能的最终值？
}\QuickQuizAnswer{
	\co{r1 == 0 && r2 == 0}的可能性已在正文中指出。
	鉴于\co{r1 == 0}意味着\co{r2 == 0}，我们知道
	\co{r1 == 0 && r2 == 1}是被禁止的。
	以下讨论将表明\co{r1 == 1 && r2 == 1}和
	\co{r1 == 1 && r2 == 0}都是可能的。
}\QuickQuizEnd

\begin{figure}
\centering
\resizebox{3in}{!}{\includegraphics{defer/RCUGuaranteeRev}}
\caption{RCU读者与更早的Grace Period}
\label{fig:defer:RCU Reader and Earlier Grace Period}
\end{figure}

RCU read-side critical section与更早的RCU grace period之间的关系
也是一种"如果-那么"关系，如
\cref{fig:defer:RCU Reader and Earlier Grace Period}所示。
如果给定critical section的任何部分后于给定grace period的结束，
那么RCU保证该critical section的全部将后于该grace period的开始。
在图中，\co{P0()}对\co{y}的访问后于\co{P1()}对同一变量的访问，
因此后于\co{P1()}调用\co{synchronize_rcu()}产生的grace period。
因此可以保证\co{P0()}对\co{x}的访问将后于\co{P1()}的访问。
在这种情况下，如果\co{r2}的最终值为1，则\co{r1}的最终值
也保证为1。

\QuickQuiz{
	如果\co{P0()}在
	\cref{fig:defer:RCU Reader and Earlier Grace Period}
	中的两次访问顺序颠倒，会发生什么？
}\QuickQuizAnswer{
	绝对不会有任何变化。
	\co{P0()}对\co{x}和\co{y}的加载在同一个
	RCU read-side critical section中这一事实就足够了；
	它们的顺序无关紧要。
}\QuickQuizEnd

\begin{figure}
\centering
\resizebox{3in}{!}{\includegraphics{defer/RCUGuaranteeMid}}
\caption{Grace Period内的RCU读者}
\label{fig:defer:RCU Reader Within Grace Period}
\end{figure}

最后，如\cref{fig:defer:RCU Reader Within Grace Period}所示，
RCU read-side critical section可以被RCU grace period完全覆盖。
在这种情况下，\co{r1}的最终值为1，\co{r2}的最终值为0。

然而，不可能出现\co{r1}的最终值为0且\co{r2}的最终值为1的情况。
这将意味着一个RCU read-side critical section完全覆盖了一个
grace period，这是被禁止的（或者至少构成RCU的一个bug）\@。
因此，RCU的等待读者保证有两个部分：
\begin{enumerate*}[(1)]
\item 如果给定RCU read-side critical section的任何部分先于
给定grace period的开始，则该critical section的全部先于该
grace period的结束。
\item 如果给定RCU read-side critical section的任何部分后于
给定grace period的结束，则该critical section的全部后于该
grace period的开始。
\end{enumerate*}
这个定义对几乎所有基于RCU的算法都是足够的，但对于想要更多的人，
简单的可执行RCU形式化模型作为Linux内核v4.17及更高版本的一部分
可供使用，如\cref{sec:formal:Axiomatic Approaches and RCU}中所述。
此外，RCU的排序属性在\cref{sec:memorder:RCU}中有更详细的讨论。

\QuickQuiz{
	如果\co{P0()}在
	\crefrange{fig:defer:RCU Reader and Later Grace Period}{fig:defer:RCU Reader Within Grace Period}
	中的访问是存储操作，会发生什么？
}\QuickQuizAnswer{
	完全相同的排序规则适用，即：
	(1)~如果\co{P0()}的RCU read-side critical section的任何部分
	先于\co{P1()}的grace period的开始，则\co{P0()}的
	RCU read-side critical section的全部将先于\co{P1()}的
	grace period的结束，
	(2)~如果\co{P0()}的RCU read-side critical section的任何部分
	后于\co{P1()}的grace period的结束，则\co{P0()}的
	RCU read-side critical section的全部将后于\co{P1()}的
	grace period的开始。

	在RCU read-side critical section中包含写操作可能看起来
	很奇怪，但这种能力不仅是被允许的，而且非常有用。
	例如，Linux内核经常执行RCU保护的链式数据结构遍历，
	然后获取目标数据元素的引用。
	因为该数据元素在此期间不能被释放，所以该元素的引用计数
	必须在遍历的RCU read-side critical section内递增。
	然而，该递增操作涉及对内存的写入。
	因此，在RCU read-side critical section中允许内存写入
	是一件非常好的事情。

	如果在RCU read-side critical section中进行写操作仍然
	看起来很奇怪，请回顾
	\cref{sec:count:Applying Exact Limit Counters}，
	其中展示了在读写锁读侧critical section中进行写操作的
	用例。
}\QuickQuizEnd

虽然RCU的等待读者能力确实有时被用于如
\crefrange{fig:defer:RCU Reader and Later Grace Period}
{fig:defer:RCU Reader Within Grace Period}
所示的变量赋值排序，但它更常被用于安全释放从链式结构中
移除的数据元素，正如在
\cref{sec:defer:Introduction to RCU}中所做的那样。
一般过程由以下伪代码说明：

\begin{enumerate}
\item	进行更改，例如，从链表中移除一个元素。
\item	等待所有已有的RCU read-side critical section完全完成
	（例如，通过使用\co{synchronize_rcu()}）。
\item	清理，例如，释放上面被替换的元素。
\end{enumerate}

\begin{figure}
\centering
\IfEbookSize{
\resizebox{2in}{!}{\includegraphics{defer/RCUGPorderingSummary}}
}{
\resizebox{2.5in}{!}{\includegraphics{defer/RCUGPorderingSummary}}
}
\caption{RCU Grace-Period排序保证总结}
\label{fig:defer:Summary of RCU Grace-Period Ordering Guarantees}
\end{figure}

这个更抽象的过程需要比
\crefrange{fig:defer:RCU Reader and Later Grace Period}
{fig:defer:RCU Reader Within Grace Period}
更抽象的图，因为那些图是针对特定石蕊测试的。
毕竟，无论RCU更新和RCU read-side critical section的形式如何，
RCU实现都必须正确工作。
\Cref{fig:defer:Summary of RCU Grace-Period Ordering Guarantees}
满足了这一需求，展示了四种可能的场景，每个场景中时间从上到下推进。
在每个场景中，RCU读者由左侧的方框栈表示，RCU更新者由右侧的
方框栈表示。

在第一个场景中，读者在更新者开始移除之前开始执行，因此该读者
可能持有对被移除数据元素的引用。
因此，更新者在读者完成之前不能释放此元素。
在第二个场景中，读者直到移除完成之后才开始执行。
该读者不可能获得对已移除数据元素的引用，因此可以在读者完成之前
释放该元素。
第三个场景类似于第二个，但说明了即使读者不可能获得对某元素的引用，
仍然允许将该元素的释放推迟到读者完成之后。
在第四个也是最后一个场景中，读者在更新者开始移除数据元素之前
开始执行，但该元素（错误地）在读者完成之前被释放。
正确的RCU实现不会允许第四种场景发生。
因此，此图说明了RCU的等待读者功能：
给定一个grace period，每个读者要么在该grace period结束之前结束，
要么在该grace period开始之后开始，或者两者兼有，
在这种情况下它完全包含在该grace period内。

因为RCU读者可以在更新进行的同时向前推进，不同的读者可能对
数据结构的状态持不同意见，这一主题将在下一节中讨论。

\subsubsection{维护最近更新对象的多个版本}
\label{sec:defer:Maintain Multiple Versions of Recently Updated Objects}

本节讨论RCU如何通过维护数据的多个版本来适应无同步读者。
因为这些无同步读者提供的时间同步非常弱，RCU用户通过空间同步
来补偿。
空间同步在\cref{chp:Partitioning and Synchronization Design}中讨论过，
并在实践中被大量使用以获得良好的性能和可扩展性。
在本节中，空间同步将被用来实现一种弱（但有用的）正确性形式，
以及出色的性能和可扩展性。

\cref{sec:defer:Minimal Insertion and Deletion}中的
\Cref{fig:defer:Deletion With Concurrent Readers}
展示了空间同步的一个简单变体，其中与\co{del_route()}
（见\cref{lst:defer:Insertion and Deletion With Concurrent Readers}）
并发运行的不同读者可能看到旧的\co{route}结构或空列表，
但无论哪种方式都能得到有效结果。
当然，仔细查看
\cref{fig:defer:Insertion With Concurrent Readers}
可以发现，\co{ins_route()}的调用也可能导致并发读者看到不同版本：
要么是初始的空列表，要么是新插入的\co{route}结构。
注意，reference counting
（\cref{sec:defer:Reference Counting}）
和hazard pointer
（\cref{sec:defer:Hazard Pointers}）
也可能导致并发读者看到不同版本，但RCU的轻量级读者使得这种情况
更为可能。

\begin{figure}
\centering
\resizebox{3in}{!}{\includegraphics{defer/multver}}
\caption{RCU数据结构的多个版本}
\label{fig:defer:Multiple RCU Data-Structure Versions}
\end{figure}

然而，维护多个弱一致性版本可能带来一些意外。
例如，考虑
\cref{fig:defer:Multiple RCU Data-Structure Versions}，
其中一个读者正在遍历一个被并发更新的链表。\footnote{
	RCU链表API可在
	\cref{sec:defer:RCU Linux-Kernel API}中找到。}
在图的第一行，读者引用数据项~A，
在第二行，它前进到~B，到目前为止看到的是A后面跟着~B\@。
在第三行，更新者移除了元素~A，在第四行，
更新者在列表末尾添加了元素~E。
在第五行也就是最后一行，读者完成了遍历，
看到的元素为~A到~E\@。

只是这样的列表从未在任何时刻存在过。
这种情况可能比
\cref{fig:defer:Deletion With Concurrent Readers}
所示的情况更令人惊讶，在那里不同的并发读者看到不同版本。
相比之下，在
\cref{fig:defer:Multiple RCU Data-Structure Versions}
中，读者看到的是一个从未实际存在过的版本！

解决这种奇怪情况的一种方法是通过较弱的语义。
读者遍历必须遇到在整个遍历期间都存在的任何数据项
（B、C和~D），但可能遇到也可能不遇到只在遍历的一部分时间
存在的数据项（A和~E）\@。
因此，在这种特定情况下，读者遍历遇到所有五个元素是完全合法的。
如果这个结果有问题，解决这种情况的另一种方法是使用更强的
同步机制，如读写锁，或巧妙地使用时间戳和版本控制，
如\cref{sec:defer:Quasi Multi-Version Concurrency Control}中所述。
当然，更强的机制会更昂贵，但工程生活就是关于选择和权衡的。

虽然这种情况看起来很奇怪，但它与现实世界完全一致。
正如我们在\cref{sec:cpu:Overheads}中看到的，
光速有限这一事实在计算机系统内不能被忽略，
在系统外部更不能被忽略。
这反过来意味着系统内表示系统外部真实世界状态的任何数据
始终是过时的，因此与真实世界不一致。
因此，序列\{A, B, C, D, E\}完全有可能在真实世界中发生过，
但由于光速延迟而从未在计算机系统的内存中体现。
在这种情况下，读者令人惊讶的遍历结果将正确反映现实。

因此，操作真实世界数据的算法必须考虑不一致的数据，
要么容忍不一致性，要么采取步骤排除或拒绝它们。
在许多情况下，这些算法也完全能够处理系统内部的不一致性。

\cref{sec:defer:Running Example}中介绍的pre-BSD数据包路由
示例就是一个很好的例证。
路由列表的内容由路由协议设置，这些协议具有显著的延迟
（秒甚至分钟级别），以避免路由不稳定。
因此，一旦路由更新到达给定系统，它可能已经将数据包发送到
错误的方向相当长时间了。
在更新传播的几微秒内再多发送几个错误方向的数据包显然不是问题，
因为处理延迟路由更新的相同高层协议操作也会处理内部不一致性。

互联网路由也不是唯一容忍不一致性的情况。
重复一下，任何系统中的数据跟踪系统外部状态的算法都必须
容忍不一致性，这包括安全策略（通常由人类委员会制定）、
存储配置和WiFi接入点，更不用说可移除硬件如麦克风、耳机、
摄像头、鼠标、打印机等等。
此外，\cref{fig:defer:RCU Usage in the Linux Kernel}所示的
大量Linux内核RCU API使用，
加上Linux内核大量使用reference counting以及其他项目越来越多地
使用\IXpl{hazard pointer}，表明对此类不一致性的容忍比人们
想象的更为常见。

这种常见情况下对不一致性的容忍的一个根本原因是，
单项查找在实践中比全数据结构遍历更为常见。
毕竟，全数据结构遍历比单项查找要昂贵得多，
因此开发者有动力避免此类遍历。
并发更新不仅比全遍历更不容易影响单项查找，
而且孤立的单项查找也无法检测到此类不一致性。
因此，在常见情况下，此类不一致性不仅是可容忍的，
实际上是不可见的。

在这种情况下，RCU读者可以被认为与更新者完全有序，
尽管这些读者可能正在执行与单线程程序完全相同的机器指令序列，
正如在\cpageref{sec:defer:Mysteries RCU}中所暗示的。
例如，回顾
\cpageref{lst:defer:Insertion and Deletion With Concurrent Readers}上的
\cref{lst:defer:Insertion and Deletion With Concurrent Readers}，
假设每个读者线程在其生命周期中恰好调用一次\co{access_route()}，
并且读者和更新者线程之间没有其他通信。
\begin{fcvref}[ln:defer:Insertion and Deletion With Concurrent Readers]
那么每次\co{access_route()}的调用可以排在产生该\co{route}结构
的\co{ins_route()}调用之后（该结构由清单中\co{access_route()}的
\clnref{access_rp}行访问），并排在任何后续的
\co{ins_route()}或\co{del_route()}调用之前。
\end{fcvref}

总之，维护多个版本正是使RCU读者具有极低开销的关键，
而且如前所述，许多算法不受多版本的影响。
然而，确实有些算法完全无法处理多个版本。
有一些技术可以使此类算法适应
RCU~\cite{PaulEdwardMcKenneyPhD}，例如
\cref{sec:together:Correlated Data Elements}中描述的sequence locking的
使用。

\paragraph{练习}
\label{sec:defer:Exercises}

这些示例假设在整个更新操作期间持有互斥锁（exclusive lock），
这意味着在给定时间最多有两个版本的列表处于活动状态。

\QuickQuizSeries{%
\QuickQuizB{
	你将如何修改删除示例以允许列表有两个以上的
	活动版本？
}\QuickQuizAnswerB{
	实现这一目标的一种方式如
	\cref{lst:defer:Concurrent RCU Deletion}所示。

\begin{listing}
\begin{VerbatimL}
spin_lock(&mylock);
p = search(head, key);
if (p == NULL)
	spin_unlock(&mylock);
else {
	list_del_rcu(&p->list);
	spin_unlock(&mylock);
	synchronize_rcu();
	kfree(p);
}
\end{VerbatimL}
\caption{并发RCU删除}
\label{lst:defer:Concurrent RCU Deletion}
\end{listing}

	注意，这意味着多个并发删除可能在
	\co{synchronize_rcu()}中等待。
}\QuickQuizEndB
%
\QuickQuizM{
	在任何给定时间，给定列表可以有多少个RCU版本
	处于活动状态？
}\QuickQuizAnswerM{
	这取决于同步设计。
	如果保护更新的信号量在grace period期间一直持有，
	则最多可以有两个版本——旧版本和新版本。

	然而，假设只有搜索、更新和\co{list_replace_rcu()}
	受锁保护，使得\co{synchronize_rcu()}在该锁之外，
	类似于\cref{lst:defer:Concurrent RCU Deletion}所示的代码。
	进一步假设大量线程大约在同一时间进行RCU替换，
	并且读者也在不断遍历数据结构。

	那么以下事件序列可能发生，从
	\cref{fig:defer:Multiple RCU Data-Structure Versions}
	的结束状态开始：

	\begin{enumerate}
	\item	线程~A遍历列表，获取对元素~C的引用。
	\item	线程~B将元素~C替换为新元素~F，
		然后等待其\co{synchronize_rcu()}调用返回。
	\item	线程~C遍历列表，获取对元素~F的引用。
	\item	线程~D将元素~F替换为新元素~G，
		然后等待其\co{synchronize_rcu()}调用返回。
	\item	线程~E遍历列表，获取对元素~G的引用。
	\item	线程~F将元素~G替换为新元素~H，
		然后等待其\co{synchronize_rcu()}调用返回。
	\item	线程~G遍历列表，获取对元素~H的引用。
	\item	前两步用新的元素快速重复，使得所有这些
		操作在任何\co{synchronize_rcu()}调用返回之前完成。
	\end{enumerate}

	因此，可以有任意数量的活动版本，仅受内存和在一个
	grace period内可以完成多少次更新的限制。
	但请注意，更新如此频繁的数据结构不太可能是RCU的
	好候选者\@。
	尽管如此，RCU在必要时可以处理高更新率。
}\QuickQuizEndM
%
\QuickQuizE{
	如何降低RCU每次更新的开销？
}\QuickQuizAnswerE{
	降低RCU每次更新开销最有效的方法是增加由
	给定grace period服务的更新数量。
	这之所以有效，是因为每个grace period的开销几乎
	与该grace period服务的更新数量无关。

	一种方法是延迟给定grace period的开始，
	希望在此期间有更多需要该grace period的更新出现。
	另一种方法是减慢grace period的执行速度，
	希望在此期间有更多需要额外grace period的更新积累。

	还有许多其他可能的优化，狂热的读者可以参考
	Linux内核的RCU实现。
}\QuickQuizEndE
}

\subsubsection{RCU基础总结}
\label{sec:defer:Summary of RCU Fundamentals}

本节描述了基于RCU的链式结构算法的三个基本组件：

\begin{enumerate}
\item	用于添加新数据的publish-subscribe机制，
	更新侧使用\co{rcu_assign_pointer()}进行发布，
	读侧使用\co{rcu_dereference()}进行订阅，

\item	一种等待已有RCU读者完成的方式，
	读者一侧由\co{rcu_read_lock()}和\co{rcu_read_unlock()}界定，
	更新者一侧通过\co{synchronize_rcu()}或
	\co{call_rcu()}等待
	（形式化描述见\cref{sec:memorder:RCU}），
	以及

\item	一种维护多版本的规范，以允许更改而不伤害或
	过度延迟并发RCU读者。
\end{enumerate}

\QuickQuiz{
	既然\co{rcu_read_lock()}和\co{rcu_read_unlock()}
	都不会自旋或阻塞，RCU更新者怎么可能延迟RCU读者？
}\QuickQuizAnswer{
	给定RCU更新者所进行的修改将导致相应CPU使包含数据的
	缓存行无效，迫使运行并发RCU读者的CPU产生昂贵的
	缓存未命中。
	（你能设计一种更改数据结构而\emph{不}对并发读者造成
	昂贵缓存未命中的算法吗？对后续读者呢？）
}\QuickQuizEnd

这三个RCU组件允许在并发读者可能正在执行与单线程实现中
读者完全相同的机器指令序列的情况下进行数据更新。
这些RCU组件可以以不同方式组合来实现种类繁多的不同类型的
基于RCU的算法，其中一些在
\cref{sec:defer:RCU Usage}中介绍。
然而，通常更好的做法是在更高的抽象层次上工作。
为此，下一节描述了Linux内核API，其中包括列表等简单数据结构。

% defer/rcuAPI.tex
% mainfile: ../perfbook.tex
% SPDX-License-Identifier: CC-BY-SA-3.0

\subsection{RCU Linux内核API}
\label{sec:defer:RCU Linux-Kernel API}
\OriginallyPublished{sec:defer:RCU Linux-Kernel API}{RCU Linux-Kernel API}{Linux Weekly News}{PaulEMcKenney2008WhatIsRCUAPI}

本节从Linux内核API的角度来审视RCU。\footnote{
	用户空间RCU的API在其他文献中有文档说明~\cite{PaulMcKenney2013LWNURCU}。}
\Cref{sec:defer:RCU API and Software Engineering}
探讨了RCU的API与软件工程之间的关系，
\cref{sec:defer:RCU has a Family of Wait-to-Finish APIs}
介绍了RCU的等待完成API，
\cref{sec:defer:RCU has Publish-Subscribe and Version-Maintenance APIs}
介绍了RCU的发布-订阅和版本维护API，
\cref{sec:defer:RCU has List-Processing APIs}
介绍了RCU的链表处理API，
\cref{sec:defer:RCU Has Diagnostic APIs}
介绍了RCU的诊断API，
\cref{sec:defer:Where Can RCU's APIs Be Used?}
描述了RCU的各种API可以在哪些上下文中使用。
最后，
\cref{sec:defer:So; What is RCU Really?}
给出了结论性评述。

对内核内部不太感兴趣的读者可以跳到
\cpageref{sec:defer:RCU Usage}上的
\cref{sec:defer:RCU Usage}，
但最好在阅读完下一节关于软件工程考量的内容之后再跳。

\subsubsection{RCU API与软件工程}
\label{sec:defer:RCU API and Software Engineering}

已经提前查看过
\cref{tab:defer:RCU Wait-to-Finish APIs,%
tab:defer:RCU Markers for Read-Side Critical Sections,%
tab:defer:SRCU Markers for Read-Side Critical Sections,%
tab:defer:RCU Polled Update-Side Primitives,%
tab:defer:RCU Publish-Subscribe and Version Maintenance APIs,%
tab:defer:RCU-Protected List APIs,%
tab:defer:RCU Diagnostic APIs}
的读者可能已经注意到，Linux内核API的完整列表
拥有100多个成员。
这与\cref{tab:defer:Core RCU API}中仅显示的六个
API成员形成了鲜明的（也许令人沮丧的）对比。
这种情况显然引出了``为什么这么多？？？''这个问题。

以下各节将更彻底地回答这个问题，
但在此期间本节其余部分概述了这些动机。

有一句古老的智慧："人非圣贤，孰能无过。"
这意味着RCU API中相当大一部分的目的是
提供诊断功能，最显著的是\cref{tab:defer:RCU Diagnostic APIs}中的内容，
但其他地方也有。

人为错误的重要原因之一是人脑的局限性，
例如短期记忆的有限容量。
本书中展示的玩具示例不会触及这些极限。
这是出于必要：
许多读者在学习新材料时会达到认知极限，
因此示例需要保持简单。

因此，这些示例将\co{rcu_dereference()}调用
保持在与外围\co{rcu_read_lock()}和
\co{rcu_read_unlock()}调用相同的函数中。
相比之下，现实世界的软件经常需要从
不同的函数甚至不同的翻译单元中调用这些API成员。
Linux内核RCU API因此已扩展以适应lockdep，
lockdep允许\co{rcu_dereference()}及其相关函数在未受
\co{rcu_read_lock()}保护时发出警告。
Linux内核RCU还检查一些双重释放错误、
RCU read-side critical section中的无限循环，
以及在RCU read-side critical section中尝试调用
quiescent state的情况。

现实世界软件适应人类认知局限性的另一种方式是
通过抽象。
因此，Linux内核API除了
\cref{tab:defer:Core RCU API}中面向指针的核心API之外，
还包括操作链表的成员。
Linux内核本身也提供了受RCU保护的哈希表和搜索树。

像Linux这样的操作系统内核运行在
\cref{fig:intro:Software Layers and Performance; Productivity; and Generality}
所示的软件栈"铁三角"的底部附近，
性能至关重要。
因此，许多RCU API都有专用于快速路径的特化变体，
例如，如
\cref{sec:defer:RCU has Publish-Subscribe and Version-Maintenance APIs}所讨论的，
\co{RCU_INIT_POINTER()}可以在受RCU保护的指针
被赋值为\co{NULL}时或该指针尚未被读者访问时
替代\co{rcu_assign_pointer()}使用。
使用\co{RCU_INIT_POINTER()}为编译器在选择指令和
执行优化方面提供了更多余地，从而提高了性能。

另一方面，错误使用\co{RCU_INIT_POINTER()}可能导致
无声的内存损坏，所以请务必小心！
是的，在某些情况下，内核可以检查在给定内核上下文中
是否不当使用了RCU API成员，但
\co{RCU_INIT_POINTER()}使用的约束尚无法被检查。

最后，在Linux内核中，上述人类认知的局限性
因Linux上运行的工作负载的多样性和严苛性而加剧。
截至v5.16，这已产生了不少于五种RCU风格，
每种都旨在为RCU读者和写者提供不同的性能、
可扩展性、响应时间和能效权衡。
这些RCU风格是下一节的主题。

\subsubsection{RCU拥有一系列等待完成API}
\label{sec:defer:RCU has a Family of Wait-to-Finish APIs}

对``什么是RCU''最直接的回答是RCU是一个API。
例如，Linux内核中使用的RCU实现由
\cref{tab:defer:RCU Wait-to-Finish APIs}概括，
其各列分别显示了RCU、"可睡眠"RCU（SRCU）、Tasks RCU、
Tasks RCU Rude、Tasks RCU Trace和通用API的等待读者部分，
\cref{tab:defer:RCU Markers for Read-Side Critical Sections}
显示了RCU读侧标记的完整范围，
\cref{tab:defer:SRCU Markers for Read-Side Critical Sections}
显示了SRCU读侧标记的完整范围，
\cref{tab:defer:RCU Polled Update-Side Primitives}
显示了完整的RCU轮询grace period API，
\cref{tab:defer:RCU Publish-Subscribe and Version Maintenance APIs}
显示了API~\cite{PaulEMcKenney2024RCUAPI}的发布-订阅部分。\footnote{
	此引用涵盖v6.10及更高版本。
	早期版本的Linux内核RCU API文档可在其他地方找到~\cite{PaulEMcKenney2008WhatIsRCUAPI,PaulEMcKenney2014RCUAPI,PaulEMcKenney2019RCUAPI}。}

\begin{sidewaystable*}[tbp]
\rowcolors{1}{}{lightgray}
\renewcommand*{\arraystretch}{1.3}
\centering
\caption{RCU等待完成API}
\label{tab:defer:RCU Wait-to-Finish APIs}
\scriptsize\IfEbookSize{\hspace*{-1.8in}}{\hspace*{-.125in}}
\ebresizeverb{0.7}{
\begin{tabularx}{8.5in}{>{\raggedright\arraybackslash}p{0.92in}
    >{\raggedright\arraybackslash}p{1.37in}
    >{\raggedright\arraybackslash}p{1.39in}
    >{\raggedright\arraybackslash}p{1.09in}
    >{\raggedright\arraybackslash}p{1.34in}
    >{\raggedright\arraybackslash}p{1.44in}}
\toprule
&
    {\bf RCU}：原始版 &
	{\bf SRCU}：可睡眠读者 &
	    {\bf Tasks RCU}：释放追踪蹦床 &
		{\bf Tasks RCU Rude}：释放空闲任务追踪蹦床 &
		    {\bf Tasks RCU Trace}：保护可睡眠BPF程序 \\
\midrule
{\bf 初始化和清理} &
    &
	\tco{DEFINE_SRCU()} \tco{DEFINE_STATIC_SRCU()}
	\tco{init_srcu_struct()} \tco{cleanup_srcu_struct()} &
	    &
		&
		    \\
{\bf Read-side critical section标记} &
    \tco{rcu_read_lock()}~! \tco{rcu_read_unlock()}~!  ~~~~~~~~~~~~~~~~~~
    （更多参见\cref{tab:defer:RCU Markers for Read-Side Critical Sections}。） &
	\tco{srcu_read_lock()} \tco{srcu_read_unlock()} ~~~~~~~~~~~~~~~~~
	（更多参见\cref{tab:defer:SRCU Markers for Read-Side Critical Sections}。） &
	    非空闲任务中的自愿上下文切换。 &
		禁用抢占的代码区域。 &
		    \tco{rcu_read_lock_trace()} \tco{rcu_read_unlock_trace()}
		    \tco{guard(rcu_tasks_trace)()}
		    \tco{scoped_guard(rcu_tasks_trace)} \\
{\bf 更新侧原语（同步） } &
    { \tco{synchronize_rcu()}
      \tco{synchronize_net()}
      \tco{synchronize_rcu_expedited()} } &
	\tco{synchronize_srcu()} \tco{synchronize_srcu_expedited()} &
	    \tco{synchronize_rcu_tasks()}
	    \tco{synchronize_rcu_mult()} &
		\tco{synchronize_rcu_tasks_rude()} &
		    \tco{synchronize_rcu_tasks_trace()}
		    \tco{rcu_trace_implies_rcu_gp()} \\
{\bf 更新侧原语（异步/callback） } &
    \tco{call_rcu()} ! ~~~~~~
    \tco{call_rcu_hurry()} &
	\tco{call_srcu()} &
	    \tco{call_rcu_tasks()} &
		\tco{call_rcu_tasks_rude()} &
		    \tco{call_rcu_tasks_trace()} \\
{\bf 更新侧原语（等待callback） } &
    \tco{rcu_barrier()} &
	\tco{srcu_barrier()} &
	    \tco{rcu_barrier_tasks()} &
		\tco{rcu_barrier_tasks_rude()} &
		    \tco{rcu_barrier_tasks_trace()} \\
{\bf 更新侧原语（轮询） } &
    \tco{get_state_synchronize_rcu()}
    \tco{start_poll_synchronize_rcu()}
    \tco{poll_state_synchronize_rcu()} ~~~~~~~~~~~~~~~~~~
    （更多参见\cref{tab:defer:RCU Polled Update-Side Primitives}。） &
	\tco{get_state_synchronize_srcu()}
	\tco{start_poll_synchronize_srcu()}
	\tco{poll_state_synchronize_srcu()} &
	    &
		&
		    \\
{\bf 更新侧原语（发起/等待）} &
    \tco{get_state_synchronize_rcu()}
    \tco{cond_synchronize_rcu()} &
	&
	    &
		&
		    \\
{\bf 更新侧原语（释放内存） } &
    \tco{kfree_rcu()}
    \tco{kfree_rcu_mightsleep()}
    \tco{kvfree_rcu()}
    \tco{kvfree_rcu_mightsleep()} &
	&
	    &
		&
		    \\
{\bf 类型安全内存 } &
    \tco{SLAB_TYPESAFE_BY_RCU} &
	&
	    &
		&
		    \\
{\bf 读侧约束 } &
    不可阻塞（仅可抢占） &
	不可对同一\tco{srcu_struct}使用\tco{synchronize_srcu()} &
	    不可自愿上下文切换 &
		既不可阻塞也不可抢占 &
			不可使用RCU tasks trace grace period \\
{\bf 读侧开销 } &
    CPU本地访问（\tco{PREEMPT=n}时为\tco{barrier()}） &
	简单指令、内存屏障 &
	    免费 &
		CPU本地访问（\tco{PREEMPT=n}时免费） &
		    CPU本地访问 \\
{\bf 异步更新侧开销 } &
    亚微秒级 &
	亚微秒级 &
	    亚微秒级 &
		亚微秒级 &
		    亚微秒级 \\
{\bf Grace period延迟 } &
    数十毫秒 &
        毫秒级 &
	    秒级 &
		毫秒级 &
		    数十毫秒 \\
{\bf 加速grace period延迟 } &
    数十微秒 &
        微秒级 &
	    不适用 &
		不适用 &
		    不适用 \\
\bottomrule
\end{tabularx}
}
\end{sidewaystable*}

\begin{table*}
\rowcolors{1}{}{lightgray}
\renewcommand*{\arraystretch}{1.15}
\footnotesize
\centering\OneColumnHSpace{-.4in}
\ebresizewidth{
\begin{tabular}{llp{2.0in}}
\toprule
获取 &
	释放 &
		说明 \\
\midrule
\tco{rcu_read_lock()} &
	\tco{rcu_read_unlock()} & \\
\tco{guard(rcu)()} & & RCU读者持续到外围作用域结束。 \\
\tco{scoped_guard(rcu)} & & RCU读者用于紧随其后的语句。 \\
\tco{rcu_read_lock_bh()} &
	\tco{rcu_read_unlock_bh()} & \\
\tco{rcu_read_lock_sched()} &
	\tco{rcu_read_unlock_sched()} & \\
\tco{rcu_read_lock_sched_notrace()} &
	\tco{rcu_read_unlock_sched_notrace()} & \\
\tco{local_bh_disable()} &
	\tco{local_bh_enable()} &
		以及其他任何禁用下半部的操作。 \\
\tco{preempt_disable()} &
	\tco{preempt_enable()} &
		以及其他任何禁用抢占的操作。 \\
\tco{local_irq_save()} &
	\tco{local_irq_restore()} &
		以及其他任何禁用中断的操作。 \\
\bottomrule
\end{tabular}
}
\caption{RCU Read-Side Critical Section标记}
\label{tab:defer:RCU Markers for Read-Side Critical Sections}
\end{table*}

\begin{table*}
\rowcolors{1}{}{lightgray}
\renewcommand*{\arraystretch}{1.15}
\footnotesize
\centering\OneColumnHSpace{-.4in}
\ebresizewidth{
\begin{tabular}{llp{2.0in}}
\toprule
获取 &
	释放 &
		说明 \\
\midrule
\tco{srcu_read_lock()} &
	\tco{srcu_read_unlock()} & \\
\tco{guard(srcu)(&my_srcu)} & & SRCU读者持续到外围作用域结束。 \\
\tco{scoped_guard(srcu, &my_srcu)} & & SRCU读者用于紧随其后的语句。 \\
\tco{srcu_read_lock_nmisafe()} &
	\tco{srcu_read_unlock_nmisafe()} & \\
\tco{srcu_read_lock_notrace()} &
	\tco{srcu_read_unlock_notrace()} & \\
\tco{srcu_down_read()} &
	\tco{srcu_up_read()} & \\
\bottomrule
\end{tabular}
}
\caption{SRCU Read-Side Critical Section标记}
\label{tab:defer:SRCU Markers for Read-Side Critical Sections}
\end{table*}

\begin{table*}
\rowcolors{1}{}{lightgray}
\renewcommand*{\arraystretch}{1.15}
\footnotesize
\centering\OneColumnHSpace{-.4in}
\ebresizewidth{
\begin{tabular}{p{2.5in}p{2.5in}}
\toprule
压缩状态 / 完整状态 &
		说明 \\
\midrule
\tco{get_completed_synchronize_rcu()}
\tco{get_completed_synchronize_rcu_full()} &
	返回始终已完成的状态。 \\
\tco{get_state_synchronize_rcu()}
\tco{get_state_synchronize_rcu_full()} &
	返回用于grace period完成检测的状态。 \\
\tco{start_poll_synchronize_rcu()}
\tco{start_poll_synchronize_rcu_full()} &
	返回用于grace period完成检测的状态，
	并在需要时启动grace period。 \\
\tco{start_poll_synchronize_rcu_expedited()}
\tco{start_poll_synchronize_rcu_expedited_full()} &
	返回用于grace period完成检测的状态，
	并在需要时启动加速的grace period。 \\
\tco{poll_state_synchronize_rcu()}
\tco{poll_state_synchronize_rcu_full()} &
	如果该状态的grace period已结束则返回true。 \\
\tco{cond_synchronize_rcu()}
\tco{cond_synchronize_rcu_full()} &
	（如需要）等待该状态的grace period结束。 \\
\tco{cond_synchronize_rcu_expedited()}
\tco{cond_synchronize_rcu_expedited_full()} &
	（如需要）以加速方式等待该状态的grace period结束。 \\
\tco{NUM_ACTIVE_RCU_POLL_OLDSTATE}
\tco{NUM_ACTIVE_RCU_POLL_OLDSTATE} &
	对应进行中grace period的最大状态数。 \\
\tco{same_state_synchronize_rcu()}
\tco{same_state_synchronize_rcu_full()} &
	两个状态是否等价？ \\
\bottomrule
\end{tabular}
}
\caption{RCU轮询更新侧原语}
\label{tab:defer:RCU Polled Update-Side Primitives}
\end{table*}

如果你是RCU新手，可以考虑只关注
\cref{tab:defer:RCU Wait-to-Finish APIs}
中的某一列，每一列概述了Linux内核RCU API系列的一个成员。
例如，如果你主要对理解RCU在Linux内核中的使用感兴趣，
``RCU''列将是起点，因为它使用得最频繁。
另一方面，如果你想为了理解RCU本身而学习，
``Tasks RCU Rude''拥有最简单的API。
你随时可以回来看其他列。

如果你已经熟悉RCU，这些表格可以
作为有用的参考。

\QuickQuiz{
	为什么\cref{tab:defer:RCU Wait-to-Finish APIs}
	中的一些单元格有感叹号（``!''）？
}\QuickQuizAnswer{
	带有感叹号的API成员（\co{rcu_read_lock()}、
	\co{rcu_read_unlock()}和\co{call_rcu()}）是
	Paul E. McKenney在90年代中期所知的
	Linux RCU API的仅有成员。
	在那个时期，他错误地认为自己
	已经了解了关于RCU的一切。
}\QuickQuizEnd

``RCU''列对应于三个Linux内核RCU
实现的合并~\cite{PaulEMcKenney2019RCUCVE,McKenney:2019:CRS:3319647.3325836}，
其中RCU read-side critical section以
\apik{rcu_read_lock()}开始、以\apik{rcu_read_unlock()}结束。
这些临界区也可以使用
\cref{tab:defer:RCU Markers for Read-Side Critical Sections}中的API，包括
\apik{rcu_read_lock_bh()}、\apik{rcu_read_lock_sched()}、
\apik{rcu_read_unlock_bh()}或\apik{rcu_read_unlock_sched()}。
此外，任何禁用下半部、中断或抢占的代码区域
也充当RCU read-side critical section。
RCU read-side critical section可以嵌套。\footnote{
	请注意\co{guard(rcu)()}和\co{scoped_guard(rcu)}
	提供的RAII变体。}
对应的同步更新侧原语
\apik{synchronize_rcu()}和\apik{synchronize_rcu_expedited()}，
以及它们的同义词\apik{synchronize_net()}，
等待任何类型的当前正在执行的RCU read-side critical section完成。
这种等待的时长被称为``\IX{grace period}''，
\apik{synchronize_rcu_expedited()}旨在以增加CPU开销和IPI
为代价减少\IXh{grace-period}{latency}。
异步更新侧原语\apik{call_rcu()}在
后续grace period之后调用指定的函数和指定的参数。
例如，\co{call_rcu(p,f);}将导致
``RCU callback'' \co{f(p)}在后续grace period之后被调用。
使用\co{CONFIG_RCU_LAZY=y}构建的内核会将
\co{call_rcu()}本来会立即启动的grace period延迟最多数秒。
在这种延迟不可接受的情况下，可以使用
\apik{call_rcu_hurry()}来立即启动grace period。
有些情况下，例如卸载使用\co{call_rcu()}或
\co{call_rcu_hurry()}的Linux内核模块时，
需要等待所有未完成的RCU callback完成~\cite{PaulEMcKenney2007rcubarrier}。
\apik{rcu_barrier()}原语完成此任务。

\QuickQuizSeries{%
\QuickQuizB{
	如何防止大量RCU read-side critical section
	无限期地阻塞\co{synchronize_rcu()}调用？
}\QuickQuizAnswerB{
	无需做任何事情来防止RCU read-side critical section
	无限期地阻塞\co{synchronize_rcu()}调用，
	因为\co{synchronize_rcu()}调用只需等待
	\emph{已有的}RCU read-side critical section。
	因此，只要每个RCU read-side critical section
	的持续时间是有限的，RCU grace period也将保持有限。
}\QuickQuizEndB
%
\QuickQuizM{
	\co{synchronize_rcu()} API会等待所有已有的
	中断处理程序完成，对吧？
}\QuickQuizAnswerM{
	在v4.20及更高版本的Linux内核中，
	是的~\cite{PaulEMcKenney2019RCUCVE,McKenney:2019:CRS:3319647.3325836}。

	但在更早的内核中不是这样，尤其是在使用
	可抢占RCU时！
	你应当改用\apik{synchronize_irq()}。
	或者，你可以在希望\co{synchronize_rcu()}等待的
	特定中断处理程序中放置\co{rcu_read_lock()}和
	\co{rcu_read_unlock()}调用。
	但即使如此也要小心，因为可抢占RCU不保证
	等待中断处理程序中\co{rcu_read_lock()}之前或
	\co{rcu_read_unlock()}之后的部分。
}\QuickQuizEndM
%
\QuickQuizM{
	既然任何禁用中断的代码区域都充当
	RCU read-side critical section，那么原子操作（atomic operation）和
	单条机器指令是否也充当微小的RCU read-side
	critical section？
}\QuickQuizAnswerM{
	是的，确实如此。

	但如果你想使用这个技巧，请三思，
	而且这次要更认真地思考。

	如果你确实别无选择只能使用此技巧，
	请仔细记录你对它的使用。
	否则，阅读代码的每个人（包括你自己）都会
	因此而恨你。

	为什么如此多的仇恨？

	要理解原因，请记住Jack Slingwine和我发明RCU的
	一大优势在于我们标记了RCU读者。
	听听所有那些独立发明了类似RCU东西但
	未能标记其read-side critical section的人的话：
	你确实需要标记你的RCU读者！
	包括你的单指令RCU读者。
}\QuickQuizEndM
%
\QuickQuizE{
	\co{synchronize_rcu()}和\co{rcu_barrier()}
	之间有什么区别？
}\QuickQuizAnswerE{
	它们等待的东西不同。
	\co{synchronize_rcu()}等待已有的RCU read-side
	critical section完成，而\co{rcu_barrier()}等待
	先前\co{call_rcu()}调用的callback被调用。

\begin{table}
\renewcommand*{\arraystretch}{1.2}
\centering
\small
\tcresizewidth{
\begin{fcvref}[ln:defer:synchonize-rcu vs rcu-barrier]
\begin{tabular}{lll}
\toprule
            & \multicolumn{2}{c}{必须等待直到（代码片段中的）：} \\
\cmidrule{2-3}
\multicolumn{1}{c}{调用于：} & \multicolumn{1}{c}{\tco{synchronize_rcu()}}
					& \multicolumn{1}{c}{\tco{rcu_barrier()}} \\
\cmidrule(r){1-1} \cmidrule{2-3}
\tco{do_something_1()} & 			 & \\
\tco{do_something_2()} & \tco{rcu_read_unlock()} (\clnref{rrul}) & \\
\tco{do_something_3()} & \tco{rcu_read_unlock()} (\clnref{rrul})
						 & \tco{f(&p->rh)} (\clnref{cb}) \\
\tco{do_something_4()} &			 & \tco{f(&p->rh)} (\clnref{cb}) \\
\tco{do_something_5()} & 			 & \\
\bottomrule
\end{tabular}
\end{fcvref}
}

\begin{fcvlabel}[ln:defer:synchonize-rcu vs rcu-barrier]
\begin{VerbatimN}[commandchars=\\\@\$]
	do_something_1();			\lnlbl@ds1$
	rcu_read_lock();			\lnlbl@rrl$
	do_something_2();			\lnlbl@ds2$
	call_rcu(&p->rh, f);			\lnlbl@cr$
	do_something_3();			\lnlbl@ds3$
	rcu_read_unlock();			\lnlbl@rrul$
	do_something_4();			\lnlbl@ds4$
	// f(&p->rh) invoked			\lnlbl@cb$
	do_something_5();			\lnlbl@ds5$
\end{VerbatimN}
\end{fcvlabel}
\caption{\tco{synchronize_rcu()} vs. \tco{rcu_barrier()}}
\label{tab:defer:synchonize-rcu vs rcu-barrier}
\end{table}

	\cref{tab:defer:synchonize-rcu vs rcu-barrier}
	说明了这一区别，其中底部的代码片段展示了
	给定CPU正在执行的代码。
	为简单起见，假设没有其他CPU在执行
	\co{rcu_read_lock()}、\co{rcu_read_unlock()}或
	\co{call_rcu()}。

	该表展示了如果与各\co{do_something_*()}
	函数并发调用，每个原语必须等待多久，
	空单元格表示无需等待。
	如你所见，\co{synchronize_rcu()}只有在
	RCU read-side critical section中才需要等待，
	在这种情况下它必须等待结束该临界区的
	\co{rcu_read_unlock()}。
	相反，RCU read-side critical section对
	\co{rcu_barrier()}没有影响。
	然而，当\co{rcu_barrier()}在\co{call_rcu()}调用之后
	执行时，它必须等到相应的RCU callback被调用。

	话虽如此，有一种特殊情况，每次
	\co{rcu_barrier()}调用都可以被直接的
	\co{synchronize_rcu()}调用替换，即
	\co{synchronize_rcu()}是用\co{call_rcu()}实现的，
	且只有一个全局callback列表的情况。
	但请勿在可移植代码中这样做！！！
}\QuickQuizEndE
}

\co{synchronize_rcu()}和\co{call_rcu()}系列API易于使用，
自动为用户提供排队、批处理和唤醒功能。
然而，有些情况下这种易用性反而成为障碍，
例如，当受RCU保护的对象缓存可能在本应允许
释放它们的grace period期间被重用时。
虽然这可以通过仔细使用与\co{call_rcu()}排队的
RCU callback交互的标志来处理，但这可能不便，
并且由于注定失败的\co{call_rcu()}调用的开销而浪费CPU时间。

在这些情况下，RCU的轮询grace period原语可以提供帮助。
\apik{get_state_synchronize_rcu()}和\apik{start_poll_synchronize_rcu()}
函数返回一个grace period状态"cookie"，它封装了
后续grace period完成后的grace period状态。
一旦这样的后续grace period结束，
\apik{poll_state_synchronize_rcu()}将返回\co{true}。
RCU的轮询grace period API中还有相当多的其他函数。
有关它们的更多信息，包括示例代码，请参见
2024年LSF/MM关于此主题的演示幻灯片第34页及
之后的内容~\cite{PaulEMcKenney2024ReclamationInteractionsWithRCU}。\footnote{
	Linux内核源代码中这些函数的kernel-doc头
	也可能非常有帮助。}

最后，RCU可用于提供
\IX{type-safe memory}~\cite{Cheriton96a}，
如\cref{sec:defer:Type-Safe Memory}所述。
在RCU的上下文中，类型安全内存保证给定的
数据元素在访问它的任何RCU read-side critical section
期间不会改变类型。
要使用基于RCU的类型安全内存，请将
\apik{SLAB_TYPESAFE_BY_RCU}传递给\apik{kmem_cache_create()}。

\cref{tab:defer:RCU Wait-to-Finish APIs}中的``SRCU''列
展示了一个专用的RCU API，允许在SRCU read-side
critical section~\cite{PaulEMcKenney2006c}中进行通用的睡眠，
这些临界区由\apik{srcu_read_lock()}和\apik{srcu_read_unlock()}界定。
然而，与RCU不同，SRCU的\apik{srcu_read_lock()}返回一个值，
必须将其传入对应的\apik{srcu_read_unlock()}。
这一区别是因为SRCU用户为每种不同的SRCU用法分配一个
\apik{srcu_struct}，因此没有方便的地方存储每任务的
读者嵌套计数。
（请记住，虽然Linux内核提供动态分配的
每CPU存储，但目前还没有动态分配的每任务存储。）
SRCU read-side critical section也可以用
\cref{tab:defer:SRCU Markers for Read-Side Critical Sections}中显示的API标记，
同样注意新的RAII变体。

给定的\co{srcu_struct}结构可以用\co{DEFINE_SRCU()}定义为
全局变量（如果该结构必须在多个翻译单元中使用），
或用\co{DEFINE_STATIC_SRCU()}定义。
例如，\co{DEFINE_SRCU(my_srcu)}将创建一个名为
\co{my_srcu}的全局变量，可被程序中的任何文件使用。
或者，\co{srcu_struct}结构可以是栈上变量
或动态分配的内存区域。
在这两种非全局变量的情况下，内存必须在首次使用前
用\co{init_srcu_struct()}初始化，
并在最后一次使用后（但在底层存储消失之前）
用\co{cleanup_srcu_struct()}清理。

无论如何创建，这些不同的\co{srcu_struct}结构
防止SRCU read-side critical section阻塞
不相关的\apik{synchronize_srcu()}和
\apik{synchronize_srcu_expedited()}调用。
当然，在SRCU read-side critical section内使用
\apik{synchronize_srcu()}或\apik{synchronize_srcu_expedited()}
可能导致自死锁，因此应当避免（v6.4及更高版本中
lockdep会检查这一点）。
与RCU一样，SRCU的\co{synchronize_srcu_expedited()}比
\co{synchronize_srcu()}减少了grace period延迟，
但代价是增加CPU开销。
同样与RCU类似，有\apik{get_state_synchronize_srcu()}、
\apik{start_poll_synchronize_srcu()}和
\apik{poll_state_synchronize_srcu()}函数允许轮询SRCU grace period API。

\QuickQuiz{
	在什么条件下可以在SRCU read-side critical section内
	安全地使用\co{synchronize_srcu()}？
}\QuickQuizAnswer{
	原则上，你可以在使用某个\co{srcu_struct}的
	SRCU read-side critical section内，
	对另一个\co{srcu_struct}使用\co{synchronize_srcu()}或
	\co{synchronize_srcu_expedited()}。
	然而，在实践中这样做几乎肯定是个坏主意。
	特别是，\cref{lst:defer:Multistage SRCU Deadlocks}
	中所示的代码仍然可能导致死锁，
	在v6.4及更高版本的Linux内核中lockdep会对此发出警告。

\begin{listing}
\begin{VerbatimL}
idx = srcu_read_lock(&ssa);
synchronize_srcu(&ssb);
srcu_read_unlock(&ssa, idx);

/* . . . */

idx = srcu_read_lock(&ssb);
synchronize_srcu(&ssa);
srcu_read_unlock(&ssb, idx);
\end{VerbatimL}
\caption{多阶段SRCU死锁}
\label{lst:defer:Multistage SRCU Deadlocks}
\end{listing}
}\QuickQuizEnd

与普通RCU类似，可以使用异步的\apik{call_srcu()}函数来
避免自死锁。
然而，使用\co{call_srcu()}时必须特别小心，
因为单个任务可能非常快速地注册SRCU callback。
鉴于SRCU读者可以阻塞任意长的时间，
这可能消耗任意大量的内存。
相比之下，使用同步的\co{synchronize_srcu()}
接口，给定任务必须完成等待当前grace period，
然后才能开始等待下一个。

同样与RCU类似，有一个\apik{srcu_barrier()}函数，
等待所有先前\co{call_srcu()} callback被调用。

换言之，SRCU通过允许开发者限制其作用域来
弥补其极弱的前向进展保证。

\cref{tab:defer:RCU Wait-to-Finish APIs}中的``Tasks RCU''列
展示了一个专用RCU API，用于协调Linux内核追踪中
使用的蹦床的释放。
这些蹦床用于将控制从被追踪代码中的某一点
转移到执行实际追踪的代码。
当然，必须确保在释放给定蹦床之前，
该蹦床中所有正在执行的代码都已完成。

对被追踪代码的更改通常仅限于单条跳转或调用指令，
因此无法容纳实现\co{rcu_read_lock()}和
\co{rcu_read_unlock()}所需的代码序列。
蹦床本身也不能包含这些\co{rcu_read_lock()}和
\co{rcu_read_unlock()}调用。
要理解这一点，考虑一个即将开始执行
给定蹦床的CPU。
因为它尚未执行\co{rcu_read_lock()}，该
蹦床可能随时被释放，这对该CPU来说将是致命的
意外。
因此，蹦床不能通过在被追踪代码或蹦床本身中
执行的同步原语来保护。
这就引出了蹦床究竟如何被保护的问题。

回答这个问题的关键在于注意到蹦床代码
从不包含直接或间接进行自愿上下文切换的代码。
此代码可能被抢占，但它永远不会直接或间接
调用\apik{schedule()}。
这提示了一种RCU变体，其中自愿上下文切换和
空闲执行是唯一的quiescent state。
此变体就是Tasks RCU。

Tasks RCU的不寻常之处在于没有读侧标记函数，
鉴于其主要用例无处放置此类标记，这倒是件好事。
相反，\co{schedule()}的调用直接充当quiescent state，
但仅限于那些对应自愿上下文切换的\co{schedule()}调用。
更新可以使用\apik{synchronize_rcu_tasks()}来等待所有
已有的蹦床执行完成，或者使用其异步对应物
\apik{call_rcu_tasks()}。
还有一个\apik{rcu_barrier_tasks()}，等待所有先前
\co{call_rcu_tasks()}调用对应的callback完成。
没有\co{synchronize_rcu_tasks_expedited()}，因为
尚未有此需求，不过实现一个有用的变体也并非
没有挑战。

\QuickQuiz{
	在使用\co{CONFIG_PREEMPT_NONE=y}构建的内核中，
	\co{synchronize_rcu()}不是会等待所有蹦床完成吗？
	毕竟抢占被禁用了，而蹦床从不直接或间接
	调用\co{schedule()}。
}\QuickQuizAnswer{
	你说得对！

	实际上，在不可抢占内核中，\co{synchronize_rcu_tasks()}
	是\co{synchronize_rcu()}的包装器。
	但请注意，具有运行时可设置抢占
	（\co{CONFIG_PREEMPT_DYNAMIC=y}）或启用了惰性抢占
	（\co{CONFIG_PREEMPT_LAZY=y}）的内核实际上是可抢占的。
}\QuickQuizEnd

``Tasks RCU Rude''列提供了
\cref{sec:defer:Toy Implementation}中展示的
玩具实现的更有效变体。
此变体使每个CPU执行一次上下文切换，
这样任何自愿上下文切换或任何可抢占的代码区域
都可以充当quiescent state。
Tasks RCU Rude变体使用Linux内核的workqueue设施来
强制并发上下文切换，与玩具实现采取的
串行逐CPU方法形成对比。
其API与Tasks RCU类似，包括缺少显式的读侧标记。

最后，``Tasks RCU Trace''列提供了一个
功能类似于SRCU的RCU实现，但具有快得多的
读侧标记。\footnote{
	因此在Tasks RCU家族中是不寻常的，因为它有
	显式的读侧标记！}
然而，这种速度是因为这些标记不执行内存屏障指令，
这意味着Tasks RCU Trace grace period通常必须向
所有CPU发送IPI，从而降低实时响应。
尽管如此，在没有读者的情况下，由此产生的
grace period延迟相当短，可与RCU相媲美。

\subsubsection{RCU拥有发布-订阅和版本维护API}
\label{sec:defer:RCU has Publish-Subscribe and Version-Maintenance APIs}

幸运的是，\cref{tab:defer:RCU Publish-Subscribe and Version Maintenance APIs}
中所示的RCU发布-订阅和版本维护原语
适用于上面讨论的所有RCU变体。
这种共性允许更多代码被共享，并减少API的增殖。
RCU发布-订阅API的最初目的是将内存屏障
嵌入这些API中，使Linux内核程序员可以使用RCU
而无需成为Linux支持的20多种CPU系列中每一种的
内存排序模型的专家~\cite{Spraul01}。

\begin{table*}
\renewcommand*{\arraystretch}{1.15}
\footnotesize
\centering\OneColumnHSpace{-.4in}
\ebresizewidth{
\begin{tabular}{llp{2.2in}}
\toprule
类别 &
	原语 &
		开销 \\
\midrule
指针发布 &
	\tco{rcu_assign_pointer()} &
		内存屏障 \\
&
	\tco{rcu_replace_pointer()} &
		内存屏障（Alpha上两个） \\
&
	\tco{rcu_pointer_handoff()} &
		简单指令 \\
&
	\tco{unrcu_pointer()} &
		简单指令 \\
&
	\tco{RCU_INIT_POINTER()} &
		简单指令 \\
&
	\tco{RCU_POINTER_INITIALIZER()} &
		编译时常量 \\
\midrule
指针订阅（遍历） &
	\tco{rcu_access_pointer()} &
		简单指令 \\
&
	\tco{rcu_dereference()} &
		简单指令（Alpha上为内存屏障） \\
&
	\tco{rcu_dereference_check()} &
		简单指令（Alpha上为内存屏障） \\
&
	\tco{rcu_dereference_bh()} &
		简单指令（Alpha上为内存屏障） \\
&
	\tco{rcu_dereference_bh_check()} &
		简单指令（Alpha上为内存屏障） \\
&
	\tco{rcu_dereference_sched()} &
		简单指令（Alpha上为内存屏障） \\
&
	\tco{rcu_dereference_sched_check()} &
		简单指令（Alpha上为内存屏障） \\
&
	\tco{rcu_dereference_protected()} &
		简单指令 \\
&
	\tco{rcu_dereference_raw()} &
		简单指令（Alpha上为内存屏障） \\
&
	\tco{rcu_dereference_raw_check()} &
		简单指令（Alpha上为内存屏障） \\
\midrule
SRCU指针订阅（遍历） &
	\tco{srcu_dereference()} &
		简单指令 \\
&
	\tco{srcu_dereference_check()} &
		简单指令 \\
&
	\tco{srcu_dereference_notrace()} &
		简单指令 \\
\bottomrule
\end{tabular}
}
\caption{RCU发布-订阅和版本维护API}
\label{tab:defer:RCU Publish-Subscribe and Version Maintenance APIs}
\end{table*}

这些原语直接操作指针，对于创建受RCU保护的
链式数据结构（如受RCU保护的数组和树）非常有用。
链表的特殊情况由
\cref{sec:defer:RCU has List-Processing APIs}
中描述的一组单独的API处理。

第一类发布指向新数据项的指针。
\apik{rcu_assign_pointer()}原语确保在弱序机器上，
任何先前的初始化保持排在指针赋值之前。
\apik{rcu_replace_pointer()}原语像\co{rcu_assign_pointer()}
一样更新指针，但也像\co{rcu_dereference_protected()}
（见下文）那样返回先前的值，包括lockdep表达式。
当更新者必须同时发布新指针并释放旧指针引用的结构时，
这种替换非常方便。

\QuickQuizSeries{%
\QuickQuizB{
	通常，任何受\co{rcu_dereference()}保护的指针
	\emph{必须}始终使用
	\cref{tab:defer:RCU Publish-Subscribe and Version Maintenance APIs}
	中的指针发布函数来更新，
	例如\co{rcu_assign_pointer()}。

	有没有例外？
}\QuickQuizAnswerB{
	一个例外是：当多元素链式数据结构
	在其他CPU不可访问时作为一个整体初始化，
	然后用单个\co{rcu_assign_pointer()}
	设置一个指向此数据结构的全局指针。
	初始化时的指针赋值不需要使用
	\co{rcu_assign_pointer()}，
	尽管在结构全局可见后发生的任何赋值
	\emph{必须}使用\co{rcu_assign_pointer()}。

	然而，除非此初始化代码处于非常热的代码路径上，
	即使理论上不必要，使用\co{rcu_assign_pointer()}
	可能也是明智的。
	一个"小"变更使你珍视的关于初始化
	私下进行的假设失效太容易了。
}\QuickQuizEndB
%
\QuickQuizE{
	这些遍历和更新原语可以与任何RCU API系列成员
	一起使用，这有什么缺点吗？
}\QuickQuizAnswerE{
	有时自动化代码检查器如``sparse''
	（或者实际上人类也是如此）
	很难确定给定的RCU遍历原语对应于
	哪种类型的RCU read-side critical section。
	例如，考虑\cref{lst:defer:Diverse RCU Read-Side Nesting}
	中所示的代码。

\begin{listing}
\begin{VerbatimL}
rcu_read_lock();
preempt_disable();
p = rcu_dereference(global_pointer);

/* . . . */

preempt_enable();
rcu_read_unlock();
\end{VerbatimL}
\caption{多样化的RCU读侧嵌套}
\label{lst:defer:Diverse RCU Read-Side Nesting}
\end{listing}

	\co{rcu_dereference()}原语是在普通的RCU
	临界区中还是在RCU Sched临界区中？
	你需要做什么才能弄清楚？

	但也许在v4.20 Linux内核中合并RCU风格之后，
	我们不再需要关心了！
}\QuickQuizEndE
}

\apik{rcu_pointer_handoff()}原语简单地返回其唯一参数，
但对于检查从RCU read-side critical section泄漏指针的
工具很有用。
使用\co{rcu_pointer_handoff()}向此类工具表明，
所讨论结构的保护已从RCU移交给某种其他机制，
如加锁或引用计数。

\apik{RCU_INIT_POINTER()}宏可用于初始化
尚未暴露给读者的受RCU保护的指针，
或者将受RCU保护的指针设为\co{NULL}。
在这些受限的情况下，不需要\co{rcu_assign_pointer()}
提供的内存屏障指令。
类似地，\apik{RCU_POINTER_INITIALIZER()}提供了一个
\GCC 风格的结构初始化器，以便于在结构中
初始化受RCU保护的指针。

第二类订阅指向数据项的指针，或者说
安全地遍历受RCU保护的指针。
同样，简单地使用C语言访问加载这些指针
可能导致看到被指向数据中的初始化前垃圾。
类似地，在RCU read-side critical section之外
以任何方式加载这些指针都可能导致被指向的对象
随时被释放。
然而，如果指针仅用于测试而不进行解引用，
被指向对象的释放不一定是问题。
在这种情况下，可以使用\apik{rcu_access_pointer()}。
但通常需要RCU读侧保护，因此
\apik{rcu_dereference()}原语使用Linux内核的\co{lockdep}
设施~\cite{JonathanCorbet2006lockdep}来验证此
\co{rcu_dereference()}调用受到
\co{rcu_read_lock()}、\co{srcu_read_lock()}或其他RCU读侧
标记的保护。
相比之下，\co{rcu_access_pointer()}原语不涉及
\co{lockdep}，因此在RCU read-side critical section
之外使用时不会引发\co{lockdep}警告。

另一种不需要保护的情况是更新侧代码
在持有更新侧锁的同时访问受RCU保护的指针。
\apik{rcu_dereference_protected()} API成员就是为此情况
提供的。其第一个参数是受RCU保护的指针，
第二个参数接受一个lockdep表达式，
描述必须持有哪些锁才能使访问安全。
从读者和更新者都调用的代码可以使用
\apik{rcu_dereference_check()}，它也接受lockdep表达式，
但也可以从不持有锁的读侧代码中调用。
在某些情况下，lockdep表达式可能非常复杂，
例如，使用细粒度加锁时，
大量锁中的任何一个都可能被持有，
要确定哪个适用可能相当困难。
在这些（希望罕见的）情况下，\apik{rcu_dereference_raw()}
提供保护但不检查是否在读者内调用或持有任何特定的锁。
\apik{rcu_dereference_raw_notrace()} API成员行为类似，
但不可被追踪，因此可以安全地被追踪代码使用。

虽然几乎任何链式结构都可以通过操作指针来访问，
但更高层的结构可以非常有帮助。
因此下一节查看Linux内核使用的各种
受RCU保护的链表。

\subsubsection{RCU拥有链表处理API}
\label{sec:defer:RCU has List-Processing APIs}

\begin{figure}
\centering
\resizebox{3in}{!}{\includegraphics{defer/Linux_list}}
\caption{Linux环形双向链表（\tco{list}）}
\label{fig:defer:Linux Circular Linked List (list)}
\end{figure}

\begin{figure}
\centering
\resizebox{3in}{!}{\includegraphics{defer/Linux_list_abbr}}
\caption{Linux链表简略图}
\label{fig:defer:Linux Linked List Abbreviated}
\end{figure}

虽然\co{rcu_assign_pointer()}和\co{rcu_dereference()}
理论上可以用于构建任何可以想象的受RCU保护的数据结构，
但在实践中使用更高层的构造通常更好。
因此，\co{rcu_assign_pointer()}和\co{rcu_dereference()}
原语已被嵌入Linux链表操作API的特殊RCU变体中。
Linux有四种双向链表变体：环形的
\co{struct list_head}和线性的
\co{struct hlist_head}/\co{struct hlist_node}、
\co{struct hlist_nulls_head}/\co{struct hlist_nulls_node}和
\co{struct hlist_bl_head}/\co{struct hlist_bl_node}
对。
前者的布局如
\cref{fig:defer:Linux Circular Linked List (list)}所示，
其中绿色（最左边）方框代表链表头，
蓝色（右边三个）方框代表链表中的元素。
这种表示法很繁琐，因此将如
\cref{fig:defer:Linux Linked List Abbreviated}所示简化，
只显示非头部的（蓝色）元素。

\begin{figure}
\centering
\resizebox{3in}{!}{\includegraphics{defer/Linux_hlist}}
\caption{Linux线性链表（\tco{hlist}）}
\label{fig:defer:Linux Linear Linked List (hlist)}
\end{figure}

Linux的\co{hlist}\footnote{
	``h''代表哈希表，与Linux的双指针环形链表
	相比，它将内存使用减半。}
是线性链表，这意味着头部只需一个指针，
而不是环形链表所需的两个，如
\cref{fig:defer:Linux Linear Linked List (hlist)}所示。
因此，使用\co{hlist}可以将大型哈希表的
哈希桶数组的内存消耗减半。
与前面一样，这种表示法很繁琐，
所以\co{hlist}结构将以与\co{list_head}风格链表
相同的方式简化，如
\cref{fig:defer:Linux Linked List Abbreviated}所示。

Linux的\co{hlist}的一个变体名为\co{hlist_nulls}，
提供多个不同的\co{NULL}指针，
但使用与\cref{fig:defer:Linux Linear Linked List (hlist)}
相同的布局。
在此变体中，低位为零的\co{->next}指针
被认为是一个指针。
然而，如果低位被设为一，则高位标识
\co{NULL}指针的类型。
此类型的链表用于允许无锁（lock-free）读者检测
节点何时从一个链表移动到另一个链表。
例如，哈希表的每个桶可以使用其索引来标记
其\co{NULL}指针。
如果读者遇到与其起始桶索引不匹配的\co{NULL}指针，
该读者就知道它正在遍历的元素在遍历期间
被移动到了其他桶，带着该读者一起。
读者可以使用\apik{is_a_nulls()}函数（当传入
\co{hlist_nulls}的\co{NULL}指针时返回\co{true}）
来确定何时到达链表末尾，
使用\apik{get_nulls_value()}函数
（返回其参数的\co{NULL}指针标识符）
来获取\co{NULL}指针的类型。
当\co{get_nulls_value()}返回意外值时，
读者可以采取纠正措施，例如从头重新开始遍历。

\QuickQuiz{
	但如果一个\co{hlist_nulls}读者被移到某个其他桶
	然后又移回来了呢？
}\QuickQuizAnswer{
	一种处理方式是始终将节点移动到目标桶的开头，
	确保当读者到达具有匹配\co{NULL}指针的链表末尾时，
	它已经搜索了整个链表。

	当然，如果在具有大量元素的哈希表中
	移动操作太多，读者可能永远无法到达链表末尾。
	在常见情况下避免这一问题的一种方法是
	保持哈希表调优良好，从而保持链表较短。
	一种检测问题并处理它的方法是让读者在
	遍历了某个大量节点后终止搜索，
	获取更新侧锁，然后重做搜索，
	但这可能引入死锁。
	另一种完全避免问题的方法是让读者
	在RCU read-side critical section内搜索，
	并在连续更新之间等待RCU grace period。
	一种折中方案是每$N$次更新等待一个RCU grace period，
	对于某个合适的$N$值。
}\QuickQuizEnd

关于\co{hlist_nulls}的更多信息可在Linux内核源代码树中
找到，\path{rculist_nulls.rst}文件
（旧内核中为\path{rculist_nulls.txt}）中提供了
有用的示例代码。

Linux的\co{hlist}的另一个变体结合了位加锁，
名为\co{hlist_bl}。
此变体使用与\cref{fig:defer:Linux Linear Linked List (hlist)}
相同的布局，但保留头指针（图中的``first''）的
低位来锁定链表。
这种方法也减少了内存使用，因为它允许
将本来需要单独的自旋锁存储与指针本身一起。

\begin{sidewaystable*}[htbp]
\rowcolors{1}{}{lightgray}
\renewcommand*{\arraystretch}{1.3}
\centering
\caption{受RCU保护的链表API}
\label{tab:defer:RCU-Protected List APIs}
\footnotesize
\newlength{\cwa}\newlength{\cwb}\newlength{\cwc}\newlength{\cwd}
\IfNimbusAvail{
  \renewcommand{\ttdefault}{NimbusMonoN}
  \setlength{\cwa}{1.9in}\setlength{\cwb}{2.1in}
  \setlength{\cwc}{1.8in}\setlength{\cwd}{1.6in}
}{
  \setlength{\cwa}{1.95in}\setlength{\cwb}{2.15in}
  \setlength{\cwc}{1.9in}\setlength{\cwd}{1.7in}
}
\ebresizewidthsw{
\begin{tabular}{>{\raggedright\arraybackslash}p{\cwa}
    >{\raggedright\arraybackslash}p{\cwb}
    >{\raggedright\arraybackslash}p{\cwc}
    >{\raggedright\arraybackslash}p{\cwd}}
\toprule
\pmb{\tco{list}}：环形双向链表 &
    \pmb{\tco{hlist}}：线性双向链表 &
	\pmb{\tco{hlist_nulls}}：带标记NULL指针的线性双向链表，最多31位标记 &
	    \pmb{\tco{hlist_bl}}：带位加锁的线性双向链表 \\
\midrule
\multicolumn{4}{l}{{\bf 结构}} \\
\tco{struct list_head} &
    \tco{struct}{\tt ~}\tco{hlist_head} ~~~~~~~~~~~~~~
    \tco{struct}{\tt ~}\tco{hlist_node} &
	\tco{struct}{\tt ~}\tco{hlist_nulls_head}
	\tco{struct}{\tt ~}\tco{hlist_nulls_node} &
	    \tco{struct}{\tt ~}\tco{hlist_bl_head}
	    \tco{struct}{\tt ~}\tco{hlist_bl_node} \\
\multicolumn{4}{l}{{\bf 初始化}} \\
\tco{INIT_LIST_HEAD_RCU()} &
    &
	&
	    \\
\multicolumn{4}{l}{{\bf 完整遍历}} \\
\tco{list_for_each_entry_rcu()}
\tco{list_for_each_entry_lockless()}
\tco{list_for_each_entry_srcu()} &
    \tco{hlist_for_each_entry_rcu()}
    \tco{hlist_for_each_entry_rcu_bh()}
    \tco{hlist_for_each_entry_rcu_notrace()}
    \tco{hlist_for_each_entry_srcu()} &
	\tco{hlist_nulls_for_each_entry_rcu()}
	\tco{hlist_nulls_for_each_entry_safe()} &
	    \tco{hlist_bl_for_each_entry_rcu()} \\
\multicolumn{4}{l}{{\bf 恢复遍历}} \\
\tco{list_for_each_entry_continue_rcu()}
\tco{list_for_each_entry_from_rcu()} &
    \tco{hlist_for_each_entry_continue_rcu()}
    \tco{hlist_for_each_entry_continue_rcu_bh()}
    \tco{hlist_for_each_entry_from_rcu()} &
	&
	    \\
\multicolumn{4}{l}{{\bf 逐步遍历}} \\
\tco{list_entry_rcu()}
\tco{list_entry_lockless()}
\tco{list_first_or_null_rcu()}
\tco{list_next_rcu()}
\tco{list_next_or_null_rcu()}
\tco{list_tail_rcu()} &
    \multicolumn{1}{p{1.2in}}{\tco{hlist_first_rcu()}
			      \tco{hlist_next_rcu()}
			      \tco{hlist_pprev_rcu()}} &
	\tco{hlist_nulls_first_rcu()}
	\tco{hlist_nulls_next_rcu()} &
	    \tco{hlist_bl_first_rcu()} \\
\multicolumn{4}{l}{{\bf 添加}} \\
\multicolumn{1}{p{1.2in}}{\tco{list_add_rcu()}
			  \tco{list_add_tail_rcu()}} &
    \tco{hlist_add_before_rcu()}
    \tco{hlist_add_behind_rcu()}
    \tco{hlist_add_head_rcu()}
    \tco{hlist_add_tail_rcu()} &
	\tco{hlist_nulls_add_head_rcu()}
	\tco{hlist_nulls_add_tail_rcu()}
	\tco{hlist_nulls_add_fake()} &
	    \tco{hlist_bl_add_head_rcu()}
	    \tco{hlist_bl_set_first_rcu()} \\
\multicolumn{4}{l}{{\bf 删除}} \\
\tco{list_del_rcu()} &
    \multicolumn{1}{p{1.2in}}{\tco{hlist_del_rcu()}
			      \tco{hlist_del_init_rcu()}} &
	\tco{hlist_nulls_del_rcu()}
	\tco{hlist_nulls_del_init_rcu()} &
	    \tco{hlist_bl_del_rcu()} \\
\multicolumn{4}{l}{{\bf 替换}} \\
\tco{list_replace_rcu()} &
    \tco{hlist_replace_rcu()} &
	&
	    \\
\multicolumn{4}{l}{{\bf 拼接}} \\
\tco{list_splice_init_rcu()}
\tco{list_splice_tail_init_rcu()} &
    \tco{hlists_swap_heads_rcu()} &
	&
	    \\
\bottomrule
\end{tabular}
}
\end{sidewaystable*}

这些链表变体的API成员在
\cref{tab:defer:RCU-Protected List APIs}中汇总。
更多信息可在Linux内核源代码树的\path{Documentation/RCU}
目录和Linux Weekly News~\cite{PaulEMcKenney2024RCUAPI}中找到。

然而，本节的其余部分将展开讨论
\apik{list_replace_rcu()}的使用，因为正是这个API成员
赋予了RCU其名称。
此API成员用于执行更复杂的更新，其中链表中间的
一个具有多个字段的元素被原子地更新，
使得给定的读者看到的要么是旧的值集合，
要么是新的值集合，而不是两个集合的混合。
例如，链表的每个节点可能有整数字段
\co{->a}、\co{->b}和\co{->c}，
可能需要将给定节点的字段从5、6、7分别
更新为5、2、3。

实现此原子更新的代码非常直接：

\begin{fcvlabel}[ln:defer:Canonical RCU Replacement Example (2nd)]
\begin{VerbatimN}[samepage=true,commandchars=\\\[\],firstnumber=15]
q = kmalloc(sizeof(*p), GFP_KERNEL);	\lnlbl[kmalloc]
*q = *p;				\lnlbl[copy]
q->b = 2;				\lnlbl[update1]
q->c = 3;				\lnlbl[update2]
list_replace_rcu(&p->list, &q->list);	\lnlbl[replace]
synchronize_rcu();			\lnlbl[sync_rcu]
kfree(p);				\lnlbl[kfree]
\end{VerbatimN}
\end{fcvlabel}

\begin{figure}
\centering
\IfEbookSize{
\resizebox{2in}{!}{\includegraphics{defer/RCUReplacement}}
}{
\resizebox{2.7in}{!}{\includegraphics{defer/RCUReplacement}}
}
\caption{链表中的RCU替换}
\label{fig:defer:RCU Replacement in Linked List}
\end{figure}

以下讨论逐步分析此代码，使用
\cref{fig:defer:RCU Replacement in Linked List}来说明
状态变化。
每个元素中的三元组分别代表字段\co{->a}、
\co{->b}和\co{->c}的值。
红色底纹的元素可能被读者引用，
由于读者不直接与更新者同步，
读者可能在整个替换过程中并发运行。
请注意，为清晰起见，反向指针和从尾部到头部的
链接被省略了。

链表的初始状态，包括指针\co{p}，
与删除示例相同，如图的第一行所示。

以下文字描述如何以使给定读者
看到这两个值之一的方式，
将\co{5,6,7}元素替换为\co{5,2,3}。

\begin{fcvref}[ln:defer:Canonical RCU Replacement Example (2nd)]
\Clnref{kmalloc}分配一个替换元素，
导致\cref{fig:defer:RCU Replacement in Linked List}
第二行所示的状态。
此时，没有读者可以持有对新分配元素的引用
（如其绿色底纹所示），且该元素未初始化
（如问号所示）。

\Clnref{copy}将旧元素复制到新元素，导致
\cref{fig:defer:RCU Replacement in Linked List}
第三行所示的状态。
新分配的元素仍然不能被读者引用，但现在已初始化。

\Clnref{update1}将\co{q->b}更新为值``2''，
\clnref{update2}将\co{q->c}更新为值``3''，
如\cref{fig:defer:RCU Replacement in Linked List}
第四行所示。
注意新分配的结构仍然对读者不可访问。

现在，\clnref{replace}执行替换，使新元素
终于对读者可见，因此显示为红色底纹，如
\cref{fig:defer:RCU Replacement in Linked List}
第五行所示。
此时，如下所示，我们有两个版本的链表。
已有的读者可能看到\co{5,6,7}元素（因此
现在显示为黄色底纹），但新读者将看到\co{5,2,3}元素。
但任何给定的读者都保证看到一组值或另一组值，
而不是两组的混合。

在\clnref{sync_rcu}上的\co{synchronize_rcu()}返回后，
一个grace period将已过去，因此所有在
\apik{list_replace_rcu()}之前开始的读取都将已完成。
特别是，任何可能持有对\co{5,6,7}元素引用的读者
都保证已退出其RCU read-side critical section，
因此被禁止继续持有引用。
因此，不再有任何读者持有对旧元素的引用，
如\cref{fig:defer:RCU Replacement in Linked List}
第六行中其绿色底纹所示。
就读者而言，我们回到了链表的单一版本，
但新元素替换了旧元素。

在\clnref{kfree}上的\apik{kfree()}完成后，
链表将如\cref{fig:defer:RCU Replacement in Linked List}
最后一行所示。
\end{fcvref}

尽管RCU以替换用例命名，
Linux内核中绝大多数RCU使用依赖于
简单的独立插入和删除，如
\cref{sec:defer:Maintain Multiple Versions of Recently Updated Objects}中
\cref{fig:defer:Multiple RCU Data-Structure Versions}所示。

\QuickQuiz{
	为什么\cref{tab:defer:RCU-Protected List APIs}
	不包含额外的链表初始化API，如\co{INIT_HLIST_HEAD()}？
}\QuickQuizAnswer{
	因为此API和类似的API不是RCU特有的，
	而是也（有时主要是）用于由其他同步机制保护的
	类似链表。
}\QuickQuizEnd

下一节将介绍帮助开发者调试
其使用RCU的代码的API。

\subsubsection{RCU拥有诊断API}
\label{sec:defer:RCU Has Diagnostic APIs}

\begin{table}
\renewcommand*{\arraystretch}{1.15}
\footnotesize
\centering
\begin{tabular}{ll}
\toprule
类别 &
	原语 \\
\midrule
标记RCU指针 &
	\tco{__rcu} \\
\midrule
调试对象支持 &
	\tco{init_rcu_head()} \\
&	\tco{destroy_rcu_head()} \\
&	\tco{init_rcu_head_on_stack()} \\
&	\tco{destroy_rcu_head_on_stack()} \\
\midrule
停滞警告控制 &
	\tco{rcu_cpu_stall_reset()} \\
\midrule
Callback检查 &
	\tco{rcu_head_init()} \\
&	\tco{rcu_head_after_call_rcu()} \\
\midrule
lockdep支持 &
	\tco{rcu_read_lock_held()} \\
&	\tco{rcu_read_lock_bh_held()} \\
&	\tco{rcu_read_lock_sched_held()} \\
&	\tco{srcu_read_lock_held()} \\
&	\tco{rcu_is_watching()} \\
&	\tco{RCU_LOCKDEP_WARN()} \\
&	\tco{rcu_sleep_check()} \\
\bottomrule
\end{tabular}
\caption{RCU诊断API}
\label{tab:defer:RCU Diagnostic APIs}
\end{table}

\Cref{tab:defer:RCU Diagnostic APIs}
展示了RCU的诊断API。

\co{__rcu}标签标记受RCU保护的指针，例如
\qtco{struct foo __rcu *p;}。
可能被传递给\co{rcu_dereference()}的指针可以被标记，
但持有\co{rcu_dereference()}返回值的指针不应被标记。
在变量、结构字段、函数参数和返回值上
提供这些标记，允许Linux内核的\co{sparse}工具
检测受RCU保护的指针被
纯C语言加载和存储错误访问的情况。

调试对象支持对于作为从Linux内核内存分配器获得的
结构一部分的任何\apik{rcu_head}结构是自动的，
但构建自己的专用内存分配器的人
可以分别在分配和释放时使用\apik{init_rcu_head()}和
\apik{destroy_rcu_head()}。
使用在函数调用栈上分配的\co{rcu_head}结构的人
（确实有这种情况！）可以在首次使用前使用
\apik{init_rcu_head_on_stack()}，在最后使用后、
函数返回前使用\apik{destroy_rcu_head_on_stack()}。
调试对象支持允许检测将同一\co{rcu_head}结构
快速连续传递给\co{call_rcu()}及其相关函数的错误，
这是\co{call_rcu()}对应于著名的双重释放类
内存分配错误的等价物。

停滞警告控制由\apik{rcu_cpu_stall_reset()}提供，
它允许调用者在当前grace period的剩余时间内
抑制RCU CPU停滞警告。
RCU CPU停滞警告有助于定位
RCU read-side critical section运行时间过长的情况，
内核调试器等可以抑制它们是有用的，
例如在遇到断点时。

Callback检查由\apik{rcu_head_init()}和
\apik{rcu_head_after_call_rcu()}提供。
前者在将\co{rcu_head}结构传递给\co{call_rcu()}之前
对其调用，然后\co{rcu_head_after_call_rcu()}将
检查callback是否已用指定函数调用。

对lockdep~\cite{JonathanCorbet2006lockdep}的支持包括
\apik{rcu_read_lock_held()}、
\apik{rcu_read_lock_bh_held()}、
\apik{rcu_read_lock_sched_held()}和
\apik{srcu_read_lock_held()}，
每个在相应类型的RCU read-side critical section内
调用时返回\co{true}。

\QuickQuiz{
	为什么没有用于Tasks RCU的\co{rcu_read_lock_tasks_held()}？
}\QuickQuizAnswer{
	因为Tasks RCU没有读侧标记。
	相反，Tasks RCU的read-side critical section
	由自愿上下文切换界定。
}\QuickQuizEnd

因为\co{rcu_read_lock()}不能在空闲循环中使用，
且能效考虑已导致空闲循环变得相当复杂，
\apik{rcu_is_watching()}在可以合法使用\co{rcu_read_lock()}
的上下文中返回\co{true}。
再次注意\co{srcu_read_lock()}可以在空闲甚至
离线的CPU上使用，这意味着\co{rcu_is_watching()}
不适用于SRCU。

\apik{RCU_LOCKDEP_WARN()}在lockdep启用且
其参数求值为\co{true}时发出警告。
例如，\co{RCU_LOCKDEP_WARN(!rcu_read_lock_held())}
在RCU read-side critical section之外调用时会发出警告。

\apik{RCU_NONIDLE()}可用于强制RCU在执行
作为唯一参数传入的语句时保持监视。
例如，\co{RCU_NONIDLE(WARN_ON(!rcu_is_watching()))}
永远不会发出警告。
然而，2020--2021年期间的变更将RCU的覆盖范围
扩展到空闲循环的更深处，这应该大大减少甚至
消除对\co{RCU_NONIDLE()}的需要。

最后，\apik{rcu_sleep_check()}在RCU、RCU-bh或
RCU-sched read-side critical section内调用时发出警告。

\subsubsection{RCU的API可以在哪里使用？}
\label{sec:defer:Where Can RCU's APIs Be Used?}

\begin{figure}
\centering
\resizebox{3in}{!}{\includegraphics{defer/RCUenvAPI}}
\caption{RCU API使用约束}
\label{fig:defer:RCU API Usage Constraints}
\end{figure}

\Cref{fig:defer:RCU API Usage Constraints}
展示了哪些API可以在哪些内核内环境中使用。
RCU读侧原语可以在任何环境中使用，包括NMI；
RCU变更和异步grace period原语可以在
NMI之外的任何环境中使用；
最后，RCU同步grace period原语只能在进程上下文中使用。
RCU链表遍历原语包括\apik{list_for_each_entry_rcu()}、
\apik{hlist_for_each_entry_rcu()}等。
类似地，RCU链表变更原语包括
\apik{list_add_rcu()}、\apik{hlist_del_rcu()}等。

请注意，作为一般规则，其他RCU系列的原语可以替代使用，
不过\co{srcu_read_lock_nmisafe()}而非
\co{srcu_read_lock()}可以在NMI上下文中使用。

\subsubsection{那么，RCU\emph{到底}是什么？}
\label{sec:defer:So; What is RCU Really?}

在其核心，RCU不多不少就是一个支持
插入的发布和订阅、等待所有RCU读者完成、
以及维护多个版本的API。
话虽如此，可以在RCU之上构建更高层的构造，
包括\cref{sec:defer:RCU Usage}中列出的
读写锁、引用计数和存在性保证构造。
此外，我毫不怀疑Linux社区将继续
发现RCU的有趣新用途，
正如他们对内核中许多同步原语所做的那样。

当然，对RCU更完整的理解还应包括
这些API可以做的所有事情。

然而，对于许多人来说，对RCU的完整理解
必须包括RCU实现的示例。
\Cref{chp:app:``Toy'' RCU Implementations}因此展示了一系列
复杂度和能力递增的"玩具"RCU实现，
不过有些人可能更喜欢经典的
``User-Level Implementations of Read-Copy
Update''~\cite{MathieuDesnoyers2012URCU}。
对于其他所有人，下一节给出了一些RCU用例的概述。

% defer/rcuusage.tex
% mainfile: ../perfbook.tex
% SPDX-License-Identifier: CC-BY-SA-3.0

\subsection{RCU的使用}
\label{sec:defer:RCU Usage}
\OriginallyPublished{sec:defer:RCU Usage}{RCU Usage}{Linux Weekly News}{PaulEMcKenney2008WhatIsRCUUsage}

本节从RCU可被应用的场景这一角度来回答``什么是RCU？''这个问题。
由于RCU最常被用来替换某些既有机制，
我们主要从它与这些机制的关系来审视它，
如\cref{tab:defer:RCU Usage}所列
以及\cref{fig:defer:Relationships Between RCU Use Cases}所展示的那样。
在本表所列各节之后，
\cref{sec:defer:RCU Usage Summary}给出了总结，
并在后续各节和其他文献~\cite{PaulEMcKenney2021RCUMysteriesFundamental,PaulEMcKenney2022RCUMysteriesUseCases}中进一步展开。

\begin{table}
\renewcommand*{\arraystretch}{1.2}
\centering
\small
\begin{tabular}{ll}
\toprule
RCU所替换的机制 & 页码 \\
\midrule
用于pre-BSD路由的RCU &
	\pageref{sec:defer:RCU for Pre-BSD Routing} \\
等待已有事务完成 &
	\pageref{sec:defer:Wait for Pre-Existing Things to Finish} \\
分阶段状态变更 &
	\pageref{sec:defer:Phased State Change} \\
仅添加链表（发布/订阅） &
	\pageref{sec:defer:Add-Only List} \\
类型安全内存 &
	\pageref{sec:defer:Type-Safe Memory} \\
存在性保证 &
	\pageref{sec:defer:Existence Guarantee} \\
轻量级垃圾收集器 &
	\pageref{sec:defer:Light-Weight Garbage Collector} \\
仅删除链表 &
	\pageref{sec:defer:Delete-Only List} \\
准读写锁 &
	\pageref{sec:defer:Quasi Reader-Writer Lock} \\
准引用计数 &
	\pageref{sec:defer:Quasi Reference Count} \\
准多版本并发控制（MVCC） &
	\pageref{sec:defer:Quasi Multi-Version Concurrency Control} \\
\bottomrule
\end{tabular}
\caption{RCU的使用}
\label{tab:defer:RCU Usage}
\end{table}

\subsubsection{用于Pre-BSD路由的RCU}
\label{sec:defer:RCU for Pre-BSD Routing}

与后面各节不同，本节聚焦于一个非常具体的用例，以便与其他机制进行比较。

\Cref{lst:defer:RCU Pre-BSD Routing Table Lookup,%
lst:defer:RCU Pre-BSD Routing Table Add/Delete}
展示了受RCU保护的Pre-BSD路由表代码
（\path{route_rcu.c}）。
前者展示了数据结构和\co{route_lookup()}，
后者展示了\co{route_add()}和\co{route_del()}。

\begin{listing}
\input{CodeSamples/defer/route_rcu@lookup.fcv}
\caption{RCU Pre-BSD路由表查找}
\label{lst:defer:RCU Pre-BSD Routing Table Lookup}
\end{listing}

\begin{listing}
\input{CodeSamples/defer/route_rcu@add_del.fcv}
\caption{RCU Pre-BSD路由表添加/删除}
\label{lst:defer:RCU Pre-BSD Routing Table Add/Delete}
\end{listing}

\begin{fcvref}[ln:defer:route_rcu:lookup]
在\cref{lst:defer:RCU Pre-BSD Routing Table Lookup}中，
\clnref{rh}添加了RCU回收所用的\co{->rh}字段，
\clnref{re_freed}添加了释放后使用检查字段\co{->re_freed}，
\clnref{lock,unlock1,unlock2}
添加了RCU读侧保护，
而\clnref{chk_freed,abort}则添加了释放后使用检查。
\end{fcvref}
\begin{fcvref}[ln:defer:route_rcu:add_del]
在\cref{lst:defer:RCU Pre-BSD Routing Table Add/Delete}中，
\clnref{add:lock,add:unlock,del:lock,%
del:unlock1,del:unlock2}添加了更新侧加锁，
\clnref{add:add_rcu,del:del_rcu}添加了RCU更新侧保护，
\clnref{del:call_rcu}使得\co{route_cb()}在
grace period过后被调用，
而\clnrefrange{cb:b}{cb:e}则定义了\co{route_cb()}。
对于一个可工作的并发实现来说，这已是最少量的附加代码。
\end{fcvref}

\begin{figure}
\centering
\resizebox{2.5in}{!}{\includegraphics{CodeSamples/defer/data/hps.2019.12.17a/perf-rcu}}
\caption{受RCU保护的Pre-BSD路由表}
\label{fig:defer:Pre-BSD Routing Table Protected by RCU}
\end{figure}

\Cref{fig:defer:Pre-BSD Routing Table Protected by RCU}
展示了在只读工作负载下的性能。
RCU的可扩展性相当好，提供了近乎理想的性能。
然而，这些数据是使用用户空间
RCU~\cite{MathieuDesnoyers2009URCU,PaulMcKenney2013LWNURCU}的\co{RCU_SIGNAL}
风格生成的，
其中\co{rcu_read_lock()}和\co{rcu_read_unlock()}
会生成少量代码。
对于\IXacr{qsbr}风格的RCU——\co{rcu_read_lock()}和\co{rcu_read_unlock()}
完全不生成任何代码——会发生什么呢？
（参见\cref{sec:defer:Introduction to RCU}，
特别是
\cref{fig:defer:QSBR: Waiting for Pre-Existing Readers}，
其中讨论了RCU QSBR。）

答案如
\cref{fig:defer:Pre-BSD Routing Table Protected by RCU QSBR}所示，
这表明RCU QSBR的性能和可扩展性实际上超过了
理想的无同步工作负载。

\QuickQuizSeries{%
\QuickQuizB{
	等等，什么？？？
	RCU QSBR怎么可能比理想情况还好？
	究竟是什么荒唐的"理想"定义才会不是
	所有可能结果中最好的？？？
}\QuickQuizAnswerB{
	这是一个极好的问题，答案是现代CPU和编译器极其复杂。
	但在深入探讨之前，非常值得注意的是，
	RCU QSBR的性能优势仅出现在
	每核一个硬件线程的情况下。
	一旦系统满载，RCU QSBR的性能就会
	回落到理想水平。

	\co{route_lookup()}搜索循环的RCU变体实际上
	比顺序版本多出一条x86指令，
	即序列\co{cmp}、\co{je}、\co{mov}、\co{cmp}、\co{lea}和\co{jne}中的\co{lea}。
	这条额外的指令是由于RCU变体的\co{route_entry}结构
	开头的\co{rcu_head}结构所致，
	因此，与顺序变体不同，RCU变体的
	\co{->re_next.next}指针有一个非零偏移量。
	回到20世纪80年代，这条额外的\co{lea}指令
	可能会可靠地导致RCU变体更慢，但我们
	现在已经是21\textsuperscript{世纪}了，80年代
	早已过去。

	但那些仔细阅读了
	\cref{sec:cpu:Pipelined CPUs}
	的读者早已知道这一切！

	这些反直觉的结果当然意味着，在现代微处理器上
	任何性能结果都必须持一定的怀疑态度。
	从理论上讲，获得优于理想的性能结果确实没有道理，
	但在现代微处理器上这确实可能发生。
	这样的结果可以被看作类似于著名的
	超线性加速（参见
	\cref{sec:SMPdesign:Beyond Partitioning}
	中的一个例子），即有趣但实际重要性有限。
	尽管如此，RCU的一大优势就是其读侧
	开销如此之低，以至于像这样的微小效应
	在实际的性能测量中清晰可见。

\begin{figure}
\centering
% Run with initial rcu_head structure in route_entry moved down.
\resizebox{2.5in}{!}{\includegraphics{CodeSamples/defer/data/hps.2019.12.17a/perf-rcu-qsbr}}
\caption{受RCU QSBR保护、\tco{rcu_head}非首字段的Pre-BSD路由表}
\label{fig:defer:Pre-BSD Routing Table Protected by RCU QSBR With Non-Initial rcu-head}
\end{figure}

	这就引出了一个问题：如果将
	\co{rcu_head}结构移动使得RCU的
	\co{->re_next.next}指针也具有零偏移量，
	与顺序变体完全相同，会发生什么。
	答案如
	\cref{fig:defer:Pre-BSD Routing Table Protected by RCU QSBR With Non-Initial rcu-head}
	所示，这会导致RCU QSBR的性能下降到
	仍然非常接近理想的水平，但不再是超理想的。
}\QuickQuizEndB
%
\QuickQuizE{
	鉴于RCU QSBR的读侧性能如此出色，为何还要费心使用
	用户空间RCU的其他风格呢？
}\QuickQuizAnswerE{
	因为RCU QSBR对整个应用程序施加了一些
	可能无法容忍的约束，
	例如，要求应用程序中的每一个线程
	都定期通过一个quiescent state。
	除此之外，这意味着RCU QSBR对
	库的编写者没有帮助，他们可能更适合使用其他
	风格的用户空间RCU~\cite{PaulMcKenney2013LWNURCU}。
}\QuickQuizEndE
}

\begin{figure}
\centering
% Run with initial rcu_head structure at beginning of route_entry structure.
% This need to run twice is the reason for the oddball directory.
% @@@ Make this generate both files with a single run?
\resizebox{2.5in}{!}{\includegraphics{CodeSamples/defer/data/hps.2019.12.02a/perf-rcu-qsbr}}
\caption{受RCU QSBR保护的Pre-BSD路由表}
\label{fig:defer:Pre-BSD Routing Table Protected by RCU QSBR}
\end{figure}

\begin{figure*}
\centering
\IfTwoColumn{
  \resizebox{5.5in}{!}{\includegraphics{defer/RCUusecases}}
}{
  \resizebox{.96\textwidth}{!}{\includegraphics{defer/RCUusecases}}
  % eb builds require .96
}
\caption{RCU各用例之间的关系}
\label{fig:defer:Relationships Between RCU Use Cases}
\end{figure*}

虽然Pre-BSD路由是一个极好的RCU用例，但还是值得
看看
\cref{fig:defer:Relationships Between RCU Use Cases}
所示的更广泛的用例之间的关系。
以下各节将承担这一任务。

在阅读这些节的同时，请思考哪个用例
最能描述Pre-BSD路由。

\subsubsection{等待已有事务完成}
\label{sec:defer:Wait for Pre-Existing Things to Finish}

如\cref{sec:defer:RCU Fundamentals}所述，
RCU的一个重要组成部分
是等待RCU读者完成的方式。
RCU的一大优势是它允许你等待数千种不同的事务完成，
而无需显式地跟踪每一个，也不会招致
使用显式跟踪的方案所固有的
性能下降、可扩展性限制、复杂的死锁
场景和内存泄漏风险。

在本节中，我们将展示\co{synchronize_sched()}的
读侧对应物（包括任何禁用抢占的操作，
以及硬件操作和
禁用中断的原语）如何允许你与
不可屏蔽中断（NMI）处理程序交互，这用锁机制来做是相当困难的。
这种方法被称为``纯RCU''~\cite{PaulEdwardMcKenneyPhD}，
在Linux内核中的少数地方有使用。

这种``纯RCU''设计的基本形式如下：

\begin{enumerate}
\item	进行一项变更，例如改变操作系统对NMI的响应方式。
\item	等待所有已有的read-side critical section
	完全结束（例如，使用
	\co{synchronize_sched()}原语）。\footnote{
		在Linux内核v5.1及更高版本中，\co{synchronize_sched()}已
		被纳入\co{synchronize_rcu()}。}
	这里的关键观察是，后续的RCU read-side critical section
	保证能看到所做的任何变更。
\item	清理，例如返回表示变更已成功完成的状态。
\end{enumerate}

本节的其余部分展示了改编自Linux内核的示例代码。
在此示例中，现已废弃的oprofile设施中的
\co{nmi_stop()}函数使用\co{synchronize_sched()}来确保
所有进行中的NMI通知在释放相关
资源之前已经完成。
此代码的简化版本如
\cref{lst:defer:Using RCU to Wait for NMIs to Finish}所示。

\begin{listing}
\begin{fcvlabel}[ln:defer:Using RCU to Wait for NMIs to Finish]
\begin{VerbatimL}[commandchars=\\\@\$]
struct profile_buffer {				\lnlbl@struct:b$
	long size;
	atomic_t entry[0];
};						\lnlbl@struct:e$
static struct profile_buffer *buf = NULL;	\lnlbl@struct:buf$

void nmi_profile(unsigned long pcvalue)		\lnlbl@nmi_profile:b$
{
	struct profile_buffer *p = rcu_dereference(buf);\lnlbl@nmi_profile:rcu_deref$

	if (p == NULL)				\lnlbl@nmi_profile:if_NULL$
		return;				\lnlbl@nmi_profile:ret:a$
	if (pcvalue >= p->size)			\lnlbl@nmi_profile:if_oor$
		return;				\lnlbl@nmi_profile:ret:b$
	atomic_inc(&p->entry[pcvalue]);		\lnlbl@nmi_profile:inc$
}						\lnlbl@nmi_profile:e$

void nmi_stop(void)				\lnlbl@nmi_stop:b$
{
	struct profile_buffer *p = buf;		\lnlbl@nmi_stop:fetch$

	if (p == NULL)				\lnlbl@nmi_stop:if_NULL$
		return;				\lnlbl@nmi_stop:ret$
	rcu_assign_pointer(buf, NULL);		\lnlbl@nmi_stop:NULL$
	synchronize_sched();			\lnlbl@nmi_stop:sync_sched$
	kfree(p);				\lnlbl@nmi_stop:kfree$
}						\lnlbl@nmi_stop:e$
\end{VerbatimL}
\end{fcvlabel}
\caption{使用RCU等待NMI完成}
\label{lst:defer:Using RCU to Wait for NMIs to Finish}
\end{listing}

\begin{fcvref}[ln:defer:Using RCU to Wait for NMIs to Finish:struct]
\Clnrefrange{b}{e}定义了一个\co{profile_buffer}结构，包含
一个大小字段和一个不定长的条目数组。
\Clnref{buf}定义了一个指向profile缓冲区的指针，
该指针大概在其他地方被初始化为指向一块动态分配的
内存区域。
\end{fcvref}

\begin{fcvref}[ln:defer:Using RCU to Wait for NMIs to Finish:nmi_profile]
\Clnrefrange{b}{e}定义了\co{nmi_profile()}函数，
它在NMI处理程序内被调用。
因此，它不能被抢占，也不能被普通的
中断处理程序中断，然而，它仍然可能由于缓存未命中、
ECC错误以及同一核心内其他硬件线程的周期窃取而延迟。
\Clnref{rcu_deref}使用\co{rcu_dereference()}原语获取
指向profile缓冲区的本地指针，以确保在
DEC Alpha上的内存排序，
\clnref{if_NULL,ret:a}在没有当前分配的
profile缓冲区时从此函数退出，而\clnref{if_oor,ret:b}
在\co{pcvalue}参数超出范围时
从此函数退出。
否则，\clnref{inc}递增以\co{pcvalue}参数为索引的
profile缓冲区条目。
注意，将大小与缓冲区一起存储可以保证
范围检查与缓冲区匹配，即使一个大缓冲区突然
被一个较小的缓冲区替换。
\end{fcvref}

\begin{fcvref}[ln:defer:Using RCU to Wait for NMIs to Finish:nmi_stop]
\Clnrefrange{b}{e}定义了\co{nmi_stop()}函数，
调用者负责互斥（例如，持有正确的锁）。
\Clnref{fetch}获取指向profile缓冲区的指针，
\clnref{if_NULL,ret}在没有缓冲区时退出函数。
否则，\clnref{NULL}将profile缓冲区指针置\co{NULL}
（使用\co{rcu_assign_pointer()}原语在弱序
机器上维持内存排序），
\clnref{sync_sched}等待一个RCU Sched grace period过去，
特别是等待所有不可抢占的代码区域
（包括NMI处理程序）完成。
一旦执行继续到\clnref{kfree}，我们就可以保证
任何获取了旧缓冲区指针的\co{nmi_profile()}
实例已经返回。
因此，可以安全地释放缓冲区，本例中使用的是
\co{kfree()}原语。
\end{fcvref}

\QuickQuiz{
	假设\co{nmi_profile()}函数是可抢占的。
	要使此示例正确工作，需要做哪些改变？
}\QuickQuizAnswer{
	一种方法是在\co{nmi_profile()}中使用
	\co{rcu_read_lock()}和\co{rcu_read_unlock()}，
	并将\co{synchronize_sched()}替换为\co{synchronize_rcu()}，
	也许如
	\cref{lst:defer:Using RCU to Wait for Mythical Preemptible NMIs to Finish}所示。
%
\begin{listing}
\begin{VerbatimL}
struct profile_buffer {
	long size;
	atomic_t entry[0];
};
static struct profile_buffer *buf = NULL;

void nmi_profile(unsigned long pcvalue)
{
	struct profile_buffer *p;

	rcu_read_lock();
	p = rcu_dereference(buf);
	if (p == NULL) {
		rcu_read_unlock();
		return;
	}
	if (pcvalue >= p->size) {
		rcu_read_unlock();
		return;
	}
	atomic_inc(&p->entry[pcvalue]);
	rcu_read_unlock();
}

void nmi_stop(void)
{
	struct profile_buffer *p = buf;

	if (p == NULL)
		return;
	rcu_assign_pointer(buf, NULL);
	synchronize_rcu();
	kfree(p);
}
\end{VerbatimL}
\caption{使用RCU等待虚构的可抢占NMI完成}
\label{lst:defer:Using RCU to Wait for Mythical Preemptible NMIs to Finish}
\end{listing}

	但NMI处理程序为什么会是可抢占的呢？？？
}\QuickQuizEnd

简而言之，RCU使得动态切换profile缓冲区变得容易
（你不妨\emph{试试}用原子操作高效地做到这一点，
或者用锁来做！）。
这是RCU纯粹形式的一种罕见使用。
RCU通常在更高的抽象层次上使用，
如后续各节所示。

\subsubsection{分阶段状态变更}
\label{sec:defer:Phased State Change}

\Cref{fig:defer:Phased State Change for Maintenance Operation}
展示了一个分阶段状态变更的时间线示例，
用于高效地处理维护操作。
如果没有正在进行的维护操作，常规操作
必须快速执行，例如不需要获取读写锁。
然而，如果有维护操作正在进行，常规操作
必须谨慎进行，考虑到因与维护操作
并发运行而增加的复杂性。
这意味着在维护操作期间，常规操作将招致更高的开销，
这也是维护操作通常被安排在
低负载时段执行的原因之一。

\begin{figure}
\centering
\resizebox{\twocolumnwidth}{!}{\includegraphics{defer/RCUphasedstatechange}}
\caption{维护操作的分阶段状态变更}
\label{fig:defer:Phased State Change for Maintenance Operation}
\end{figure}

在图中，这些看似矛盾的需求通过
在维护操作之前设置准备阶段、在之后设置清理阶段来解决，
在这些阶段中常规操作可以快速执行或谨慎执行。

\begin{listing}
\begin{fcvlabel}[ln:defer:Phased State Change for Maintenance Operations]
\begin{VerbatimL}[commandchars=\\\@\$]
bool be_careful;

void cco(void)
{
	rcu_read_lock();			\lnlbl@rrl$
	if (READ_ONCE(be_careful))		\lnlbl@if$
		cco_carefully();		\lnlbl@careful$
	else
		cco_quickly();			\lnlbl@quick$
	rcu_read_unlock();			\lnlbl@rul$
}

void maint(void)
{
	WRITE_ONCE(be_careful, true);		\lnlbl@tocareful$
	synchronize_rcu();			\lnlbl@sr1$
	do_maint();				\lnlbl@maint$
	synchronize_rcu();			\lnlbl@sr2$
	WRITE_ONCE(be_careful, false);		\lnlbl@toquick$
}
\end{VerbatimL}
\end{fcvlabel}
\caption{维护操作的分阶段状态变更}
\label{lst:defer:Phased State Change for Maintenance Operations}
\end{listing}

\begin{fcvref}[ln:defer:Phased State Change for Maintenance Operations]
此分阶段状态变更的伪代码示例如
\cref{lst:defer:Phased State Change for Maintenance Operations}所示。
常规操作由\co{cco()}在从\clnref{rrl}到\clnref{rul}的
RCU read-side critical section内执行。
其中\clnref{if}检查全局的\co{be_careful}标志，
相应地调用\co{cco_carefully()}或\co{cco_quickly()}。
\end{fcvref}

\begin{fcvref}[ln:defer:Phased State Change for Maintenance Operations]
这允许\co{maint()}函数在\clnref{tocareful}设置\co{be_careful}标志，
并在\clnref{sr1}等待一个RCU grace period。
当控制到达\clnref{maint}时，所有看到\co{be_careful}
为\co{false}值（因此可能调用
\co{cco_quickly()}函数）的\co{cco()}函数将已经完成其操作，
从而所有当前正在执行的\co{cco()}函数都将调用
\co{cco_carefully()}。
这意味着此时调用\co{do_maint()}函数是安全的。
\Clnref{sr2}接着等待所有可能与\co{do_maint()}并发
运行的\co{cco()}函数完成，最后
\clnref{toquick}将\co{be_careful}标志重新设为\co{false}。
\end{fcvref}

\QuickQuizSeries{%
\QuickQuizB{
	\cref{lst:defer:Phased State Change for Maintenance Operations}中
	\co{maint()}函数第二次调用\co{synchronize_rcu()}
	的意义何在？
	在清理阶段让\co{cco()}调用\co{cco_carefully()}或
	\co{cco_quickly()}不是都可以吗？
}\QuickQuizAnswerB{
	问题在于\co{cco()}函数从\co{be_careful}的加载
	与\co{cco_quickly()}函数执行的任何内存加载之间
	没有排序保证。
	由于没有排序保证，如果没有第二次调用
	\co{synchronize_rcu()}，内存排序可能导致
	\co{cco_quickly()}中的加载与\co{do_maint()}中的存储重叠。

	另一种替代方案是通过将\co{READ_ONCE()}
	改为\co{smp_load_acquire()}、将
	\co{WRITE_ONCE()}改为\co{smp_store_release()}来
	弥补去掉第二次\co{synchronize_rcu()}调用的影响，
	从而恢复所需的排序。
}\QuickQuizEndB

\QuickQuizE{
	你怎么能确定
	\cref{lst:defer:Phased State Change for Maintenance Operations}中
	\co{maint()}所示的代码确实有效？
}\QuickQuizAnswerE{
	按照一种流行的学派观点，你无法确定。

	但在这个例子中，那些愿意跳到
	\cref{chp:Formal Verification}
	和
	\cref{chp:Advanced Synchronization: Memory Ordering}
	的读者可能会发现几个LKMM litmus测试很有趣
	（\path{C-RCU-phased-state-change-1.litmus}和
	\path{C-RCU-phased-state-change-2.litmus}）。
	可以说这些测试证明了这段代码
	及其一个变体确实能正常工作。
}\QuickQuizEndE
}% End of \QuickQuizSeries

分阶段状态变更允许频繁的操作使用轻量级检查，
无需昂贵的锁获取或原子读-修改-写操作，
并在Linux内核中以\co{rcu_sync}~\cite{OlegNesterov2013rcusync}的
形式被使用，以实现具有轻量级读者的
读写信号量的变体。
分阶段状态变更只是在等待完成用例
（\cref{sec:defer:Wait for Pre-Existing Things to Finish}）之上
增加了一个受检查的状态变量，
因此也处于较低的抽象层次。

\subsubsection{仅添加链表}
\label{sec:defer:Add-Only List}

仅添加数据结构，以仅添加链表为代表，可用于
数量惊人的常见用例，其中最常见的可能是
变更日志。
仅添加数据结构是对RCU底层
发布/订阅机制的纯粹使用。

Pre-BSD路由表的仅添加变体可以从
\cref{lst:defer:RCU Pre-BSD Routing Table Lookup,lst:defer:RCU Pre-BSD Routing Table Add/Delete}
派生而来。
由于没有删除操作，\co{route_del()}和\co{route_cb()}
函数可以省略，连同\co{route_entry}结构的
\co{->rh}和\co{->re_freed}字段、
\co{route_lookup()}函数中的
\co{rcu_read_lock()}和\co{rcu_read_unlock()}调用，
以及所有剩余函数中\co{->re_freed}字段的所有使用。

当然，如果有许多\co{route_add()}函数的并发调用，
\co{routelock}将面临严重的争用，
如果使用无锁技术，\co{routelist}也将面临严重的内存争用。
避免这种争用的通常方法是使用一种并发友好的
数据结构，如哈希表（参见\cref{chp:Data Structures}）。
另一种方法是将每CPU数据结构定期合并
到一个全局数据结构中。

另一方面，如果永远不进行删除操作，
长时间的多个\co{route_add()}并发调用最终将
耗尽所有可用内存。
因此，大多数受RCU保护的数据结构也实现了删除操作。

\subsubsection{类型安全内存}
\label{sec:defer:Type-Safe Memory}

许多无锁算法不要求给定的数据元素
在引用它的整个RCU read-side critical section期间
保持相同的身份——但前提是该数据元素保持
相同的类型。
换言之，这些无锁算法可以容忍给定数据元素
在被引用期间被释放并重新分配为相同类型的结构，
但必须禁止类型更改。
这种保证在学术文献中被称为``\IX{type-safe memory}''~\cite{Cheriton96a}，
它比\cref{sec:defer:Existence Guarantee}中讨论的
\IXpl{existence guarantee}更弱，
因此使用起来相当困难。
Linux内核中的类型安全内存算法利用slab缓存，
通过\co{SLAB_TYPESAFE_BY_RCU}特别标记这些缓存，
以便在将释放的slab返回到系统内存时使用RCU。
这种RCU的使用保证了这种slab中任何正在使用的元素
将在所有已有的RCU read-side critical section持续期间
保留在该slab中，从而保持其类型。

\QuickQuiz{
	但如果多个线程中有任意长的RCU read-side critical section
	序列，使得在任何时间点系统中总有至少一个线程
	在执行RCU read-side critical section呢？
	这不会阻止\co{SLAB_TYPESAFE_BY_RCU} slab中的
	任何数据被归还给系统，可能导致
	OOM事件吗？
}\QuickQuizAnswer{
	确实可能存在任意长的时间段，
	在此期间总有至少一个线程处于RCU read-side critical section中。
	然而，\cref{sec:defer:Type-Safe Memory}中
	描述的关键词是``正在使用''和``已有的''。
	请记住，给定的RCU read-side critical section
	在概念上只被允许获取对在该临界区期间
	对读者可见的数据元素的引用。
	此外，请记住slab不能被归还给系统，
	直到其所有数据元素都被释放——事实上，
	RCU grace period不能在它们全部被释放之前开始。

	因此，slab缓存只需等待那些在最后一个元素
	释放之前开始的RCU read-side critical section。
	这反过来意味着，在最后一个元素释放之后
	开始的任何RCU grace period都可以——在该grace period
	结束后，slab可以被归还给系统。
}\QuickQuizEnd

重要的是要注意，\co{SLAB_TYPESAFE_BY_RCU}
\emph{绝不会}阻止\co{kmem_cache_alloc()}立即
重新分配刚通过\co{kmem_cache_free()}释放的内存！
事实上，\co{rcu_dereference()}刚返回的受
\co{SLAB_TYPESAFE_BY_RCU}保护的数据结构
可能被释放和重新分配任意多次，
即使在\co{rcu_read_lock()}的保护下也是如此。
相反，\co{SLAB_TYPESAFE_BY_RCU}通过阻止
\co{kmem_cache_free()}将完全释放的数据结构slab
归还给系统来运作，直到一个RCU grace period过去之后。
简而言之，虽然给定的RCU read-side critical section
可能看到给定的\co{SLAB_TYPESAFE_BY_RCU}数据元素
被任意频繁地释放和重新分配，
但在该临界区完成之前，该元素的类型保证不会改变。

因此，这些算法通常使用一个验证步骤来检查
确保新引用的数据结构确实是
所请求的那个~\cite[Section~2.5]{LaninShasha1986TSM}。
这些验证检查要求数据结构的部分内容
在释放-重新分配过程中保持不变。
这样的验证检查通常很难做对，并且可能
隐藏微妙而难以发现的错误。

因此，虽然基于类型安全的无锁算法
在极少数困难情况下可以极为有用，
但你应当尽可能使用存在性保证。
毕竟，更简单几乎总是更好！
另一方面，基于类型安全的无锁算法可以
提供更好的缓存局部性，从而提高性能。
这种改善的缓存局部性源于这样一个事实：
此类算法可以立即重新分配刚释放的内存块。
相比之下，基于存在性保证的算法必须等待
所有已有的读者完成后才能重新分配内存，
届时该内存可能已被从CPU缓存中淘汰。

如\cref{fig:defer:Relationships Between RCU Use Cases}所示，
RCU的类型安全内存用例结合了等待完成
和发布/订阅两个组件，但在Linux内核中还包括
slab分配器通过\co{SLAB_TYPESAFE_BY_RCU}标志
指定的延迟回收。

\subsubsection{存在性保证}
\label{sec:defer:Existence Guarantee}

Gamsa等人~\cite{Gamsa99}
讨论了\IXpl{existence guarantee}，并描述了一种
类似RCU的机制如何能够提供这些存在性保证
（参见PDF第7页的第5节），
\cref{sec:locking:Lock-Based Existence Guarantees}
讨论了如何通过加锁来保证存在性，
以及这样做所带来的缺点。
其效果是，如果在RCU read-side critical section中
访问了任何受RCU保护的数据元素，
则该数据元素在整个RCU read-side critical section
期间都保证继续存在。

\begin{listing}
\begin{fcvlabel}[ln:defer:Existence Guarantees Enable Per-Element Locking]
\begin{VerbatimL}[commandchars=\\\@\$]
int delete(int key)
{
	struct element *p;
	int b;

	b = hashfunction(key);			\lnlbl@hash$
	rcu_read_lock();			\lnlbl@rdlock$
	p = rcu_dereference(hashtable[b]);
	if (p == NULL || p->key != key) {	\lnlbl@chkkey$
		rcu_read_unlock();		\lnlbl@rdunlock1$
		return 0;			\lnlbl@ret_0:a$
	}
	spin_lock(&p->lock);			\lnlbl@acq$
	if (hashtable[b] == p && p->key == key) {\lnlbl@chkkey2$
		rcu_read_unlock();		\lnlbl@rdunlock2$
		rcu_assign_pointer(hashtable[b], NULL);\lnlbl@remove$
		spin_unlock(&p->lock);		\lnlbl@rel1$
		synchronize_rcu();		\lnlbl@sync_rcu$
		kfree(p);			\lnlbl@kfree$
		return 1;			\lnlbl@ret_1$
	}
	spin_unlock(&p->lock);			\lnlbl@rel2$
	rcu_read_unlock();			\lnlbl@rdunlock3$
	return 0;				\lnlbl@ret_0:b$
}
\end{VerbatimL}
\end{fcvlabel}
\caption{存在性保证实现每元素加锁}
\label{lst:defer:Existence Guarantees Enable Per-Element Locking}
\end{listing}

\begin{fcvref}[ln:defer:Existence Guarantees Enable Per-Element Locking]
\Cref{lst:defer:Existence Guarantees Enable Per-Element Locking}
展示了基于RCU的存在性保证如何通过一个
从哈希表中删除元素的函数来实现每元素加锁。
\Clnref{hash}计算哈希函数，\clnref{rdlock}进入一个
RCU read-side critical section。
如果\clnref{chkkey}发现哈希表的对应桶
为空或其中的元素不是我们要删除的，
则\clnref{rdunlock1}退出RCU read-side critical section，
\clnref{ret_0:a}返回失败。
\end{fcvref}

\QuickQuiz{
	如果我们需要删除的元素不是
	\cref{lst:defer:Existence Guarantees Enable Per-Element Locking}
	\clnrefr{ln:defer:Existence Guarantees Enable Per-Element Locking:chkkey}行上
	链表的第一个元素怎么办？
}\QuickQuizAnswer{
	与（有缺陷的）
	\cref{lst:locking:Per-Element Locking Without Existence Guarantees (Buggy!)}
	一样，这是一个非常简单的无链接哈希表，因此给定桶中
	唯一的元素就是第一个元素。
	读者再次被邀请将此示例改编为具有
	完整链接的哈希表。
	不那么精力充沛的读者可以参阅
	\cref{chp:Data Structures}。
}\QuickQuizEnd

\begin{fcvref}[ln:defer:Existence Guarantees Enable Per-Element Locking]
否则，\clnref{acq}获取更新侧自旋锁，
\clnref{chkkey2}然后检查该元素是否仍然是我们想要的。
如果是，\clnref{rdunlock2}离开RCU read-side critical section，
\clnref{remove}将其从表中移除，\clnref{rel1}释放锁，
\clnref{sync_rcu}等待所有已有的RCU read-side critical section完成，
\clnref{kfree}释放新移除的元素，
\clnref{ret_1}返回成功。
如果该元素已不再是我们想要的，\clnref{rel2}释放锁，
\clnref{rdunlock3}离开RCU read-side critical section，
\clnref{ret_0:b}返回删除指定键失败。
\end{fcvref}

\QuickQuizSeries{%
\QuickQuizB{
	\begin{fcvref}[ln:defer:Existence Guarantees Enable Per-Element Locking]
	为什么在
	\cref{lst:defer:Existence Guarantees Enable Per-Element Locking}的
	\clnref{rdunlock2}行退出RCU read-side critical section，
	然后在\clnref{rel1}才释放锁，是可以的？
	\end{fcvref}
}\QuickQuizAnswerB{
	\begin{fcvref}[ln:defer:Existence Guarantees Enable Per-Element Locking]
	首先请注意，\clnref{chkkey2}上的第二次检查是必要的，
	因为在我们等待获取锁的期间，
	其他CPU可能已经移除了此元素。
	然而，我们在获取锁时处于RCU read-side critical section中，
	这保证了此元素不可能被重新分配
	并重新插入到此哈希表中。
	此外，一旦我们获取了锁，锁本身就保证了
	元素的存在，因此我们不再需要处于
	RCU read-side critical section中。
	% The question as to whether it is necessary to re-check the
	% element's key is left as an exercise to the reader.
	% A re-check is necessary if the key can mutate or if it is
	% necessary to reject deleted entries (in cases where deletion
	% is recorded by mutating the key.
	\end{fcvref}
}\QuickQuizEndB
%
\QuickQuizM{
	\begin{fcvref}[ln:defer:Existence Guarantees Enable Per-Element Locking]
	为什么不在
	\cref{lst:defer:Existence Guarantees Enable Per-Element Locking}的
	\clnref{rdunlock3}行退出RCU read-side critical section之前，
	先在\clnref{rel2}行释放锁呢？
	\end{fcvref}
}\QuickQuizAnswerM{
	假设我们颠倒这两行的顺序。
	那么这段代码就容易受到以下事件序列的影响：
	\begin{enumerate}
	\begin{fcvref}[ln:defer:Existence Guarantees Enable Per-Element Locking]
	\item	CPU~0调用\co{delete()}，找到了要删除的元素，
		执行到\clnref{rdunlock2}。
		它尚未实际删除该元素，但即将这样做。
	\item	CPU~1并发调用\co{delete()}，试图
		删除同一个元素。
		然而，CPU~0仍持有锁，因此CPU~1
		在\clnref{acq}处等待。
	\item	CPU~0执行\clnref{remove,rel1}，
		并在\clnref{sync_rcu}处阻塞，等待CPU~1
		退出其RCU read-side critical section。
	\item	CPU~1现在获取了锁，但\clnref{chkkey2}上的测试
		失败了，因为CPU~0已经移除了该元素。
		CPU~1现在执行\clnref{rel2}
		（为了本快速测验的目的，我们已将其与\clnref{rdunlock3}
		交换）
		并退出其RCU read-side critical section。
	\item	CPU~0现在可以从\co{synchronize_rcu()}返回，
		因此执行\clnref{kfree}，将元素发送到
		空闲列表。
	\item	CPU~1现在试图为一个已被释放的元素释放锁，
		更糟糕的是，该元素可能已被
		重新分配为其他类型的数据结构。
		这是一个致命的内存损坏错误。
	\end{fcvref}
	\end{enumerate}
}\QuickQuizEndM
%
\QuickQuizE{
	\cref{lst:defer:Existence Guarantees Enable Per-Element Locking}
	中所示的基于RCU的算法看起来与
	\cref{lst:locking:Per-Element Locking With Lock-Based Existence Guarantees}
	中的非常相似，那为什么基于RCU的方法应该更好呢？
}\QuickQuizAnswerE{
	\Cref{lst:defer:Existence Guarantees Enable Per-Element Locking}
	将
	\cref{lst:locking:Per-Element Locking With Lock-Based Existence Guarantees}
	中所示的每元素\co{spin_lock()}和\co{spin_unlock()}
	替换为开销低得多的\co{rcu_read_lock()}和\co{rcu_read_unlock()}，
	从而极大地提高了性能和可扩展性。
	更多详情请参阅
	\cref{sec:datastruct:RCU-Protected Hash Table Performance}。
}\QuickQuizEndE
}

细心的读者会认识到这只是原始的等待完成主题
（\cref{sec:defer:Wait for Pre-Existing Things to Finish}）
的一个轻微变体，添加了发布/订阅、链式结构、
堆分配器（通常）和延迟回收，如
\cref{fig:defer:Relationships Between RCU Use Cases}所示。
他们可能还会注意到，与
\cref{sec:locking:Lock-Based Existence Guarantees}
中讨论的基于锁的存在性保证相比，
这具有死锁免疫的优势。

\subsubsection{轻量级垃圾收集器}
\label{sec:defer:Light-Weight Garbage Collector}

初次了解RCU的人常发出的一个惊叹是``RCU有点像
垃圾收集器！''~\cite{MattKline2023LockBasedLocklessFresh}。
这一惊叹有很大的道理，特别是当使用RCU来
实现依赖垃圾收集的非阻塞算法时，
但有时也可能会产生误导。

也许思考RCU与自动垃圾收集器（GC）之间关系的最佳方式是：
RCU类似于GC，因为收集的\emph{时机}是自动
确定的，但RCU与GC的不同之处在于：
\begin{enumerate*}[(1)]
\item 程序员必须手动指明给定的数据结构
何时有资格被收集，且
\item 程序员必须手动标记可能持有引用的
RCU read-side critical section。
\end{enumerate*}

尽管存在这些差异，两者的相似性确实很深。
事实上，我所知的第一个类RCU机制使用了
基于引用计数的垃圾收集器来处理
grace period~\cite{Kung80}，RCU与垃圾收集之间的联系
也在较近期被注意到~\cite{HarshalSheth2016goRCU}。

轻量级垃圾收集器用例与存在性保证用例非常相似，
只是在混合中添加了所需的非阻塞算法。
此轻量级垃圾收集器用例也可以与下一节
描述的存在性保证结合使用。

\subsubsection{仅删除链表}
\label{sec:defer:Delete-Only List}

仅删除链表是\cref{sec:defer:Add-Only List}中
介绍的仅添加链表的较不常用的对应物，
可以被看作是没有发布/订阅组件的存在性保证用例，
如\cref{fig:defer:Relationships Between RCU Use Cases}所示。
仅删除链表可用于链表的所有可能成员
在初始化时已知、且成员可以被移除的情况。
例如，链表的元素可能代表系统中可能发生故障
但在重启前无法修复或替换的硬件组件。

Pre-BSD路由表的仅删除变体可以从
\cref{lst:defer:RCU Pre-BSD Routing Table Lookup,lst:defer:RCU Pre-BSD Routing Table Add/Delete}
派生而来。
由于没有添加操作，可以省略\co{route_add()}函数，
或者将其使用限制在初始化阶段。
理论上，\co{route_lookup()}函数可以使用非RCU的迭代器，
不过在Linux内核中这会引起调试代码的警告。
此外，RCU迭代器的增量成本通常可以忽略不计。

因此，仅删除的情况通常使用为
添加和删除而设计的算法和数据结构。

\subsubsection{准读写锁}
\label{sec:defer:Quasi Reader-Writer Lock}

也许Linux内核中最常见的RCU使用方式是
在读密集场景中替换读写锁。
然而，RCU的这种使用方式在一开始对我来说并不明显。
事实上，我在20世纪90年代早期
实现了一个轻量级读写锁~\cite{WilsonCHsieh92a}\footnote{
	类似于2.4 Linux内核中的\co{brlock}和
	更新版本Linux内核中的\co{lglock}。}，
然后才实现了通用的RCU实现。
我为轻量级读写锁设想的每一个使用场景
都改用了RCU来实现。
事实上，轻量级读写锁过了三年多才
迎来它的第一次使用。
真是让我大出洋相！

RCU与读写锁之间的关键相似之处在于
两者都有可以并发执行的读侧临界区。
事实上，在某些情况下，可以机械地将RCU API
成员替换为对应的读写锁API成员。
但首先，为什么要这样做？

RCU的优势包括性能、
死锁免疫和实时延迟。
当然，RCU也有局限性，包括
读者和更新者并发运行的事实、低优先级的RCU读者
可以阻塞等待grace period过去的高优先级线程、
以及grace period延迟可达多个毫秒。
以下段落将讨论这些优势和局限性。

\paragraph{性能}

\begin{figure}
\centering
\resizebox{2.5in}{!}{\includegraphics{CodeSamples/defer/data/rcuscale.hps.2020.05.28a/rwlockRCUperf}}
\caption{RCU相比读写锁的性能优势}
\label{fig:defer:Performance Advantage of RCU Over Reader-Writer Locking}
\end{figure}

\cref{fig:defer:Performance Advantage of RCU Over Reader-Writer Locking}
展示了Linux内核RCU相对于读写锁的读侧性能优势，
数据在一台448 CPU、2.10\,GHz的Intel x86系统上生成。

\QuickQuizSeries{%
\QuickQuizB{
	什么鬼？
	在时钟周期为2.10\,GHz（几乎500皮秒）的情况下，
	你怎么能指望我相信RCU的开销
	不到300皮秒？
}\QuickQuizAnswerB{
	首先，考虑用于此测量的内部循环如下：

\begin{VerbatimN}
	for (i = nloops; i >= 0; i--) {
		rcu_read_lock();
		rcu_read_unlock();
	}
\end{VerbatimN}

	接下来，考虑\co{rcu_read_lock()}和
	\co{rcu_read_unlock()}的有效定义：

\begin{VerbatimN}
#define rcu_read_lock()   barrier()
#define rcu_read_unlock() barrier()
\end{VerbatimN}

	这些定义约束了编译器涉及内存引用的
	代码移动优化，但本身不生成任何指令。
	然而，如果循环变量保存在寄存器中，
	对\co{i}的访问将不算作内存引用。
	此外，编译器可以进行循环展开，
	使得生成的代码通过将\co{i}递增
	某个大于1的值来"执行"循环体的多次遍历。

	因此267皮秒的"测量值"不过是计时测量的
	固定开销除以包含\co{rcu_read_lock()}和
	\co{rcu_read_unlock()}调用的内部循环遍历次数，
	再加上操作\co{i}的代码除以循环展开因子。
	所以这个测量值确实有误——事实上，
	它将开销夸大了任意多个数量级。
	毕竟，就生成的机器指令而言，
	\co{rcu_read_lock()}和\co{rcu_read_unlock()}
	各自的实际开销精确为零。

	267皮秒的计时测量结果竟然是高估——
	这种事可不是天天都能碰到！
}\QuickQuizEndB

\QuickQuizM{
	本书的早期版本不是显示RCU读侧开销
	在亚皮秒范围内吗？
	发生了什么？？？
}\QuickQuizAnswerM{
	好记性！！！
	一些早期版本中的开销确实大约为
	100飞秒。

	发生的事情是RCU的使用在Linux内核中更广泛地传播，
	包括传播到会产生页面错误的代码中。
	在那个时候，\co{rcu_read_lock()}和\co{rcu_read_unlock()}
	在\co{CONFIG_PREEMPT=n}内核中是完全的空操作。
	不幸的是，这种情况允许编译器将
	产生页面错误的内存访问重排到RCU read-side critical section中。
	当然，页面错误可以阻塞，这会破坏这些临界区。

	而且这不是理论上的问题：
	2019年确实出现了一次故障。
	\ppl{Herbert}{Xu}追踪到了这个故障，
	\ppl{Linus}{Torvalds}因此
	排队了一个提交，将\co{rcu_read_lock()}和
	\co{rcu_read_unlock()}无条件升级为包含对
	\co{barrier()}的调用~\cite{LinusTorvalds2019:RCUreader.barrier}。
	虽然\co{barrier()}不生成代码，但它确实约束了
	编译器优化。
	因此，RCU广泛使用的代价是略高的
	\co{rcu_read_lock()}和\co{rcu_read_unlock()}开销。
	由此，Linux内核的RCU可谓是
	自身成功的受害者。

	当然，较早的结果也是在不同的系统上获得的，
	与\cref{fig:defer:Performance Advantage of RCU Over Reader-Writer Locking}
	中所示的不同。
	那么哪个变化影响更大——Linus的提交还是系统的更换？
	这个问题留给读者作为练习。
}\QuickQuizEndM

\QuickQuizE{
	为什么
	\cref{fig:defer:Performance Advantage of RCU Over Reader-Writer Locking}
	中\co{RCU}曲线的变化如此之大？
}\QuickQuizAnswerE{
	请记住这是对数-对数图，因此那些看似很大的
	\co{RCU}方差实际上只跨越了几百皮秒。
	如此短的时间内，任何因素都可能导致这种变化。
	然而，考虑到方差在CPU数量较少和较多时
	都会减小，一个假设是这种变化
	是由于从一个CPU迁移到另一个CPU造成的。

	是的，这些测量是在中断禁用的情况下进行的，
	但它们也是在客户操作系统内进行的，因此在
	hypervisor层面仍然可能发生抢占。
	此外，该系统具有超线程特性，
	运行此RCU工作负载的单个硬件线程能够消耗
	超过核心一半的资源。
	因此，总体吞吐量取决于给定客户操作系统
	有多少CPU共享核心。
	尝试通过以实时优先级运行客户操作系统
	（如Joel Fernandes所建议）来减少这些变化
	留给读者作为练习。
}\QuickQuizEndE
}                 % End of \QuickQuizSeries

请注意，在单个CPU上读写锁比RCU慢了一个数量级以上，
在192个CPU上则慢了\emph{四}个数量级以上。
相比之下，RCU的可扩展性相当好。
在两种情况下，误差条覆盖了30次运行中
测量值的全部范围，线条为中位数。

\begin{figure}
\centering
\resizebox{2.5in}{!}{\includegraphics{CodeSamples/defer/data/rcuscale.hps.2020.05.28a/rwlockRCUperfPREEMPT}}
\caption{可抢占RCU相比读写锁的性能优势}
\label{fig:defer:Performance Advantage of Preemptible RCU Over Reader-Writer Locking}
\end{figure}

从\co{CONFIG_PREEMPT}内核可以获得更温和的观点，
不过RCU仍然在单个CPU上以七倍的优势、
在192个CPU上以三个数量级的优势击败读写锁，如
\cref{fig:defer:Performance Advantage of Preemptible RCU Over Reader-Writer Locking}所示，
数据在同一台448 CPU、2.10\,GHz的x86系统上生成。
请注意读写锁在较多CPU数量时的高变异性。
误差条覆盖了数据的全部范围。

\QuickQuiz{
	既然系统有不少于448个硬件线程，
	为什么只测到192个CPU？
}\QuickQuizAnswer{
	因为生成此数据的脚本（\path{rcuscale.sh}）为
	每组采集的数据点启动一个客户操作系统，
	而在这个特定的系统上，\co{qemu}和KVM都限制了
	可以配置到给定客户操作系统中的CPU数量。
	是的，可以多运行几个CPU，但
	从二进制的角度来看，在256不可行的情况下，
	192是一个很好的整数。
}\QuickQuizEnd

\begin{figure}
\centering
\resizebox{2.5in}{!}{\includegraphics{CodeSamples/defer/data/rcuscale.hps.2020.05.28a/rwlockRCUperfwt}}
\caption{RCU与读写锁随临界区持续时间变化的比较，192个CPU}
\label{fig:defer:Comparison of RCU to Reader-Writer Locking as Function of Critical-Section Duration}
\end{figure}

当然，
\cref{fig:defer:Performance Advantage of RCU Over Reader-Writer Locking,%
fig:defer:Performance Advantage of Preemptible RCU Over Reader-Writer Locking}
中读写锁的低性能被不切实际的零长度临界区所夸大了。
RCU的性能优势随着临界区开销的增加而减小，如
\cref{fig:defer:Comparison of RCU to Reader-Writer Locking as Function of Critical-Section Duration}所示，
在同一系统上运行。
这里，y轴表示读侧原语的开销与临界区开销之和，
x轴表示以纳秒为单位的临界区开销。
但请注意对数刻度的y轴，这意味着曲线之间的
小间距仍然代表着显著的差异。
此图展示的是不可抢占RCU，但鉴于可抢占RCU的
读侧开销仅约三纳秒，其图形将与
\cref{fig:defer:Comparison of RCU to Reader-Writer Locking as Function of Critical-Section Duration}
几乎相同。

\QuickQuiz{
	为什么
	\cref{fig:defer:Comparison of RCU to Reader-Writer Locking as Function of Critical-Section Duration}
	中亚微秒持续时间的误差范围更大？
}\QuickQuizAnswer{
	因为较小的干扰对较小的测量值造成
	更大的相对误差。
	此外，Linux内核的\co{ndelay()}纳秒级原语
	（截至2020年）不如用于
	微秒及以上持续时间数据的\co{udelay()}原语精确。
	将其与
	\cref{fig:defer:Performance Advantage of RCU Over Reader-Writer Locking}
	中所示的零长度情况进行比较很有启发性。
}\QuickQuizEnd

读写锁有三条曲线，最上面的是100个CPU，
接下来是10个CPU，最下面的是1个CPU。
CPU数量越多、临界区越短，
RCU的性能优势越大。
以下事实进一步凸显了这些性能优势：
100 CPU的系统已不再罕见，
而且许多系统调用（因此它们包含的任何RCU read-side
critical section）在微秒内完成。

此外，如下一段所讨论的，
RCU读侧原语几乎完全免疫死锁。


\paragraph{死锁免疫}

虽然RCU为读密集工作负载提供了显著的性能优势，
创建RCU的主要原因之一实际上是它对读侧
死锁的免疫。
这种免疫源于RCU读侧原语不阻塞、不自旋、
甚至不进行后向分支的事实，因此它们的执行时间是确定的。
因此它们不可能参与死锁循环。

\QuickQuiz{
	这种死锁免疫有没有例外？如果有的话，
	什么事件序列可能导致死锁？
}\QuickQuizAnswer{
	导致涉及RCU读侧原语的死锁循环的一种方式是
	通过以下（非法的）语句序列：

\begin{VerbatimU}
rcu_read_lock();
synchronize_rcu();
rcu_read_unlock();
\end{VerbatimU}

	\co{synchronize_rcu()}不能在所有已有的
	RCU read-side critical section完成之前返回，
	但它被包含在一个RCU read-side critical section中，
	而该临界区不能在\co{synchronize_rcu()}返回之前完成。
	结果是经典的自死锁——当你试图在
	读持有读写锁时进行写获取时，效果相同。

	请注意，此自死锁场景不适用于RCU QSBR，
	因为\co{synchronize_rcu()}执行的上下文切换
	将充当此CPU的quiescent state，
	从而允许grace period完成。
	然而，如果说这更糟也不为过，因为
	RCU read-side critical section使用的数据可能
	由于grace period完成而被释放。
	而且Linux内核的lockdep设施会对你发出警告。

	简而言之，不要在RCU read-side critical section内
	调用同步的RCU更新侧原语，
	这些原语列在
	\cref{tab:defer:RCU Wait-to-Finish APIs}中。

	此外，在Linux内核中，RCU使用调度器，
	调度器也使用RCU。
	在某些情况下，RCU和调度器都必须小心
	以避免死锁。
}\QuickQuizEnd

RCU读侧死锁免疫的一个有趣后果是，
可以无条件地将RCU读者升级为RCU更新者。
使用读写锁尝试这样的升级会导致死锁。
以下代码片段展示了RCU读到更新的升级：

\begin{VerbatimN}[samepage=true]
rcu_read_lock();
list_for_each_entry_rcu(p, &head, list_field) {
	do_something_with(p);
	if (need_update(p)) {
		spin_lock(&my_lock);
		do_update(p);
		spin_unlock(&my_lock);
	}
}
rcu_read_unlock();
\end{VerbatimN}

请注意\co{do_update()}在锁的保护\emph{和}
RCU读侧保护下执行。

RCU死锁免疫的另一个有趣后果是它对
大量优先级反转问题的免疫。
例如，低优先级的RCU读者不能阻止高优先级的
RCU更新者获取更新侧锁。
类似地，低优先级的RCU更新者不能阻止高优先级的
RCU读者进入RCU read-side critical section。

\QuickQuiz{
	既免疫死锁又免疫优先级反转？？？
	听起来好得不像真的。
	我凭什么相信这是可能的？
}\QuickQuizAnswer{
	它确实有效。
	毕竟，如果它不能工作，Linux内核就无法运行。
}\QuickQuizEnd

\paragraph{实时延迟}

因为RCU读侧原语既不自旋也不阻塞，
它们提供了极好的实时延迟。
此外，如前所述，这意味着它们
免疫涉及RCU读侧原语和锁的优先级反转。

然而，RCU容易受到更微妙的优先级反转场景的影响，
例如，在\rt\ 内核中，等待RCU grace period过去的
高优先级进程可能被低优先级的RCU读者阻塞。
这可以通过使用RCU优先级
提升~\cite{PaulEMcKenney2007BoostRCU,DinakarGuniguntala2008IBMSysJ}来解决。

然而，使用RCU优先级提升要求\co{rcu_read_unlock()}
执行降优先级操作，这需要获取调度器锁。
因此，在调度器和RCU中需要一些小心来避免
死锁，截至v5.15 Linux内核，这要求RCU在持有
任何RCU锁时避免调用调度器。

这反过来意味着当启用RCU优先级提升时，
\co{rcu_read_unlock()}并不总是无锁的。
然而，如果其临界区没有被优先级提升，
\co{rcu_read_unlock()}仍然是无锁的。
此外，除非临界区被抢占，或者在-rt内核中获取了
非原始自旋锁，否则不会对临界区进行优先级提升。
这意味着从任何给定CPU上运行的最高优先级任务的
角度来看，\co{rcu_read_unlock()}通常是无锁的。

\paragraph{RCU读者和更新者并发运行}

因为RCU读者永远不自旋也不阻塞，而且更新者
不受任何形式的回滚或中止语义的约束，
RCU读者和更新者确实可以并发运行。
这意味着RCU读者可能访问过时数据，甚至可能
看到不一致性，这两者都可能使从读写锁
到RCU的转换变得不简单。

\begin{figure}
\centering
\resizebox{3in}{!}{\includegraphics{defer/rwlockRCUupdate}}
\caption{RCU与读写锁的响应时间对比}
\label{fig:defer:Response Time of RCU vs. Reader-Writer Locking}
\end{figure}

然而，在数量惊人的场景中，不一致性和
过时数据并不是问题。
经典的例子是网络路由表。
因为路由更新可能需要相当长的时间才能到达给定
系统（数秒甚至数分钟），在更新到达时，
系统早已按错误的方式发送了数据包相当长一段时间。
继续按错误的方式多发送几毫秒数据包
通常不是问题。
此外，因为RCU更新者可以在不等待RCU读者完成的情况下
进行更改，RCU读者很可能比
批量公平的读写锁读者更快地看到更改，如
\cref{fig:defer:Response Time of RCU vs. Reader-Writer Locking}所示。
这种更快的RCU响应时间后来也被Markov模型
分析所证实（虽然可能有些不切实际）~\cite[Figures 3 and 5]{VishakhaRamani2023LockBasedLocklessFresh}。

\QuickQuiz{
	但究竟有多少其他算法真正容忍过时和
	不一致的数据？
}\QuickQuizAnswer{
	相当多！

	请记住，光速有限意味着
	到达给定计算机系统的数据在到达时至少略有过时，
	对于天文数据来说则极其过时。
	光速有限也对来自不同来源或通过
	不同路径到达的数据的一致性设定了严格的限制。

	你不妨正视这样一个事实：物理定律
	与对完美新鲜度和一致性的朴素观念
	是不兼容的。
}\QuickQuizEnd

一旦收到更新，rwlock写者不能在最后一个读者完成之前
继续，后续读者也不能在写者完成之前继续。
然而，这些后续读者保证能看到新值，
如最右边方框的绿色底纹所示。
相比之下，RCU读者和更新者不会互相阻塞，这允许
RCU读者更快地看到更新后的值。
当然，因为它们的执行与RCU更新者重叠，
\emph{所有}RCU读者都很可能看到更新后的值，包括
在更新之前开始的三个读者。
尽管如此，只有绿色底纹的最右边RCU读者
才\emph{保证}能看到更新后的值。

读写锁和RCU只是提供不同的保证。
对于读写锁，在写者开始之后开始的任何读者
都保证能看到新值，而在写者自旋期间
试图开始的任何读者可能看到也可能看不到新值，
取决于所涉及的rwlock实现的读者/写者优先级。
相比之下，对于RCU，在更新者完成之后开始的任何读者
都保证能看到新值，而在更新者开始之后完成的
任何读者可能看到也可能看不到新值，取决于时序。

这里的关键点是，虽然读写锁确实
在计算机系统的范围内保证了一致性，
但有些情况下这种一致性的代价是
与外部世界增加的\emph{不一致性}，
这要归因于有限的光速和原子的非零大小。
换句话说，读写锁以对外部世界数据的
无声过时为代价获取了内部一致性。

请注意，如果在读持有读写锁时计算了一个值，
然后在释放该锁后使用该值，则
这种读写锁的使用场景就是在使用过时数据。
毕竟，该值所基于的量可以在锁释放后
随时更改。
这类读写锁使用场景通常很容易转换为RCU，
如\cref{lst:defer:Converting Reader-Writer Locking to RCU: Data,%
lst:defer:Converting Reader-Writer Locking to RCU: Search,%
lst:defer:Converting Reader-Writer Locking to RCU: Deletion}
及其附带文字所示。

\paragraph{低优先级RCU读者可阻塞高优先级回收者}

在Realtime RCU~\cite{DinakarGuniguntala2008IBMSysJ}或
SRCU~\cite{PaulEMcKenney2006c}中，
被抢占的读者会阻止grace period完成，即使
有高优先级任务在等待该grace period完成而阻塞。
Realtime RCU可以通过用\co{call_rcu()}
替代\co{synchronize_rcu()}或通过使用RCU优先级
提升~\cite{PaulEMcKenney2007BoostRCU,DinakarGuniguntala2008IBMSysJ}
来避免这个问题。
将来可能有必要为SRCU和RCU Tasks Trace
增加优先级提升，但在明确的现实需求被证明之前不会这样做。

\QuickQuiz{
	如果Tasks RCU Trace将来可能被优先级提升，
	为什么不对Tasks RCU和Tasks RCU Rude也这样做呢？
}\QuickQuizAnswer{
	也许可以，但可能性较小。

	对于Tasks RCU，请回忆quiescent state是
	自愿的上下文切换。
	因此，所有在自愿上下文切换后未被阻塞的任务
	可能都需要被提升，而降优先级的机制
	可能不会很优雅。

	对于Tasks RCU Rude，与旧的RCU Sched一样，
	任何可抢占的代码区域都是quiescent state。
	因此，唯一可能需要提升的任务是那些
	当前在禁用抢占的情况下运行的任务。
	但提升禁用抢占的任务的优先级没有任何效果。
	因此，优先级提升被引入Tasks RCU Rude的
	可能性加倍地小，至少在其当前形式下是如此。
}\QuickQuizEnd

\paragraph{RCU Grace Period持续数毫秒}

除了用户空间
RCU~\cite{MathieuDesnoyers2009URCU,PaulMcKenney2013LWNURCU}、
加速的grace period以及
\cref{chp:app:``Toy'' RCU Implementations}中描述的
几种"玩具"RCU实现外，
RCU grace period持续数毫秒。
虽然有许多技术可以使这种长延迟变得无害，
包括使用异步接口
（\co{call_rcu()}和\co{call_rcu_bh()}）或轮询接口
（\co{get_state_synchronize_rcu()}、\co{start_poll_synchronize_rcu()}
和\co{poll_state_synchronize_rcu()}），
但这种情况是RCU应当用于读密集场景这一
经验法则的主要原因。

如\cref{sec:defer:RCU Linux-Kernel API}所述，在Linux
内核中，可以通过加速的grace period获得更短的grace period，
例如，调用\co{synchronize_rcu_expedited()}
而不是\co{synchronize_rcu()}。
加速的grace period可以将延迟减少到数十微秒，
但代价是更高的CPU使用率和IPI。
额外的IPI在某些实时工作负载中尤其不受欢迎。

\paragraph{代码：
		 读写锁 vs.\@ RCU}

在最佳情况下，从读写锁到RCU的转换
非常简单，如
\cref{lst:defer:Converting Reader-Writer Locking to RCU: Data,%
lst:defer:Converting Reader-Writer Locking to RCU: Search,%
lst:defer:Converting Reader-Writer Locking to RCU: Deletion}所示，
均来自Wikipedia~\cite{WikipediaRCU}。

\begin{listing*}
{ \scriptsize
\begin{verbbox}
 1 struct el {                           1 struct el {
 2   struct list_head lp;                2   struct list_head lp;
 3   long key;                           3   long key;
 4   spinlock_t mutex;                   4   spinlock_t mutex;
 5   int data;                           5   int data;
 6   /* Other data fields */             6   /* Other data fields */
 7 };                                    7 };
 8 DEFINE_RWLOCK(listmutex);             8 DEFINE_SPINLOCK(listmutex);
 9 LIST_HEAD(head);                      9 LIST_HEAD(head);
\end{verbbox}
}
\hspace*{0.9in}\OneColumnHSpace{-0.5in}
\IfEbookSize{\hspace*{-1.05in}}{}\theverbbox
\caption{将读写锁转换为RCU：
						    数据}
\label{lst:defer:Converting Reader-Writer Locking to RCU: Data}
\end{listing*}

\begin{listing*}
{ \scriptsize
\begin{verbbox}
 1 int search(long key, int *result)     1 int search(long key, int *result)
 2 {                                     2 {
 3   struct el *p;                       3   struct el *p;
 4                                       4
 5   read_lock(&listmutex);              5   rcu_read_lock();
 6   list_for_each_entry(p, &head, lp) { 6   list_for_each_entry_rcu(p, &head, lp) {
 7     if (p->key == key) {              7     if (p->key == key) {
 8       *result = p->data;              8       *result = p->data;
 9       read_unlock(&listmutex);        9       rcu_read_unlock();
10       return 1;                      10       return 1;
11     }                                11     }
12   }                                  12   }
13   read_unlock(&listmutex);           13   rcu_read_unlock();
14   return 0;                          14   return 0;
15 }                                    15 }
\end{verbbox}
}
\hspace*{0.9in}\OneColumnHSpace{-0.5in}
\IfEbookSize{\hspace*{-1.05in}}{}\theverbbox
\caption{将读写锁转换为RCU：
						    搜索}
\label{lst:defer:Converting Reader-Writer Locking to RCU: Search}
\end{listing*}

\begin{listing*}
{ \scriptsize
\begin{verbbox}
 1 int delete(long key)                  1 int delete(long key)
 2 {                                     2 {
 3   struct el *p;                       3   struct el *p;
 4                                       4
 5   write_lock(&listmutex);             5   spin_lock(&listmutex);
 6   list_for_each_entry(p, &head, lp) { 6   list_for_each_entry(p, &head, lp) {
 7     if (p->key == key) {              7     if (p->key == key) {
 8       list_del(&p->lp);               8       list_del_rcu(&p->lp);
 9       write_unlock(&listmutex);       9       spin_unlock(&listmutex);
                                        10       synchronize_rcu();
10       kfree(p);                      11       kfree(p);
11       return 1;                      12       return 1;
12     }                                13     }
13   }                                  14   }
14   write_unlock(&listmutex);          15   spin_unlock(&listmutex);
15   return 0;                          16   return 0;
16 }                                    17 }
\end{verbbox}
}
\hspace*{0.9in}\OneColumnHSpace{-0.5in}
\IfEbookSize{\hspace*{-1.05in}}{}\theverbbox
\caption{将读写锁转换为RCU：
						    删除}
\label{lst:defer:Converting Reader-Writer Locking to RCU: Deletion}
\end{listing*}

然而，转换并不总是这么直接。
这是因为\cref{lst:defer:Converting Reader-Writer Locking to RCU: Deletion}
中的\co{spin_lock()}和\co{synchronize_rcu()}
都不排斥\cref{lst:defer:Converting Reader-Writer Locking to RCU: Search}
中的读者。
首先，\co{spin_lock()}与\co{rcu_read_lock()}和
\co{rcu_read_unlock()}没有任何交互，因此不排斥它们。
其次，虽然\co{write_lock()}和\co{synchronize_rcu()}
都等待已有读者完成，但只有\co{write_lock()}阻止
后续读者开始。\footnote{
	感谢向Paul指出这一点的人。}
因此，\co{synchronize_rcu()}无法排斥读者。
尽管如此，使用读写锁的大量场景
都可以转换为RCU。

更复杂的用RCU替换读写锁的案例可以在其他地方找到~\cite{NeilBrown2015PathnameLookup,NeilBrown2015RCUwalk}。

\paragraph{语义：
		      读写锁 vs.\@ RCU}

在上一节的基础上扩展，读写锁的语义可以
粗略且非正式地概括为以下三个时间约束：

\begin{enumerate}
\item	写侧获取等待任何读持有者释放锁。
\item	写侧获取等待任何写持有者释放锁。
\item	读侧获取等待任何写持有者释放锁。
\end{enumerate}

RCU完全消除了约束\#3，并将其他两个弱化如下：

\begin{enumerate}
\item	写者在进入更新的破坏性阶段（通常是
	释放内存）之前等待任何已有的读持有者完成。
\item	写者之间根据需要进行同步。
\end{enumerate}

正是这种弱化使得RCU实现能够获得
出色的性能和可扩展性。
它也允许RCU实现前面提到的无条件
读到写升级，这在读写锁中既吸引人又容易死锁。
使用RCU的代码可以通过数量惊人的
方式来弥补这种弱化，但最常见的是施加空间约束：

\begin{enumerate}
\item	新数据被放置在新分配的内存中。
\item	旧数据被释放，但仅在以下条件满足后：
	\begin{enumerate}
	\item	该数据已被断开链接使其对后续读者不可访问，且
	\item	随后的一个RCU grace period已经过去。
	\end{enumerate}
\end{enumerate}

当然，有些读写锁用例的语义不适合
RCU的弱化语义，但Linux内核中的经验表明
超过80\pct\ 的读写锁实际上可以被RCU替换。
例如，一个常见的读写锁用例是在持有锁时计算某个值，
然后在释放该锁后使用该值。
这种用例导致使用过时数据，因此通常能够适应
RCU的较弱语义。

\begin{listing}
\input{CodeSamples/defer/singleton@get.fcv}
\caption{RCU单例获取}
\label{lst:defer:Singleton Get}
\end{listing}

\begin{listing}
\input{CodeSamples/defer/singleton@set.fcv}
\caption{RCU单例设置}
\label{lst:defer:Singleton Set}
\end{listing}

\begin{fcvref}[ln:defer:singleton:get]
时间约束和空间约束的这种交互通过
\cref{fig:defer:Insertion With Concurrent Readers,fig:defer:Deletion With Concurrent Readers}
中所示的RCU单例数据结构来说明。
此结构定义在\cref{lst:defer:Singleton Get}的
\clnrefrange{myconfig.b}{myconfig.e}，包含两个整数字段
\co{->a}和\co{->b}（\path{singleton.c}）。
此结构的当前实例由\clnref{myconfig.e}上定义的
\co{curconfig}指针引用。
\end{fcvref}

\begin{fcvref}[ln:defer:singleton:get]
当前结构的字段通过\co{cur_a}和\co{cur_b}参数
传回给定义在\clnrefrange{get_config.b}{get_config.e}上的
\co{get_config()}函数。
这两个字段可以略微过时，但它们绝对必须
彼此一致。
\co{get_config()}函数在从\clnref{rrl}开始、
在\clnref{rrul1}或\clnref{rrul2}结束的
RCU read-side critical section内提供这种一致性，
从而提供所需的时间同步。
\Clnref{rderef}获取指向当前\co{myconfig}结构的指针。
无论由于\co{set_config()}函数的调用而发生的任何并发更改，
此结构都将被使用，从而提供所需的空间同步。
如果\clnref{nullchk}确定\co{curconfig}指针为\co{NULL}，
\clnref{retfail}返回失败。
否则，\clnref{copya,copyb}复制出\co{->a}和\co{->b}字段，
\clnref{retsuccess}返回成功。
这些\co{->a}和\co{->b}字段来自同一个\co{myconfig}
结构，而RCU read-side critical section阻止此结构
被释放，从而保证这两个字段彼此一致。
\end{fcvref}

\begin{fcvref}[ln:defer:singleton:set]
该结构由\cref{lst:defer:Singleton Set}中所示的
\co{set_config()}函数更新。
\Clnrefrange{allocinit.b}{allocinit.e}分配并初始化
一个新的\co{myconfig}结构。
\Clnref{xchg}原子地将指向此新结构的指针与
\co{curconfig}中指向旧结构的指针交换，
同时在\co{xchg()}操作前后提供完整的内存排序，
从而一方面提供所需的更新者/读者空间同步，
另一方面提供所需的更新者/更新者同步。
如果\clnref{if}确定指向旧结构的指针确实
非\co{NULL}，\clnref{sr}等待一个grace period（从而提供
所需的读者/更新者时间同步），\clnref{free}
释放旧结构，确信不再有任何读者仍在引用它。
\end{fcvref}

\begin{figure*}
\centering
\resizebox{\onecolumntextwidth}{!}{\includegraphics{defer/RCUspacetime}}
\caption{RCU空间/时间同步}
\label{fig:defer:RCU Spatial/Temporal Synchronization}
\end{figure*}

\Cref{fig:defer:RCU Spatial/Temporal Synchronization}在左右两侧
展示了\co{get_config()}的简略表示，
在中间展示了\co{set_config()}的类似简略表示。
时间从上到下推进，\co{curconfig}引用的对象
的地址空间从左到右推进。
带有逗号分隔数字的方框各代表一个\co{myconfig}
结构，其约束是\co{->b}是\co{->a}的平方。
每个蓝色的点划线箭头代表与旧结构
（左侧，包含``5,25''）的交互，每个绿色虚线箭头代表
与新结构（右侧，包含``9,81''）的交互。

黑色虚线箭头代表左右RCU读者与中间RCU grace period
之间的时间关系，每个箭头从较早的事件指向较晚的事件。
\co{synchronize_rcu()}的调用在最左边的\co{rcu_read_lock()}之后，
因此该\co{synchronize_rcu()}调用必须在
对应的\co{rcu_read_unlock()}之后才能返回。
相比之下，\co{synchronize_rcu()}的调用先于
最右边的\co{rcu_read_lock()}，
这允许同一\co{synchronize_rcu()}的返回忽略
对应的\co{rcu_read_unlock()}。
这些时间关系阻止了\co{myconfig}结构在
RCU读者仍在访问时被释放。

两条水平灰色虚线代表不同读者获得不同结果的
时间段，但每个读者只会看到两个对象中的一个。
在第一条水平线之前结束的所有读者将看到
最左边的\co{myconfig}结构，而在第二条水平线之后
开始的所有读者将看到最右边的结构。
在两条线之间，即在grace period期间，
不同的读者可能看到不同的对象，但只要每个读者
只加载\co{curconfig}指针一次，每个读者将看到
其\co{myconfig}结构的一致视图。

\QuickQuiz{
	但RCU grace period不是在调用\co{synchronize_rcu()}
	之后的某个时刻才开始的，而不是在
	\co{xchg()}语句的中间开始的吗？
}\QuickQuizAnswer{
	具体是哪个grace period？

	更新者被要求等待至少一个在移除操作
	（在本例中为\co{xchg()}）之时或之后开始的grace period。
	因此在
	\cref{fig:defer:RCU Spatial/Temporal Synchronization}中，
	所示的grace period尽可能早地开始，
	并延伸到\co{synchronize_rcu()}返回。
	这是一个完全合法的grace period，
	对应于该\co{xchg()}语句所进行的更改。
}\QuickQuizEnd

简而言之，当操作合适的链式数据结构时，RCU
结合时间和空间同步来近似读写锁，
其中RCU read-side critical section充当
读写锁的读者，如
\cref{fig:defer:Relationships Between RCU Use Cases,fig:defer:RCU Spatial/Temporal Synchronization}所示。
RCU的时间同步由读侧标记提供，
例如\co{rcu_read_lock()}和\co{rcu_read_unlock()}，以及
更新侧grace period原语，例如\co{synchronize_rcu()}
或\co{call_rcu()}。
空间同步由读侧的\co{rcu_dereference()}系列原语提供，
每个原语订阅一个由\co{rcu_assign_pointer()}发布的版本。\footnote{
	最好是将\co{rcu_dereference()}和
	\co{rcu_assign_pointer()}都嵌入到更高层的API中。}
RCU将时间和空间同步相结合，这与
\cref{sec:SMPdesign:Code Locking,sec:SMPdesign:Data Locking,sec:locking:Inefficiency}
中展示的方案形成对比，在那些方案中，
时间和空间同步分别由锁和静态数据结构布局提供。

\QuickQuiz{
	RCU是唯一以这种方式结合时间和空间同步的
	同步机制吗？
}\QuickQuizAnswer{
	完全不是。

	Hazard pointer可以被认为以类似的方式
	结合了时间和空间同步。
	参考\cref{lst:defer:Hazard-Pointer Recording and Clearing}，
	\co{hp_record()}函数获取引用的操作
	同时提供空间和时间同步，订阅
	一个版本并标记引用的开始。
	因此，此函数结合了RCU的
	\co{rcu_read_lock()}和\co{rcu_dereference()}的效果。
	现在参考\cref{lst:defer:Hazard-Pointer Scanning and Freeing}，
	\co{hp_clear()}函数释放引用的操作
	提供标记引用结束的时间同步，
	因此类似于RCU的\co{rcu_read_unlock()}。
	\co{hazptr_free_later()}函数退役
	一个受hazard pointer保护的对象提供时间同步，
	类似于RCU的\co{call_rcu()}。
	用于变更受hazard pointer保护的结构的原语
	提供空间同步，类似于RCU的
	\co{rcu_assign_pointer()}。

	另外，也可以通过与引用计数的类比
	来理解hazard pointer。
	这样做就会发现，引用计数也以上述
	hazard pointer描述的方式结合了
	时间和空间同步。
}\QuickQuizEnd

\subsubsection{准引用计数}
\label{sec:defer:Quasi Reference Count}

因为在有RCU read-side critical section正在进行时
不允许grace period完成，RCU读侧原语可以
用作一种受限的引用计数机制。
例如，考虑以下代码片段：

\begin{VerbatimN}
rcu_read_lock();  /* acquire reference. */
p = rcu_dereference(head);
/* do something with p. */
rcu_read_unlock();  /* release reference. */
\end{VerbatimN}

\co{rcu_read_lock()}和\co{rcu_dereference()}
原语的组合可以被看作是获取了对\co{p}的引用，
因为在\co{rcu_dereference()}赋值给\co{p}之后
开始的grace period不可能在我们到达匹配的
\co{rcu_read_unlock()}之前结束。
这种引用计数方案是受限的，因为在RCU read-side
critical section内禁止等待RCU grace period，
也禁止将一个RCU read-side critical section
的引用从一个任务传递给另一个任务。

不管这些限制如何，
以下代码可以安全地删除\co{p}：

\begin{VerbatimN}
spin_lock(&mylock);
p = head;
rcu_assign_pointer(head, NULL);
spin_unlock(&mylock);
/* Wait for all references to be released. */
synchronize_rcu();
kfree(p);
\end{VerbatimN}

对\co{head}的赋值阻止了任何将来对\co{p}的
引用被获取，而\co{synchronize_rcu()}
等待任何先前获取的引用被释放。

\QuickQuiz{
	但等等！
	这与将RCU视为读写锁替代品时
	可能使用的代码完全相同！
	这是怎么回事？
}\QuickQuizAnswer{
	这是玩具示例定律的效果：
	超过某一点后，代码片段看起来都一样。
	唯一的区别在于我们如何思考代码。
	例如，一个\co{atomic_inc()}操作做什么？
	它可能是在获取我们已有引用的对象的
	另一个显式引用，可能是递增
	一个经常读/偶尔更新的统计计数器，
	可能是签入一个HPC风格的barrier，
	或者许多其他事情中的任何一个。

	然而，这些区别可能极其重要。
	仅举一个重要性的例子：如果我们将RCU
	视为受限的引用计数方案，
	我们就永远不会被误导认为更新会排斥
	RCU read-side critical section。

	然而，将RCU视为读写锁的替代品通常很有用，
	例如，当你正在用RCU替换读写锁时。
}\QuickQuizEnd

当然，RCU也可以与传统引用计数结合使用，
如\cref{sec:together:Refurbish Reference Counting}中所讨论的。

\begin{figure}
\centering
\resizebox{2.5in}{!}{\includegraphics{CodeSamples/defer/data/rcuscale.hps.2020.05.28a/refcntRCUperf}}
\caption{RCU与引用计数的性能对比}
\label{fig:defer:Performance of RCU vs. Reference Counting}
\end{figure}

\begin{figure}
\centering
\resizebox{2.5in}{!}{\includegraphics{CodeSamples/defer/data/rcuscale.hps.2020.05.28a/refRCUperfPREEMPT}}
\caption{可抢占RCU与引用计数的性能对比}
\label{fig:defer:Performance of Preemptible RCU vs. Reference Counting}
\end{figure}

但为什么要费心呢？
同样，部分答案是性能，如
\cref{fig:defer:Performance of RCU vs. Reference Counting,%
fig:defer:Performance of Preemptible RCU vs. Reference Counting}所示，
同样展示了在448 CPU、2.1\,GHz的Intel x86系统上
分别对不可抢占和可抢占的Linux内核RCU采集的数据。
不可抢占RCU相对于\IXalt{reference counting}{reference count}的
优势从1个CPU上超过一个数量级
到192个CPU上约四个数量级。
可抢占RCU的优势从1个CPU上约三倍
到192个CPU上约三个数量级。

\begin{figure}
\centering
\resizebox{2.5in}{!}{\includegraphics{CodeSamples/defer/data/rcuscale.hps.2020.05.28a/refRCUperfwt}}
\caption{RCU与引用计数的响应时间对比，192个CPU}
\label{fig:defer:Response Time of RCU vs. Reference Counting}
\end{figure}

然而，与读写锁一样，RCU的性能优势
在短持续时间的临界区和大量CPU时最为显著，如
\cref{fig:defer:Response Time of RCU vs. Reference Counting}
中在同一系统上所示。
此外，与读写锁一样，许多系统调用
（因此它们包含的任何RCU read-side critical section）
在几微秒内完成。

虽然传统引用计数通常与特定的数据结构相关联，
或者可能与特定的一组数据结构相关联，
但这种方法确实有一些缺点。
例如，为大量不同的数据结构维护一个
全局引用计数器通常会导致包含引用计数的
缓存行反复弹跳。
正如我们在
\crefrange{fig:defer:Performance of RCU vs. Reference Counting}{fig:defer:Response Time of RCU vs. Reference Counting}
中看到的，这种缓存行弹跳会严重降低性能。

相比之下，RCU轻量级的\co{rcu_read_lock()}、
\co{rcu_dereference()}和\co{rcu_read_unlock()}读侧原语
允许极其频繁的读侧使用，性能降低可以忽略不计。
只不过\co{rcu_dereference()}的调用并没有做任何
特定于获取被指向对象引用的事情。
繁重的工作实际上由\co{rcu_read_lock()}和
\co{rcu_read_unlock()}原语及其与RCU grace period的
交互来完成。

忽略这些\co{rcu_dereference()}的调用，
RCU可以被看作一种"批量引用计数"机制，
其中每次调用\co{rcu_read_lock()}就获取了
对每一个受RCU保护的对象的引用，而且几乎没有开销。
然而，RCU附带的限制可能相当繁重。
例如，在很多情况下，Linux内核禁止在RCU
read-side critical section中睡眠的规定会破坏
整个目的。
这种情况可能更适合使用
\cref{sec:defer:Hazard Pointers}中描述的hazard pointer机制。
代码很少睡眠的情况已通过以下方式处理：
在常见的不睡眠情况下使用RCU作为引用计数，
在需要睡眠时桥接到显式引用计数器。

另外，当引用必须由单个任务在
包含睡眠的代码段中持有时，可以使用Sleepable
RCU（SRCU）~\cite{PaulEMcKenney2006c}来适应。
这无法涵盖引用从一个任务"传递"给另一个任务的
不少见情况，例如，在启动I/O时获取引用
并在相应的完成中断处理程序中释放引用。
同样，这种情况可能更适合用显式引用计数器
或hazard pointer来处理。

当然，SRCU也有自身的限制，即
\co{srcu_read_lock()}的返回值必须传入
对应的\co{srcu_read_unlock()}，并且不得从
硬件中断处理程序或\IXacrf{nmi}处理程序中
调用SRCU原语。
关于这一限制带来多大问题以及如何最好地处理它，
目前尚无定论。

然而，在引用被限制在单个CPU或任务范围内的
常见情况下，RCU可以用作高性能且高度
可扩展的引用计数机制。

如\cref{fig:defer:Relationships Between RCU Use Cases}所示，
准引用计数将RCU读者作为个别或批量
引用计数添加，并可能在边界情况下
桥接到引用计数器。

\subsubsection{准多版本并发控制}
\label{sec:defer:Quasi Multi-Version Concurrency Control}

RCU也可以被看作是一种简化的多版本并发控制
（MVCC）机制，具有较弱的一致性准则。
多版本方面在
\cref{sec:defer:Maintain Multiple Versions of Recently Updated Objects}
中已有涉及。
然而，在其原始形式中，RCU仅在给定的
受RCU保护的数据元素内提供版本一致性。

尽管如此，有些情况需要跨多个数据元素的
一致性和新鲜数据。
幸运的是，有许多方法可以避免不一致性
和过时数据，包括以下几种：

\begin{enumerate}
\item	将RCU读者封装在序列锁读者内，强制
	在发生更新时重试RCU读者，
	如\cref{sec:together:Correlated Data Elements}
	和\cref{sec:together:Atomic Move}中所述。
\item	将必须一致的数据放入链式数据结构的
	单个元素中，并避免更新对RCU读者可见的
	元素中的那些字段。
	获取对任何此类元素的引用的RCU读者
	因此保证能看到一致的值。
	更多细节参见\cref{sec:together:Correlated Fields}。
\item	使用一个保护"已删除"标志的每元素锁来允许
	RCU读者拒绝过时
	数据~\cite{PaulEdwardMcKenneyPhD,Arcangeli03}。
\item	提供一个被所有数据元素引用的存在标志，
	这些数据元素的更新对RCU读者来说应当
	表现为原子的~\cite{PaulEMcKennneyAtomicTreeN4037,PaulEMcKennneyAtomicTreeCPPCON2014,PaulEMcKenneyIssaquahUpdate2015,PaulEMcKenney2016IssaquahACMApp,PaulEMcKenney2016IssaquahCPPCON}。
\item	使用多种基于计数器的
	方法~\cite{PaulEMcKenney2008cyclicRCU,PaulEMcKenney2010cyclicRCU,PaulEMcKenney2011cyclicparallelRCU,PaulEMcKenney2014cyclicRCU,Matveev:2015:RLS:2815400.2815406,Kim:2019:MSR:3297858.3304040}。
	在这些方法中，更新者维护一个版本号并维护
	指向给定数据旧版本的链接。
	读者获取当前版本号的快照，如有必要，
	沿着链接找到与该快照一致的版本。
\end{enumerate}

简而言之，当使用RCU来近似多版本并发控制时，
你只需为你实际需要的一致性级别付费。

如\cref{fig:defer:Relationships Between RCU Use Cases}所示，
准多版本并发控制基于存在性保证，
添加了读侧快照操作和对读者与写者的约束，
约束的确切形式由一致性需求决定，
如上所述。

\subsubsection{RCU使用总结}
\label{sec:defer:RCU Usage Summary}

在其核心，RCU不多不少就是一个提供以下功能的API：

\begin{enumerate}
\item	用于添加新数据的发布-订阅机制，
\item	一种等待已有RCU读者完成的方式，以及
\item	维护多个版本以允许更改而不伤害或过度延迟
	并发RCU读者的规范。
\end{enumerate}

话虽如此，可以在RCU之上构建更高层次的结构，
包括前面各节描述的各种用例。
此外，我毫不怀疑新的用例将继续被发现，
无论是对RCU还是对许多其他同步原语。
因此，RCU的用例在概念上比RCU本身更复杂，
如\cpageref{sec:defer:Mysteries RCU Use Cases}所暗示的那样。

\QuickQuiz{
	在\cref{sec:defer:RCU for Pre-BSD Routing}中的
	Pre-BSD路由示例最符合哪个用例？
}\QuickQuizAnswer{
	Pre-BSD路由可以被认为适合准读写锁、
	准引用计数或准多版本并发控制中的任何一个。
	无论如何，代码都是一样的。
	这类似于\co{atomic_inc()}之类的东西，
	另一种可以有多种用途的工具。
}\QuickQuizEnd

\begin{figure}
\centering
\resizebox{3in}{!}{\includegraphics{defer/RCUApplicability}}
\caption{RCU的适用领域}
\label{fig:defer:RCU Areas of Applicability}
\end{figure}

同时，\cref{fig:defer:RCU Areas of Applicability}
展示了RCU最有用之处的一些粗略经验法则。

如图右上角蓝色方框所示，RCU
在读密集且允许过时和不一致数据的情况下效果最好
（但有关过时和不一致数据的更多信息请参见下文）。
Linux内核中此案例的典型例子是路由表。
因为路由更新可能需要数秒甚至数分钟
才能在Internet上传播，所以系统已经
按错误的方式发送数据包相当长一段时间了。
继续以较小的概率按错误的方式多发送一些数据包
几毫秒几乎从来不是问题。

如果你有一个读密集的工作负载，其中需要一致的数据，
RCU也能很好地工作，如绿色的"读密集，需要一致数据"
方框所示。
此案例的一个例子是Linux内核中从用户级
System-V信号量ID到相应内核内数据结构的映射。
信号量的使用频率远高于其创建和销毁频率，
因此这个映射是读密集的。
然而，对已被删除的信号量执行信号量操作
将是错误的。
这种对一致性的需求通过使用
内核内信号量数据结构中的锁以及在删除信号量时
设置的"已删除"标志来处理。
如果用户ID映射到一个设置了"已删除"标志的
内核内数据结构，则该数据结构被忽略，
用户ID被标记为无效。

虽然这要求读者获取代表信号量本身的
数据结构的锁，但它允许它们省略
映射数据结构的锁。
因此读者无锁地遍历用于从ID映射到
数据结构的树，这反过来极大地改善了性能、
可扩展性和实时响应。

如黄色的"读写"方框所示，RCU对于
需要一致数据的读写工作负载也可能有用，
尽管通常需要结合许多其他同步原语。
例如，近期Linux内核中的目录项缓存使用RCU
结合序列锁、每CPU锁和每数据结构锁，
在常见情况下允许路径名的无锁遍历。
虽然RCU在此读写情况下可以非常有益，
但相应的代码通常比读密集情况更复杂。

最后，如图左下角红色方框所示，
需要一致数据的更新密集工作负载很少是
使用RCU的好场所，不过也有一些
例外~\cite{MathieuDesnoyers2012URCU}。
例如，如\cref{sec:defer:Type-Safe Memory}所述，
在Linux内核中，\co{SLAB_TYPESAFE_BY_RCU}
slab分配器标志为RCU读者提供类型安全内存，
这可以极大地简化\IXacrl{nbs}和其他无锁算法。
此外，如果罕见的读者处于实时系统的关键代码路径上，
为这些读者使用RCU可能带来实时响应方面的
收益，足以弥补增加的更新侧开销，
如\cref{sec:advsync:The Role of RCU}中所讨论的。

简而言之，RCU是一个API，包括用于添加新数据的
发布-订阅机制、等待已有RCU读者完成的方式，
以及维护多个版本以允许更新不伤害或过度
延迟并发RCU读者的规范。
此RCU API最适合读密集的场景，尤其是在
应用程序可以容忍过时和不一致数据的情况下。

% defer/rcurelated.tex
% mainfile: ../perfbook.tex
% SPDX-License-Identifier: CC-BY-SA-3.0

\subsection{RCU 相关工作}
\label{sec:defer:RCU Related Work}
\OriginallyPublished{sec:defer:RCU Related Work}{RCU Related Work}{Linux Weekly News}{PaulEMcKenney2014ReadMostly}
\OriginallyPublished{sec:defer:RCU Related Work}{RCU Related Work}{Linux Weekly News}{PaulEMcKenney2015ReadMostly}

已知最早提及类似 RCU 机制的是
\ppl{Donald}{Knuth}~\cite[page 413 of Fundamental Algorithms]{Knuth73}
针对\ppl{Joseph}{Weizenbaum}的 FORTRAN SLIP 列表处理
工具~\cite{Weizenbaum:1963:SLP:367593.367617}的一个 bug 报告。
Knuth 报告这个 bug 是正当的，因为 SLIP 没有任何
grace-period 保证的概念。

已知最早的非 bug 报告形式的类似 RCU 机制的提及出现在
\pplsur{H. T.}{Kung}和\pplsur{Philip L.}{Lehman}的
里程碑式论文~\cite{Kung80}中。
学术界有一些关于此技术的额外使用~\cite{Manber82,Manber84,BarbaraLiskov1988ArgusCACM,Pugh90,Andrews91textbook,Pu95a,Cowan96a,Rastogi:1997:LPV:645923.671017,Gamsa99}，
但该领域的大量工作反而由
从业者完成~\cite{RichardRashid87a,Hennessy89,Jacobson93,AjuJohn95,Slingwine95,Slingwine97,Slingwine98,McKenney98}。

\QuickQuiz{
	垃圾回收器？
	被动序列化？
	系统参考点？
	Quiescent states？
	老化？
	代？
	为什么这些早期论文的研究人员就不能统一一下术语呢？？？
}\QuickQuizAnswer{
	有多个独立发明的机制大致类似于 RCU\@。
	每组发明者都不知道其他人的存在，
	因此每组理所当然地创造了自己的术语。
	而不同的术语使得任何一组都很难找到其他人。

	抱歉，但生活有时候就是这样！
}\QuickQuizEnd

到 2000 年，主导权已经转移到开源项目，
最显著的是 Linux 内核
社区~\cite{RustyRussell2000a,RustyRussell2000b,McKenney01b,McKenney01a,McKenney02a,Arcangeli03}。\footnote{
	在本书的 {\LaTeX} 源码中的 \co{bib/RCU.bib}
	可以找到一个包含超过 200 个条目的引用列表。}
RCU 于 2002 年末被 Linux 内核接受，
随后在可扩展性、健壮性、实时响应、
能源效率和专用使用场景方面进行了许多改进。
截至 2023 年，Linux 内核 RCU 仍在积极开发中。

\QuickQuiz{
	为什么 Kung 和 Lehman 的论文没有导致 RCU 被立即使用？
}\QuickQuizAnswer{
	一个原因是 Kung 和 Lehman 确实走在了时代前面。
	另一个原因是他们的方法虽然具有开创性，
	但没有考虑到许多软件工程和性能问题。

	要看出他们走在时代前面，考虑一下在他们的论文发表三年后，
	Paul 正在一个运行 BSD 2.8 的 PDP-11 系统上工作。
	这个系统缺乏任何自动配置，
	这意味着任何硬件修改，包括添加新磁盘驱动器，
	都需要手动编辑和重建内核。
	此外，这是一个单 CPU 系统，
	这意味着全系统同步只需简单地禁用中断。

	快进若干年，允许运行时更改硬件配置的多核系统
	变得司空见惯。
	这意味着 1980 年代内核源代码中隐式表示的
	硬件配置数据现在是一个可变的数据结构，
	每次 I/O 都会被访问。
	这样的数据结构很少更改，但可以随时更改。
	而这种读多属性适用于许多其他新时代的数据结构，
	包括关于网络（1980 年代很少见）、
	安全策略（1980 年代用物理锁）、
	软件配置（1980 年代运行时不可变）等等的数据结构。
	因此，1990 年代和 2000 年代比 1980 年代
	有更多机会让 RCU 展示其优势。

	Kung 和 Lehman 的软件工程方面的不足包括
	未标记读者（因此带来调试困难）、
	未提供干净的 RCU API（因此将其机制绑定到特定数据结构），
	以及不允许 grace period 后操作除释放内存外的任何其他操作
	（因此排除了许多 RCU 使用场景）。

	Kung 和 Lehman 提出了两种垃圾回收策略。
	第一种等待在给定时间运行的所有进程终止，
	这代表了另一个软件工程不足，
	排除了其机制在无限期运行的软件中的使用。
	第二种使用每对象引用计数，
	这大大复杂化了其读侧代码（因此代表又一个软件工程不足），
	并且在现代硬件上导致严重的缓存未命中开销
	（因此代表性能不足，参见例如
	\cref{fig:defer:Performance of RCU vs. Reference Counting,fig:defer:Performance of Preemptible RCU vs. Reference Counting}）。

	尽管软件工程和性能方面有这么长的不足清单，
	Kung 和 Lehman 的论文仍然是一项真正令人印象深刻的工作，
	特别是考虑到后来的许多工作
	（无论是独立的还是非独立的）
	都犯了同样的不足，甚至还有其他不足。
}\QuickQuizEnd

然而，在 2010 年代中期，RCU 的研究和开发在许多
社区和机构中出现了可喜的
回升~\cite{FransKaashoek2015ParallelOSHistory}。
\Cref{sec:defer:RCU Uses}描述 RCU 的使用，
\cref{sec:defer:RCU Implementations}描述 RCU 实现
（以及既创建又使用实现的工作），
最后，
\cref{sec:defer:RCU Validation}描述 RCU 及其使用的
验证和确认。

\subsubsection{RCU 的使用}
\label{sec:defer:RCU Uses}

\ppl{Phil}{Howard}和\ppl{Jon}{Walpole}（波特兰州立大学，PSU）
将 RCU 应用于红黑
树~\cite{PhilHowardPhD,PhilHoward2011RCUTMRBTree}，
结合使用软件事务内存同步的更新。
\ppl{Josh}{Triplett}和\ppl{Jon}{Walpole}（同样来自 PSU）
将 RCU 应用于可调整大小的
哈希表~\cite{JoshTriplettPhD,Triplett:2011:RPHash,JonCorbet2014RCUhash1,JonCorbet2014RCUhash2}。
其他 RCU 保护的可调整大小哈希表由
\ppl{Herbert}{Xu}~\cite{HerbertXu2010RCUResizeHash}和
\ppl{Mathieu}{Desnoyers}~\cite{PaulMcKenney2013LWNURCUhash}创建。

\ppl{Austin}{Clements}、\ppl{Frans}{Kaashoek}和\ppl{Nickolai}{Zeldovich}
（MIT）创建了一种 RCU 优化的平衡二叉树
（Bonsai）~\cite{AustinClements2012RCULinux:mmapsem}，
并将此树应用于 Linux 内核的 VM 子系统，
以减少 Linux 内核 \co{mmap_sem} 上的读侧竞争。
这项工作在一个以大量次要缺页错误为特征的微基准测试中
产生了数量级的加速，可扩展到至少 80 个 CPU。
这类似于\ppl{Peter}{Zijlstra}~\cite{PeterZijlstra2014SpeculativePageFault}
较早开发的补丁，两者都受到当时文件系统数据结构
对 RCU 读者不安全这一事实的限制。
\pplsur{Austin}{Clements}等人通过仅优化匿名页面的
缺页路径来避免此限制。
最近，文件系统数据结构已经对 RCU 读者变得
安全~\cite{JonathanCorbet2010dcacheRCU,JonathanCorbet2011dcacheRCUbug}，
因此也许这项工作可以为所有页面类型实现，
而不仅仅是匿名页面——\ppl{Peter}{Zijlstra}实际上最近已经
做了原型，\ppl{Laurent}{Dufour}和\ppl{Michel}{Lespinasse}
继续沿这些方向工作。
\ppl{Matthew}{Wilcox}、\ppl{Liam}{Howlett}和
\ppl{Suren}{Baghdasaryan}正在致力于使用 RCU 来实现对其他
内存管理数据结构的细粒度锁定和无锁访问。

\ppl{Yandong}{Mao}和\ppl{Robert}{Morris}（MIT）以及
\ppl{Eddie}{Kohler}（哈佛大学）创建了另一种名为
Masstree~\cite{Mao:2012:CCF:2168836.2168855}的 RCU 保护树，
结合了 B+ 树和 trie 的思想。
虽然此树比 RCU 保护的哈希表慢约 2.5 倍，
但它支持对键范围的操作，这与哈希表不同。
此外，Masstree 支持具有长共享键前缀的对象的高效存储，
并且还通过日志记录到大容量存储提供持久性。

该论文指出 Masstree 的性能可与 memcached 媲美，
即使 Masstree 持久存储更新而 memcached 没有。
该论文还将 Masstree 的性能与持久数据存储
MongoDB、VoltDB 和 Redis 进行了比较，
报告了 Masstree 的显著性能优势，
在某些情况下超过两个数量级。
另一篇论文~\cite{Tu:2013:STM:2517349.2522713}，
由\ppl{Stephen}{Tu}、\ppl{Wenting}{Zheng}、\ppl{Barbara}{Liskov}
和\ppl{Samuel}{Madden}（MIT）以及\pplsur{Eddie}{Kohler}撰写，
将 Masstree 应用于名为 Silo 的内存数据库，
在一个著名的事务处理基准测试上实现了每秒 700K 事务
（每分钟 42M 事务）。
有趣的是，Silo 在不持有锁的情况下保证线性化，
同时不承担 grace periods 的开销。

\ppl{Maya}{Arbel}和\ppl{Hagit}{Attiya}（以色列理工学院）
采用了更严格的
方法~\cite{MayaArbel2014RCUtree}来实现 RCU 保护的搜索树，
与 Masstree 一样允许并发更新。
该论文包括正确性证明，包括该树上所有操作
都是\IX{linearizable}的证明。
不幸的是，此实现通过在持有锁时承担 grace-period 等待的
全部延迟来实现线性化，这降低了仅更新工作负载的可扩展性。
解决此问题的一种方法是放弃
线性化~\cite{AndreasHaas2012FIFOisnt,PaulEMcKennneyAtomicTreeN4037}，
然而 Arbel 和 Attiya 创建了一种减少低端
grace-period 延迟的 RCU 变体。
当然，没有免费的午餐，这种 RCU 变体似乎在约 32 个 CPU 处
达到可扩展性极限。
虽然放弃线性化以获得性能和可扩展性有很多可取之处，
但看到学者们尝试替代 RCU 实现是非常好的。

\subsubsection{RCU 实现}
\label{sec:defer:RCU Implementations}

\ppl{Timothy}{Harris}创建了一种基于时间的用户空间
RCU~\cite{TimothyLHarris2001}，改进了之前由
Jacobson~\cite{Jacobson93}和 John~\cite{AjuJohn95}创建的版本。
之前这两种基于时间的方法各自假设读者持续时间有严格上限，
这在硬实时系统中可以正确工作。
在非实时系统中，当读者被中断、抢占或以其他方式延迟时，
这种方法容易失败。
然而，这种易失败的实现被独立发明两次的事实
表明了对类似 RCU 机制的深层需求。
\ppl{Timothy}{Harris}通过要求每个读者在开始其
读侧遍历之前获取全局时间基准的快照来改进这两项早期工作。
延迟释放读者可见对象，直到所有进程的读者快照
指示该对象移除之后的时间。
然而，全局时间基准在某些系统上可能既昂贵又不准确。

\ppl{Keir}{Fraser}创建了一种名为\IXacr{ebr}的用户空间 RCU，
用于非阻塞同步和软件事务
内存~\cite{KeirAnthonyFraserPhD,UCAM-CL-TR-579,KeirFraser2007withoutLocks}。
这项工作通过用软件计数器替换全局时钟来改进\ppl{Timothy}{Harris}的工作，
从而消除了当时商品系统全局时钟的大部分开销和所有不准确性。
有趣的是，这项工作一方面引用了 Linux 内核 RCU，
但也为原始的不可抢占 Linux 内核 RCU 实现
启发了\acr{qsbr}的名称。

\ppl{Mathieu}{Desnoyers}创建了用于跟踪的用户空间
RCU~\cite{MathieuDesnoyers2009URCU,MathieuDesnoyersPhD,MathieuDesnoyers2012URCU,PaulMcKenney2013LWNURCU,PaulMcKenney2013LWNURCUhash,PaulMcKenney2013LWNURCUhashAPI,PaulMcKenney2013LWNURCUqueuestack,PaulMcKenney2013LWNURCUqueuestackAPI,PaulMcKenney2013LWNURCUatomicop,PaulMcKenney2013LWNURCUmenagerie,PaulMcKenney2013LWNURCUAPI,PaulMcKenney2013LWNURCUlist,PaulMcKenney2013LWNURCUmenagerieRCU}，
已在许多项目中使用~\cite{MikeDay2013RCUqemu}。

布拉格查理大学的研究人员也在研究 RCU 实现，
包括\ppl{Andrej}{Podzimek}~\cite{AndrejPodzimek2010masters}和
\ppl{Adam}{Hraska}~\cite{AdamHraska2013RCUHelenOS}的论文。

\ppl{Yujie}{Liu}（利哈伊大学）、\ppl{Victor}{Luchangco}（Oracle 实验室）和
\ppl{Michael}{Spear}（同样利哈伊大学）~\cite{Liu:2013:MSA:2549695.2549732}
将可扩展非零指示器
（SNZI）~\cite{FaithEllen:2007:SNZI}用作 grace-period 机制。
其预期用途是实现软件事务内存
（见\cref{sec:future:Transactional Memory}），
这施加了线性化要求，反过来似乎限制了可扩展性。

类似 RCU 的机制也在进入 Java。
\pplsur{KC}{Sivaramakrishnan}等人~\cite{Sivaramakrishnan:2012:ERB:2258996.2259005}
使用类似 RCU 的机制消除了与 Java 垃圾回收器
交互时所需的读屏障，从而显著提高了性能。

\ppl{Ran}{Liu}、\ppl{Heng}{Zhang}和\ppl{Haibo}{Chen}（上海交通大学）
创建了 RCU 的专用变体，用于优化的
``被动读写锁''~\cite{RanLiu2014PassiveRWLock}，
类似于\ppl{Gautham}{Shenoy}~\cite{GauthamShenoy2006RCUrwlock}和
\ppl{Srivatsa}{Bhat}~\cite{SrivatsaSBhat2014RCUrwlock}创建的版本。
Liu 等人的论文从多个角度来看都很有趣~\cite{PaulEMcKenney2014ReadMostly}。

\ppl{Mike}{Ash}发布了~\cite{MikeAsh2015Apple}Apple Objective-C
运行时中类似 RCU 原语的描述。
此方法通过指定的代码范围标识读侧临界区，
因此有资格作为另一种实现零读侧开销的方法，
虽然对跨多个函数的大型读侧临界区提出了一些有趣的实际挑战。

\ppl{Pedro}{Ramalhete}和\ppl{Andreia}{Correia}~\cite{PedroRmalhete2015PoorMansRCU}
产生了``穷人的 RCU''，尽管使用了一对
\IXhpl{reader-writer}{lock}，仍然为
读者提供了\IXalt{lock-free}{lock free}的\IXpl{forward-progress guarantee}~\cite{PaulEMcKenney2015ReadMostly}。

\ppl{Maya}{Arbel}和\ppl{Adam}{Morrison}~\cite{Arbel:2015:PRR:2858788.2688518}
产生了``谓词 RCU''，它努力缩短 grace-period 持续时间，
以高效支持在 grace periods 期间持有更新侧锁的算法。
这导致减少了更新批处理到 grace periods 中的数量和可扩展性，
但确实成功提供了短 grace periods。

\QuickQuiz{
	为什么不在等待 grace period 之前释放锁，
	或者使用像 \co{call_rcu()} 这样的东西代替等待 grace period？
}\QuickQuizAnswer{
	作者希望支持\IX{linearizable}的树操作，
	使得对树的并发添加、删除和搜索看起来
	以某种全局一致的顺序执行。
	在他们的搜索树中，这需要在 grace periods 期间持有锁。
	（在大多数情况下最好放弃线性化作为要求，
	但线性化是一个出人意料地流行的（且昂贵的！\@）要求。）
}\QuickQuizEnd

\ppl{Alexander}{Matveev}（MIT）、\ppl{Nir}{Shavit}（MIT 和特拉维夫大学）、
\ppl{Pascal}{Felber}（纳沙泰尔大学）和\ppl{Patrick}{Marlier}（同样纳沙泰尔大学）~\cite{Matveev:2015:RLS:2815400.2815406}
产生了一种类似 RCU 的机制，可以被认为是
显式标记只读事务的软件事务内存。
他们的使用场景要求在 grace periods 期间持有锁，
这限制了可扩展性~\cite{PaulEMcKenney2015ReadMostly,PaulEMcKenney2015ReadMostlySidebar}。
这似乎是第一个很好地利用 \co{rcutorture} 测试套件的
学术 RCU 相关工作，也是第一个向 Linux 内核 RCU
提交性能改进并被接受到 v4.4 的工作。

\ppl{Alexander}{Matveev}的 RLU 之后是
\ppl{Jaeho}{Kim}等人的 MV-RLU~\cite{Kim:2019:MSR:3297858.3304040}。
这项工作通过允许多个并发更新、避免在 grace periods 期间
持有锁以及使用异步 grace periods
（例如 \co{call_rcu()} 代替 \co{synchronize_rcu()}）
来改善 RLU 的可扩展性。
该论文还做了一些有趣的性能评估选择，
在\cpageref{sec:future:Deferred Reclamation}上的
\cref{sec:future:Deferred Reclamation}中有进一步讨论。

\ppl{Adam}{Belay}等人创建了一个 RCU 实现，
保护其 IX 操作系统中 TCP/IP 地址解析协议（ARP）
使用的数据结构~\cite{Belay:2016:IOS:3014162.2997641}。

\ppl{Geoff}{Romer}和\ppl{Andrew}{Hunter}（均在 Google）
为将基于 cell 的 API 用于单例数据结构的 RCU 保护
提出了纳入 C++ 标准的
建议~\cite{GeoffRomer2018C++DeferredReclamationP0561R4}。

\ppl{Dimitrios}{Siakavaras}等人将
HTM 和 RCU 应用于搜索树~\cite{Siakavaras2017CombiningHA,DimitriosSiakavaras2020RCU-HTM-B+Trees}，
\ppl{Christina}{Giannoula}等人使用 HTM 和 RCU 进行
图着色~\cite{ChristinaGiannoula2018HTM-RCU-graphcoloring}，
\ppl{SeongJae}{Park}等人使用 HTM 和 RCU 来优化
\IXacr{numa}系统上的高竞争锁。

\ppl{Alex}{Kogan}等人将 RCU 应用于构建可扩展
地址空间的范围锁~\cite{AlexKogan2020RCUrangelocks}。

2023 年 6 月 17 日，ISO C++ 标准委员会投票将 RCU
纳入 C++26~\cite{PaulEMcKennney2023C++26RCU4}。

RCU 的生产使用列在\cref{sec:defer:Production Uses of RCU}中。

\subsubsection{RCU 验证}
\label{sec:defer:RCU Validation}

在 2017 年初，普遍认识到几乎任何 bug 都是潜在的安全漏洞，
因此验证和确认是一等关注事项。

石溪大学的研究人员产生了一个 RCU 感知的数据竞争
检测器~\cite{AbhinavDuggal2010Masters,JustinSeyster2012PhD,Seyster:2011:RFA:2075416.2075425}。
\ppl{Alexey}{Gotsman}（IMDEA）、\ppl{Noam}{Rinetzky}（特拉维夫大学）
和\ppl{Hongseok}{Yang}（牛津大学）发表了
一篇论文~\cite{AlexeyGotsman2012VerifyGraceExtended}，
用分离逻辑表达 RCU 的形式语义，并继续研究并发的其他方面。

\ppl{Joseph}{Tassarotti}（卡内基梅隆大学）、\ppl{Derek}{Dreyer}（马普
软件系统研究所）和\ppl{Viktor}{Vafeiadis}
（同样 MPI-SWS）~\cite{JosephTassarotti2015RCUproof}
产生了用户空间 RCU 的\IXacrf{qsbr}变体的
手动形式正确性证明~\cite{MathieuDesnoyers2009URCU,MathieuDesnoyers2012URCU}。
\ppl{Lihao}{Liang}（牛津大学）、\pplmdl{Paul E.}{McKenney}（IBM）、
\ppl{Daniel}{Kroening}和\ppl{Tom}{Melham}
（均为牛津大学）~\cite{LihaoLiang2016VerifyTreeRCU}
使用\IXacrf{cbmc}~\cite{EdmundClarke2004CBMC}
产生了 Linux 内核 Tree RCU 重要部分的机械正确性证明。
\ppl{Lance}{Roy}~\cite{LanceRoy2017CBMC-SRCU}使用 CBMC 产生了
Linux 内核 sleepable RCU（SRCU）~\cite{PaulEMcKenney2006c}
重要部分的类似正确性证明。
最后，\ppl{Michalis}{Kokologiannakis}和\ppl{Konstantinos}{Sagonas}
（雅典国家技术大学）~\cite{MichalisKokologiannakis2017NidhuggRCU,MichalisKokologiannakis2019RCUstatelessModelCheck}
使用 Nidhugg 工具~\cite{CarlLeonardsson2014Nidhugg}
产生了 Linux 内核 Tree RCU 更大部分的机械正确性证明。

这些工作都没有发现除了专门注入以测试验证工具的 bug 之外的任何 bug。
相比之下，\ppl{Alex}{Groce}（俄勒冈州立大学）、\ppl{Iftekhar}{Ahmed}、
\ppl{Carlos}{Jensen}（均为 OSU）和\pplmdl{Paul E.}{McKenney}
（IBM）~\cite{Groce:2015:VMC:2916135.2916190}
自动变异 Linux 内核 RCU 的源代码以测试
\co{rcutorture} 测试套件的覆盖率。
该工作发现了此套件覆盖率中的几个漏洞，
其中一个隐藏了 Tiny RCU 中的一个真实 bug（已修复）。

如果幸运的话，所有这些验证工作最终将产生
更多更好的工具来验证并发代码。

